



\documentclass{article}

\usepackage{iclr2027/iclr2027_conference,times}

\usepackage[utf8]{inputenc} % allow utf-8 input
\usepackage[T1]{fontenc}    % use 8-bit T1 fonts
\usepackage{hyperref}       % hyperlinks
\usepackage{url}            % simple URL typesetting
\usepackage{booktabs}       % professional-quality tables
\usepackage{amsfonts}       % blackboard math symbols
\usepackage{amsmath}
\usepackage{nicefrac}       % compact symbols for 1/2, etc.
\usepackage{microtype}      % microtypography
\usepackage{xcolor}         % colors
\usepackage{graphicx}
\usepackage{amsthm}
\usepackage{placeins}
% Upright (non-italic) body for all theorem-like environments, per style rule.
\theoremstyle{definition}
\newtheorem{theorem}{Theorem}
\newtheorem{proposition}{Proposition}
\newtheorem{lemma}{Lemma}
\newtheorem{corollary}{Corollary}
\newtheorem{definition}{Definition}
\theoremstyle{remark}
\newtheorem{remark}{Remark}
\newcommand{\nnote}[1]{\textcolor{green}{\textbf{[Note:]}}}
\newcommand{\mnote}[1]{\textcolor{blue}{\textbf{[Note:]}}}
\newcommand{\iclrTODO}[1]{%
  \par\noindent\fcolorbox{red}{yellow!15}{%
    \parbox{\dimexpr\linewidth-2\fboxsep-2\fboxrule\relax}{%
      \centering\color{red}\Large\bfseries ICLR TODO: #1}}\par}
\newcommand{\deprecated}[1]{\textcolor{red}{\textbf{[DEPRECATED FOR ICLR: #1]}}}

\title{OptEvolve: Hardware-Aware Search for Certified Optimization Solvers}


% The \author macro works with any number of authors. There are two commands
% used to separate the names and addresses of multiple authors: \And and \AND.
%
% Using \And between authors leaves it to LaTeX to determine where to break the
% lines. Using \AND forces a line break at that point. So, if LaTeX puts 3 of 4
% authors names on the first line, and the last on the second line, try using
% \AND instead of \And before the third author name.


\author{%
  Miria Feng\thanks{Use footnote for providing further information
    about author (webpage, alternative address), not for acknowledging
    funding agencies.} \\
  Department of Electrical Engineering\\
  Stanford University\\
  Stanford, CA 94305 \\
  \texttt{miria00@stanford.edu} \\
  \And
  Nikhil Devanathan \\
  Department of Electrical Engineering \\
  Stanford University \\
  Stanford, CA 94305 \\
  \texttt{ndev@stanford.edu} \\
  % examples of more authors
  % \And
  % Coauthor \\
  % Affiliation \\
  % Address \\
  % \texttt{email} \\
  % \AND
  % Coauthor \\
  % Affiliation \\
  % Address \\
  % \texttt{email} \\
  % \And
  % Coauthor \\
  % Affiliation \\
  % Address \\
  % \texttt{email} \\
  % \And
  % Coauthor \\
  % Affiliation \\
  % Address \\
  % \texttt{email} \\
}

% Keep this commented for an anonymous submission. Uncomment only for the
% camera-ready version after acceptance.
%\iclrfinalcopy

\begin{document}


\maketitle


\begin{abstract}
Solving a convex problem repeatedly on fixed hardware requires choosing both a
mathematical iteration and the way it is executed. We show that time to a
certified solution factorises as $T = N(\text{method}, \text{cadence},
\text{tolerance}) \cdot t(\text{execution}, \text{fusion}, \text{size},
\text{device})$, and that the certificate cadence is the only recipe gene
appearing in both factors. The factorisation predicts that the two halves can
be tuned separately, which is what we measure, and it names in advance the one
axis on which that separation fails.

We then attack each factor. For the per-iteration cost we localise a factor of
two to a single implementation choice, an inner loop whose trip count is a
traced rather than a static value, and we build a temporally fused CUDA kernel
that advances several primal-dual iterations per launch from shared memory. It
is admitted against a sequential reference to $4.6 \times 10^{-15}$ and is
$1.86\times$ faster end to end on total-variation denoising.

For the iteration count, the step-size ratio is predictable from one measured
property of the problem, the ratio of the cost norm to the constraint norm. A
rule fitted on min-cost flow alone, applied unchanged to optimal transport and
to multi-commodity flow, lands within $1.13\times$ of a per-instance oracle
across 162 held-out instances, against $1.83\times$ for the best single fixed
choice.

We measure two GPU architectures, and we report the settings at which each
invariance fails.
\end{abstract}

% ---------------------------------------------------------------------------
% CENTRAL THESIS. Stated once, immediately after the abstract, so a reviewer
% cannot miss it and so every later section can refer back to one sentence.
% The abstract above is UNCHANGED from the 6 September manuscript and does not
% yet reflect these results; reconciling it is a coauthor decision.
% ---------------------------------------------------------------------------
\begin{center}
\fbox{\begin{minipage}{0.93\linewidth}
\textbf{Central thesis.} For a certified first-order convex solver, time to a
checked solution factorises as
\[
T \;=\; \underbrace{N(\text{method},\, \text{cadence},\, \text{tolerance})}_{\text{iterations}}
\;\times\;
\underbrace{t(\text{execution structure},\, \text{fusion},\, \text{size},\, \text{device})}_{\text{cost per iteration}} ,
\]
and the certificate cadence is the only recipe gene that appears in both
factors. The two halves can therefore be searched separately, which is why
joint search does not beat staged search, and the one axis where that
separation fails is the one the factorisation names in advance. We measure
each factor, identify the mechanism behind the execution term, build a fused
kernel that attacks it, and show that the method term is predictable from a
single measured property of the problem.
\end{minipage}}
\end{center}

% Section 1: Introduction and contributions
% OptEvolve, ICLR 2027 submission draft, 10 September 2026.
% Every number carries a "% source:" comment pointing at a file under
% /home/miria/OptEvolve/v2/results/. Do not add a number without one.
% STYLE CONTRACT: no em-dashes, no italics, plain declarative prose.
% PREAMBLE REQUIREMENT: \usepackage{booktabs} and \usepackage{amsmath}.
% This file is a section body only. No preamble, no \begin{document}.

\section{Introduction}
\label{sec:intro}

Two axes are routinely conflated in work on fast first-order solvers, and we
separate them before anything else. Arithmetic width is the floating-point
format in which iterates are stored and updated, float32 or float64.
Termination tolerance is the threshold a certificate must meet before a solve
is declared finished, here a relative duality gap of $10^{-4}$ or $10^{-6}$. On
total-variation denoising of retinal OCT scans at $n=256$, float32 iterates
with a float64 certificate reach mean iteration counts of $584.0$, $597.3$ and
$682.7$ at certificate cadences $K=16$, $64$ and $256$, certifying 6 of 6. The
per-instance counts are identical to float64 at $K=64$ and $K=256$ under both
schemes and at $K=16$ under Condat-Vu; under PDHG at $K=16$ float32 gives
$584.0$ against $578.7$ in float64. At $10^{-6}$ the same configuration stops
certifying: Condat-Vu solves 0 of 6 and PDHG 1 of 6, against 5 of 6 for float64
at every cadence.
% source: results/tv_ffi_20260906/f32row_4090/factorial.json (the float32 rows)
% and results/tv_ffi_20260906/cadence_4090_v1/factorial.json (the float64 rows
% at K=16, 64 and 256; f32row_4090 carries a float64 row only at K=64).
Width left the trajectory unchanged in most measured cells and moved it by at
most $1.9\%$ in the rest: the per-instance counts are identical at $n=256$ for
$K=64$ and $K=256$, while PDHG at $K=16$ moves $578.7$ to $584.0$, and on the
second architecture at $K=64$ the mean moves from $576.0$ to $586.7$ at $n=128$
and from $1066.7$ to $1077.3$ at $n=496$. What width changed is the reachable
tolerance.
% source: results/bw_20260909/sz128_bw/factorial.json,
% results/bw_20260909/sz496_bw/factorial.json,
% results/tv_ffi_20260906/f32row_4090/factorial.json
Work that labels loose tolerances as low precision lets one phrase carry both
meanings, so we never let the two share a word.

That separation is the first case of a general one. We measure that wall-clock
time to a certified solution factorises as
\begin{equation}
T \;=\; N(\text{method},\, \text{cadence},\, \text{tolerance}) \;\times\;
        t(\text{hardware},\, \text{execution},\, \text{fusion},\, \text{size}),
\label{eq:factorisation}
\end{equation}
with $N$ the iteration count and $t$ the cost per iteration. Neither factor is
assumed. The float64 per-instance counts behind the triple above are exactly
equal on a second GPU architecture at $n=256$ at all three cadences, giving the
same $584.0$, $597.3$ and $682.7$ under Condat-Vu; the float32 arm was run on
that architecture only at $K=64$, where it also gives $597.3$.
% source: results/bw_20260909/sz256_bw/factorial.json against
% results/tv_ffi_20260906/cadence_4090_v1/factorial.json
On grid-topology min-cost flow all 24 static-versus-traced loop pairs and all
12 cross-scheme instances have exactly equal $N$, of which 22 pairs and 11
instances converged; the remaining instance stops at the 20000-iteration cap
under both loop structures and both schemes.
% source: results/netflow_20260909/grid_4090.json (sizes 12, 16, 24, 32)
Conversely $t$ is insensitive to scheme identity. Two separate three-way
decompositions carry that, and they have different responses and different
factor sets. On log penalised SGM-1 time over algorithm, precision and cadence,
the algorithm term is $0.00\%$ of the variance at $10^{-4}$ and $0.34\%$ at
$10^{-6}$; that null is measured over two levels whose recorded step sizes and
relaxation are identical, $\tau=0.05$, $\sigma=1.0$, $\rho=1.9$, and whose
update maps differ by $7.1\times 10^{-4}$ in the recorded map difference
against $1.5\times 10^{-2}$ at the balanced setting, so it is a null over two
nearly identical operators. On log per-iteration cost over execution, precision
and cadence, cadence is the largest single term at $47.43\%$, over a 12-cell
design one of whose cells, host phases at float32, is a non-catalog diagnostic
rather than a shipped-genome run.
% source: results/artifacts_20260910/fanova.json, decompositions
% alg_x_prec_x_cadence__logSGM1__tol0.0001 (algorithm 0.0012 percent),
% alg_x_prec_x_cadence__logSGM1__tol1e-06 (algorithm 0.3444 percent) and
% exec_x_prec_x_cadence__logT__tol0.0001 (cadence 47.4302 percent).
% source: results/method_axis_20260909/map_difference.json and factorial.json
% source for the non-catalog cell: fanova.json records host_phases|f32 =
% tv_ffi_20260906/px999999_4090; results/tv_ffi_20260906/DECOMPOSITION_REPORT.md
% lines 128-130 disclaim it as a /tmp policy-registry patch, not the genome.
Cadence is the only gene appearing in both factors, so it is the sole coupling
point of the two.

\subsection{Contributions}
\label{sec:intro:contributions}

\paragraph{(a) The factorisation, measured, with cadence as the sole coupling
point.}
Inside the deployed solver the eight end-to-end arms at each size record
identical mean iteration counts of $554.7$, $640.0$, $810.7$ and $1024.0$ from
$n=128$ to $496$ while their run times differ by up to $2.629\times$, at
$n=128$, between the traced-bound baseline and the $8\times 4$ tile.
% source: results/kernel_fusion_20260909/e2e/full_4090_sz{128,256,384,496}.json
% CAVEAT for the experiments section: scheme identity is inert at K=256
% everywhere and at K=64 at 1e-4. It splits only at K=16, only at n=256, and
% only in float64 (condat_vu 584.0 against pdhg 578.7, on both devices); in
% f32_cert64 at K=16 both schemes give 584.0, and at n=128 and n=496 on the
% Blackwell f64_K16 gives 552.0 and 1040.0 for both schemes. Never write
% "identical across both schemes" without the cadence and size qualifiers.
% source: results/tv_ffi_20260906/cadence_4090_v1/factorial.json,
% results/tv_ffi_20260906/f32row_4090/factorial.json,
% results/bw_20260909/sz128_bw/factorial.json, sz496_bw/factorial.json
We also settle the mechanism of the largest effect on $t$. A traced rather than
static inner-loop trip count costs $2.03\times$ at $n=256$, the geometric mean
over the six measured cadences, spanning $1.94\times$ to $2.13\times$, and
$2.34\times$ at $n=128$, spanning $2.30\times$ to $2.37\times$. The penalty is
per iteration and not proportional to the number of outer trips: at $n=256$ it
falls only from $12943$\,ns at 32 trips to $11531$\,ns at one, so a 32-fold
reduction in outer trips removes $11\%$ of it, and at $n=128$ from $11972$\,ns
to $10769$\,ns, removing $10\%$. The certificate term, 24 differences computed
from the 48 measured cells, spans $-240$ to $+845$\,ns and is concentrated at
the shortest cadences: $845$ and $697$\,ns at $K=16$ and $n=128$, against
$-240$\,ns at $K=256$ and $n=256$. That is the $1/K$ signature of a real
per-check cost, but the largest value is $7\%$ of the traced penalty in the
same cell, so the check cannot account for the penalty, and the penalty does
not scale with outer trips. Neither a per-certificate-check nor a
pure-unrolling explanation survives.
% source: results/finish_20260910/loop_mechanism.json (48 cells over
% n in {128,256}, K in {16,32,64,128,256,512}, traced/static, cert on/off)

\paragraph{(b) A temporally fused kernel, numerically admitted, whose optimum
moves with size and device.}
We run $T$ PDHG iterations per launch in shared memory with a halo of width
$T$.
% source: results/kernel_fusion_20260909/tv_temporal.cu, 149 lines
Against $T$ sequential reference launches the result deviates by at most
$4.649\times 10^{-15}$ over eight fused iterations at $\rho=1.9$; the worst
case is the two $T=8$ tiles at a $100\times 130$ non-square size with the
proximal term off. The archived admission run covers the RTX 4090 only.
% source: results/artifacts_20260910/kernel_admission_4090.log (106 lines, 102
% measured rows over 17 tile variants x 3 sizes x 2 proximal settings, all at
% rho=1.9; worst deviation 4.649e-15; PASS). Header records device sm_89 only.
Table~\ref{tab:intro-e2e} gives the end-to-end result. The tile is a search
dimension, not a default: the worst 4090 tile gives $0.926\times$ at $n=384$
and $0.743\times$ at $n=496$, slower than not fusing,
% source: results/kernel_fusion_20260909/e2e/full_4090_sz384.json (0.926x) and
% full_4090_sz496.json (0.743x); both are the 8x4 tile against the
% traced-bound baseline
and on the Blackwell at $n=496$ even the best tile reaches $0.925\times$ of the
static-loop path, so the correct decision there is not to fuse.
% source: results/kernel_fusion_20260909/e2e_bw/full_bw_sz496.json
% (best tile tf16x2 at 0.17379 s against the static loop at 0.16067 s)
The gain over the static loop is U-shaped in size. Over six raw OCT crops it
falls from $1.571\times$ at 128\,px to $1.240\times$ at 496\,px, and over four
montage problems it rises from $1.060\times$ at 992\,px through $1.346\times$
at 1488\,px to $1.410\times$ at 1984\,px. The two limbs use different instance
populations, four montages against six crops, and the trough sits at that
boundary.
% source: results/kernel_fusion_20260909/e2e/full_4090_sz{128,256,384,496}.json
% for 1.571x (2.6294/1.6733), 1.386x, 1.266x and 1.240x, and
% large/large_m2.log, large_m3.log, large_m4.log for 1.060x, 1.346x and 1.410x
% Sizes above 496 px are four montage problems per size, tiled 2x2, 3x3 and 4x4
% from 16, 36 and 64 distinct scans; sizes 128 to 496 are six raw crops. State
% both counts in the text of the experiments section.
% NOTE: the 1.410x here is fusion at 1984 px. It is NOT the fp32/fp64 envelope
% number the literature review told us not to lean on. Keep them apart in text.

\begin{table}[t]
\centering
\caption{Time to a certified $10^{-4}$ gap over six OCT scans, $K=256$, median
of 3 repeats per instance summed over the six instances, against the shipped
traced-bound baseline. Every arm certified 6 of 6 at every size on both
devices. Best tile is $B \times T$. The Blackwell entry at $n=496$ is the best
tile, and the static loop beats it there, $1.243\times$ against
$1.149\times$.}
\label{tab:intro-e2e}
\begin{tabular}{lrrrr}
\toprule
$n$ & 128 & 256 & 384 & 496 \\
\midrule
RTX 4090 baseline (s)  & 0.06647 & 0.07733 & 0.13486 & 0.22564 \\
RTX 4090 speedup       & 2.629$\times$ & 1.867$\times$ & 1.634$\times$ & 1.356$\times$ \\
RTX 4090 best tile     & 8$\times$4 & 24$\times$4 & 24$\times$4 & 32$\times$4 \\
\midrule
Blackwell baseline (s) & 0.05732 & 0.07299 & 0.13750 & 0.19964 \\
Blackwell speedup      & 1.732$\times$ & 1.450$\times$ & 1.625$\times$ & 1.149$\times$ \\
Blackwell best tile    & 8$\times$4 & 24$\times$4 & 32$\times$4 & 16$\times$2 \\
\bottomrule
\end{tabular}
\end{table}
% source: results/kernel_fusion_20260909/e2e/full_4090_sz{128,256,384,496}.json
% source: results/kernel_fusion_20260909/e2e_bw/full_bw_sz{128,256,384,496}.json
% The n=384 Blackwell baseline 0.13750 is from the JSON; REPORT.md section 5
% omits that row. Use the JSON.
% PROTOCOL: all fusion numbers in this paper are median of 3. A median-of-9
% re-measurement one day later gives 1.940x at n=256 against 1.867x and 1.428x
% at n=496 against 1.356x, and it also moves the baseline (0.07733 to 0.07885
% at n=256, 0.22564 to 0.22805 at n=496) and the static-loop gain at n=496
% (1.093x to 1.179x). The derived ratios move by 4 to 8 percent, so the
% reported figures are conservative but are not resolved to four figures.
% Never mix the two protocols.
% source: results/finish_20260910/fusion_r9_sz256.log, fusion_r9_sz496.log

\paragraph{(c) A symbolic construction gate with a family-dependent admissible
region.}
Admissibility is an exact function of $\theta = \tau\sigma\|L\|^2$ and not of
the instance: over 2106 settings evaluated the gate admits 351 of 351 at each
$\theta \in \{0.1, 0.3, 0.6, 0.9, 0.99\}$ and 0 of 351 at $\theta = 1.02$. The
351 rejected settings never ran, so 2106 counts settings evaluated, not solves.
% source: results/structure_20260909/structure_summary_4090.json,
% gate_is_function_of_theta_only and totals.runs
The region is family dependent. One setting is admitted on TV and certifies
6 of 6, and is rejected before any iteration runs on every grid-topology
min-cost flow instance, at gate values $1.02267$, $1.03040$, $1.03577$ and
$1.03781$ for sizes 12, 16, 24 and 32, while the same setting is admitted on
all nine regular-topology cells at gate values $0.51994$ to $0.96969$. On these
instances the linear objective makes the smooth term vanish, so the Condat-Vu
and PDHG conditions record the same number.
% source: results/structure_20260909/structure_summary_4090.json (cells block,
% raw_t025_s025_admit and raw_t025_s025_gate_lhs) and
% results/netflow_20260909/grid_4090.json (the two error strings)
% check: 0.25*0.25*4.045095034381878^2 = 1.0226746148238077 at size 12
For optimal transport the recorded operator-norm bound is the implementation's
$1.02$ safety margin on the exact norm $\sqrt{m+n}$, recovered to within
$4\times 10^{-14}$ relative at all three measured sizes, $8.160000$,
$11.539983$ and $14.133535$ against $8$, $11.313708$ and $13.856406$, so the
admissible product contracts as $1/(m{+}n)$.
% source: results/finish_20260910/omega_ot_4090.json (Lnorm per cell)
Step sizes and relaxation together are worth $4.31\times$ on mean iteration
count at $K=256$ and $4.84\times$ at $K=64$, and they decide feasibility at
$10^{-6}$, where the balanced setting solves 0 of 6 and the banked setting
5 of 6 at both cadences.
% source: results/method_axis_20260909/factorial.json (2944/682.6667 = 4.3125,
% 2890.6667/597.3333 = 4.8393)
% CAVEAT: banked is (tau=0.05, sigma=1.0, rho=1.9), balanced is
% (tau=0.25, sigma=0.25, rho=1.5). rho differs too, so 4.31x is the value of
% the whole parameter setting, not of the step sizes alone.

\paragraph{(d) A step-size rule read off a pre-solve quantity, fitted on one
family and transferred to three held-out sets.}
The scale ratio $\omega_0 = \|c\|/\|b\|$ varies little within a family, which
is why the recipe looked unpredictable: over 24 min-cost-flow cell-and-seed
pairs at unit rescaling it spans only $6.969$ to $9.125$, and over the nine
optimal-transport cell-and-seed pairs only $0.446$ to $0.793$.
% source: results/finish_20260910/omega_4090.json, omega_bw.json,
% omega_ot_4090.json (rows with alpha = 1)
Rescaling $c \mapsto \alpha c$ is exactly
$(\tau,\sigma) \mapsto (\alpha\tau, \sigma/\alpha)$, which fixes the gate
product and moves $r = \tau/\sigma$ by $\alpha^2$, so it sweeps $\omega_0$ over
decades without touching admissibility. Across the 1944 rows of the two
original sweeps $\tau\sigma\|L\|^2$ reproduces the requested $0.99$ to within
$3.3\times 10^{-16}$. Since $\theta$ is a fixed input to those sweeps, the
evidence that admissibility depends on $\theta$ alone is the separate $\theta$
sweep of contribution (c), not this reproduction. Sweeping $\alpha$ over four
decades on 8 cells gives 1510 converged runs and 209 of 216 instances with a
converged ratio.
% source: results/finish_20260910/omega_4090.json and omega_bw.json
% (args.theta = 0.99 in both; 1944 rows, 1510 converged; exact float equality
% to 0.99 holds in only 774 of the 1944 rows, the rest agree to 3.3e-16)
Table~\ref{tab:intro-omega} is the result. A rule of the form
$r = C/\omega_0^{\,p}$, with $p = 2$ and $C = 0.977$ fitted on the min-cost
flow sweep alone, cuts geometric mean regret against a per-instance oracle from
$2.168\times$ to $1.131\times$ pooled over 162 held-out instances it never
saw, a gain of $1.92\times$. The pair $(p, C) = (2.00, 0.977)$ is the joint
argmin of geometric mean regret on the fitted set, over $p$ on a $0.05$ grid
and $C$ on a $0.01$ dex grid, and the same pair is used in every later section.
Two of the three held-out sets are different problem families and the third is
min-cost flow at topologies not used for fitting. The rule certifies all 54
instances of every held-out set, where the best fixed ratio for that set
certifies 54 of 54 for optimal transport, 48 for multi-commodity flow and 44
for the new topologies.
% source: results/finish_20260910/omega_wide_bw.json (fitted), omega_ot_wide.json,
% omega_mcf.json, omega_holdout.json (held out), all recomputed

\begin{table}[t]
\centering
\caption{Geometric mean iteration-count regret against a per-instance oracle,
under the single protocol used in every table of this paper that reports these
numbers. An instance is one cell, seed and rescaling $\alpha$; the oracle is
the minimum certified $N$ over that instance's own ratio grid; a run that fails
to certify inside the 60000-iteration cap is charged 60000 iterations, so every
policy is scored on every instance of the set. The best fixed ratio is the
single grid ratio minimising geometric mean regret on the evaluation set
itself, chosen separately for each set, which favours the baseline. The rule is
$r = 0.977/\omega_0^{\,2}$ snapped to the nearest grid ratio in log, fitted on
the first row only. The pooled row is the geometric mean over the 162 held-out
instances, each family carrying its own best fixed ratio. Certified counts are
the instances each policy solved inside the cap, out of the instance count.}
\label{tab:intro-omega}
\begin{tabular}{lrrrrrr}
\toprule
evaluation set & instances & $\omega_0$ range & best fixed $r$ & rule & rule worst & gain \\
\midrule
min-cost flow, fitted on      & 108 & 0.070 to 825   & 2.779$\times$ & 1.186$\times$ & 2.125$\times$ & 2.34$\times$ \\
optimal transport             &  54 & 0.0054 to 79.3 & 1.428$\times$ & 1.179$\times$ & 1.562$\times$ & 1.21$\times$ \\
multi-commodity flow          &  54 & 0.023 to 312   & 2.840$\times$ & 1.032$\times$ & 1.262$\times$ & 2.75$\times$ \\
min-cost flow, new topologies &  54 & 0.074 to 806   & 2.514$\times$ & 1.190$\times$ & 1.825$\times$ & 2.11$\times$ \\
\midrule
pooled held out               & 162 & 0.0054 to 806  & 2.168$\times$ & 1.131$\times$ & 1.825$\times$ & 1.92$\times$ \\
\bottomrule
\end{tabular}
\end{table}
% source: results/finish_20260910/omega_wide_bw.json (min-cost flow, fitted on:
% cells grid:16, grid:24, regular:256:4, regular:512:4),
% omega_ot_wide.json (ot:32:32, ot:64:64), omega_mcf.json (mcf:16:2, mcf:16:4),
% omega_holdout.json (regular:256:2, grid:32, neither used for fitting).
% Best fixed ratio under the protocol in the caption: r = 0.03 on the fitted
% set, r = 100 on OT, r = 0.1 on MCF, r = 0.003 on the new topologies.
% Certified instances, fixed ratio then rule: 100/108 and 108/108 on the fitted
% set; 54/54 and 54/54 on OT; 48/54 and 54/54 on MCF; 44/54 and 54/54 on the
% new topologies; 146/162 and 162/162 pooled. A single fixed ratio shared
% across all three held-out families is worse still, r = 3 at 2.820x, so the
% pooled row understates the gain.
% Rule to four decimals: 1.1858 / 1.1786 / 1.0318 / 1.1899 / 1.1311; fixed
% 2.7790 / 1.4279 / 2.8395 / 2.5136 / 2.1681.

The exponent is not identified more finely than a bracket. Two independent
estimators approach each other from opposite directions as censoring by the
ratio grid falls, and they have not met. The log-log slope of the argmin ratio
on $\omega_0$ moves from $-0.94$ on the original four-decade grid, where 102 of
209 optima sit on an edge of the grid, $43.5\%$ on the floor and $5.3\%$ on the
ceiling, to $-1.54$ on a six-decade grid with 33 of 108 on an edge, $-1.58$ on
a seven-decade eight-seed grid with 46 of 144, $-1.68$ on a finer 19-point grid
with 15 of 80, and $-1.74$ on a ten-decade grid with 14 of 88. The
regret-minimising exponent barely moves: $2.10$, $2.00$, $1.90$, $2.00$ and
$2.00$ on those same five grids. On the ten-decade grid, the least censored,
geometric mean regret against a per-instance oracle is $1.070\times$ at $p=1.8$
and $1.062\times$ at $p=2.0$. The two estimators bracket the exponent between
about $1.7$ and $2.1$. That is consistent with the value 2 implied by PDLP's
primal-weight heuristic, and at its lower end below it; our grid resolution
does not separate the two, and we do not claim the exponent is exactly 2.
% source: results/finish_20260910/omega_4090.json and omega_bw.json (original),
% omega_wide_bw.json, omega_seeds8.json, omega_fine_4090.json (19 points, 1e-5
% to 300, spaced by a factor of about 2.6), omega_ultrawide_bw.json (ten
% decades, ratios 1e-6 to 1e4, 88 instances).
% CONVENTIONS, the same two used in section 5 and section 6 and named in the
% caption of the exponent table there: an instance counts as censored when any
% ratio attaining its minimum certified N sits at an endpoint of that instance's
% grid, and tied argmin ratios are averaged geometrically before the log-log
% fit. Slopes to four decimals under that convention: -0.9365, -1.5412, -1.5781,
% -1.6834, -1.7380. The other two tie rules move each slope by at most 0.011.
% Regret profile on the ten-decade grid, C on a 0.01 dex grid and p on a 0.05
% grid: 1.360 at p=1.0, 1.098 at 1.6, 1.070 at 1.8, 1.062 at 2.0, 1.067 at 2.1,
% 1.096 at 2.5, 1.174 at 3.0. The 19-point grid has run and is included above;
% it does not move either estimator off the bracket.

Cheap a-priori predictors fail on the same task. A roofline model picks the
best fusion tile 1 of 7 times, whether or not a fitted per-launch cost of
$10^{-5}$\,s is added, while mean Spearman rises from $0.355$ for the plain
model to $0.567$ with that cost. The pick rate does not order with the
correlation.
% source: results/diagnostics_20260909/roofline_4090.json scored against
% results/kernel_fusion_20260909/e2e/full_4090_sz{128,256,384,496}.json and
% large/large_4090_m{2,3,4}.json, over seven sizes and the six timed tiles.
% Per-size Spearman for the plain model: -0.943, -0.143, +0.143, +0.943 at
% 128 to 496 px and +0.829 at 992, 1488 and 1984 px. Both models pick 32x4
% every time; only at n=496 is that the measured best.
Empirical sharpness pools to a log-log correlation of $+0.943$ against
iteration count over 17 converged instances, but that is Simpson's paradox:
$-0.539$ over the 13 min-cost flow instances, $-0.842$ over the six regular
topology instances and $-0.120$ over the seven grid topology instances. The one
family where the sign does not invert is TV, at $+0.983$ over four instances.
% source: results/diagnostics_20260909/sharpness_4090.json, converged rows only
% Also available for the experiments section: a leakage-free fixed-index
% estimate reaches r=0.620 only after 5120 iterations, median solve 9600.

\subsection{What this paper does not show}
\label{sec:intro:limits}

We show no advantage for language-model-guided search. The arms were not
independent: 336 of 336 fallback-sourced candidates in the model-guided arm are
byte-identical to candidates of the matched random run, and budgets were
unequal at 76 to 136 fresh solver evaluations against 160 per random run.
% source: results/sprint/phase2/analysis/gate1_verdict.md lines 101-106
We show no advantage for joint search over sequential per-axis search:
backend-only, sequential, reverse-sequential and joint selection all land on one
configuration. That catalog diagnostic ran on the RTX 4090 and the RTX 5090, at
$0.392627$ and $0.379070$ penalised holdout seconds, and the 5090 appears
nowhere else in this paper.
% source: results/hardware_codex/REPORT.md lines 74-77 and 102-105
% CAVEAT: those selections are over an exhaustively measured finite catalog and
% are not budget-matched evolutionary arms.
Equation~\ref{eq:factorisation} explains that null instead of hiding it: when
one factor is set by cadence and the other by execution, per-axis search
recovers the joint optimum. Every other result in this paper is measured on two
GPU architectures, the RTX 4090 (sm\_89) and the RTX PRO 6000 Blackwell Max-Q
(sm\_120), and we claim nothing about accelerator portability in general.
% source: results/kernel_fusion_20260909/REPORT.md lines 3-5
Machines were not exclusive during timing. The process-to-process spread of a
re-measured float64 reference cell is $0.66\%$ over 12 processes, spanning
$20.297$ to $20.432$ microseconds per iteration; that figure does not bound the
fusion timings, which move by 4 to $8\%$ between a median-of-3 and a
median-of-9 protocol.
% source: results/tv_ffi_20260906/DECOMPOSITION_REPORT.md lines 5-8 and
% results/finish_20260910/fusion_r9_sz256.log, fusion_r9_sz496.log
The frozen final holdout has never been evaluated.
% source: results/hardware_codex/REPORT.md line 7,
% results/attack_plan_20260906/REPORT.md line 121

% Section 2. Related work.
% Style contract for this draft: no em-dash characters, no italics, plain
% declarative prose, no hedging adverbs.
%
% PROVENANCE RULE. Every number describing OUR measurements carries a
% "% source:" comment naming a file under /home/miria/OptEvolve/v2/results/.
% Nothing here is drawn from atlas/ or from the 2026 manuscript.
%
% CITATION STATUS, READ BEFORE COMPILING. The literature verdicts positioned
% against below come from the four-lane survey recorded in
% docs/ICLR_SUBMISSION_ATTACK_PLAN.md. That survey is prose in this repo and has
% no artifact under results/, so it sits outside the claim-ledger discipline and
% needs a citation-verification pass of its own. Of the keys cited here, only
% Applegate2021Practical, Lu2024HPRLP, ORToolsPDLP2026, Gleixner2021MIPLIB,
% Stellato2020OSQP and Goulart2024Clarabel exist in refs.bib today. These keys
% must be created and checked against the primary source before submission:
% Lu2026cuPDLP, Lu2026cuPDLPx, Lu2026MPAX, PDOT2024, cuOpt2026,
% CarsonHigham2018, Brisebarre2026, CuClarabel2024, LiangYang2026, cuOSQP2020,
% KempkeKoch2025, PushakHoos2022, Hutter2014fANOVA, Xu2008SATzilla,
% BischlASlib2016, Zheng2020Ansor, KernelFoundry2026, Su2012clSpMV,
% KernelBenchVerified2026, Wolf1991TemporalBlocking.
% No third-party numeric claim is printed except the OR-Tools default of 64,
% which the coauthor should confirm against the source file before submission.
%
% COAUTHOR DECISION FLAGGED: the KernelBenchVerified2026 mention in
% subsection 2.5 is a live scoop risk (arXiv:2607.16241, June 2026,
% unpublished). It is cited as concurrent work, and the section leads with the
% factorisation rather than with kernels. Cut that sentence if the coauthor
% prefers no acknowledgement.

\section{Related work}

Two phrases here name different quantities. Arithmetic width is the floating
point format the iterates are stored and updated in. Termination tolerance is
the residual level at which a certificate is accepted. Work titled low precision
sometimes means the second and not the first \citep{KempkeKoch2025}. We report
the two apart throughout.

\subsection{GPU first-order solvers for linear programming}

PDLP established restarted averaged PDHG with presolve, diagonal
preconditioning, adaptive step sizes and a primal weight
\citep{Applegate2021Practical}. The accelerator line ports that method to GPUs:
cuPDLP and cuPDLPx \citep{Lu2026cuPDLP,Lu2026cuPDLPx}, MPAX as a JAX
reimplementation \citep{Lu2026MPAX}, HPR-LP with Halpern iteration and
duplicated dual kernels \citep{Lu2024HPRLP}, PDOT for optimal transport
\citep{PDOT2024}, and cuOpt as a shipped product \citep{cuOpt2026}. Each reports
wall clock for one default configuration and attributes the result to the
method. What is left open is the attribution.
% source: results/tv_ffi_20260906/cadence_4090_v1/factorial.json (f64, device
% loop), f32row_4090/factorial.json (f32_cert64, device loop),
% f64host_4090/factorial.json (schedule, host phases),
% results/bw_20260909/sz256_bw/factorial.json (second architecture). Train fold
% of six scans, data seed 20260907, drawn from a pool of twelve
% (cadence_4090_v1/protocol.json, cadence_4090_v1/data.json). Per-instance
% vectors [1008,512,608,416,320,640], [1024,512,640,448,320,640] and
% [1024,512,768,512,512,768] agree element by element on the two devices.
On total variation denoising at $n=256$ with tolerance $10^{-4}$, over a train
fold of six scans drawn with data seed 20260907 from a pool of twelve,
Condat-Vu with banked steps needs a mean of 584.0, 597.3 and 682.7 iterations
at cadences $K=16$, 64 and 256, and the same triple recurs under both
arithmetic widths, under both execution structures and on a second GPU
architecture, where the per-instance vectors agree element by element.
% source: results/tv_ffi_20260906/kernelhost_4090/factorial.json and
% catalog_4090/factorial.json. Neither archives a cadence-16 cell.
The hand written kernel is archived only at $K=64$ and $K=256$, where it
returns 597.3 and 682.7 like the XLA arm.
% source: results/bw_20260909/sz128_bw/factorial.json and sz496_bw/factorial.json
That invariance is a scope statement and not a law. It covers two GPU
architectures, one problem family, one size and one tolerance, and its width
half already fails one size away: at $K=64$ on the second architecture, float32 iterates
need 586.7 against 576.0 at $n=128$ and 1077.3 against 1066.7 at $n=496$.
% source: results/method_axis_20260909/factorial.json
The method axis moves that count through its step sizes and not its name:
Condat-Vu and PDHG return byte-identical per-instance counts at $K=256$, while
balanced steps cost 4.3125x on iteration count and solve 0 of 6 instances at
$10^{-6}$.

\subsection{Reduced arithmetic width in optimisation}

The accuracy ceiling of a reduced width is a theorem
\citep{CarsonHigham2018,Brisebarre2026}, reduced-width iterates under
full-precision certification already ship in cuOpt and CuClarabel, and cuOSQP
measured the float32 against float64 envelope on GPU and attributed it to memory
boundedness \citep{cuOpt2026,CuClarabel2024,LiangYang2026,cuOSQP2020,Stellato2020OSQP,Goulart2024Clarabel}.
We claim none of that. What is left open is whether width is a per-iteration
speed knob at all.
% source: cadence_4090_v1 (device f64 0.013491 s), f32row_4090 (device f32
% 0.013489 s), f64host_4090 (host f64 0.009042 s),
% px999999_4090/factorial.json (host f32 0.005763 s), all condat_vu_banked at
% K=256, N=682.7, mean of the six train instances
On one device, at one cadence, at $n=256$, float32 buys 1.000x per iteration
inside the traced device loop and 1.569x once the loop is chunked, so the payoff
belongs to the loop structure that the width gene happened to select.
% source: results/tv_ffi_20260906/f32row_4090/factorial.json (4090, K=16, 64
% and 256), results/bw_20260909/sz256_bw/factorial.json (second architecture,
% K=64 only at this size)
The ceiling is also sharper than slower convergence. At $n=256$, Condat-Vu with
float32 iterates and float64 certification solves 0 of 6 instances at
$10^{-6}$, and PDHG solves 1 of 6. On the 4090 that holds at each of the three
archived cadences $K=16$, 64 and 256; on the second architecture only $K=64$ is
archived at this size, and it gives the same 0 of 6 and 1 of 6.

\subsection{Certificate cadence as established engineering}

Checking the certificate less often is a standard implementation choice.
OR-Tools sets \texttt{termination\_check\_frequency} to 64 and annotates it as
extra work \citep{ORToolsPDLP2026}, cuPDLPx captures one check interval as a
replayed CUDA graph, and HPR-LP keeps a second dual kernel to avoid paying for
the check. The field tunes this constant. What is left open is its status.
% source: results/artifacts_20260910/fanova.json, decompositions
% alg_x_prec_x_cadence__logSGM1__tol0.0001 and __tol1e-06, regenerated from
% results/tv_ffi_20260906/cadence_4090_v1/factorial.json. Algorithm 0.0012 and
% 0.3444 percent; cadence 39.4447 and 60.9751; precision by cadence 39.7328 and
% 35.5144. Same values as DECOMPOSITION_REPORT.md section 4.
An fANOVA decomposition \citep{Hutter2014fANOVA} over the twelve-cell algorithm
by width by cadence factorial assigns 0.00\% of the variance to algorithm at
$10^{-4}$ and 0.34\% at $10^{-6}$, against 39.44\% and 60.98\% for the cadence
main effect and 39.73\% and 35.51\% for the width by cadence interaction.
Cadence is the only gene entering both factors of the factorisation.
% The per-iteration loop-structure cost that separates the cadence constant from
% the check itself (11972 ns/iter at 32 outer trips against 10769 ns at one, with
% the certificate term between -240 and +845 ns) is reported in Section 5.
% source: results/finish_20260910/loop_mechanism.json

\subsection{Algorithm selection, portfolios and AutoML nulls}

SATzilla, ASlib and the MIPLIB feature line select among solvers using instance
features \citep{Xu2008SATzilla,BischlASlib2016,Gleixner2021MIPLIB}. What is left
open is whether the selected object carries variance. Scheme identity is inert
to the last iteration here, so a portfolio over scheme names has nothing to
win.
% source: results/diagnostics_20260909/sharpness_4090.json
The obvious replacement feature also fails: in logs, empirical sharpness
correlates with iteration count at a pooled $r=+0.943$ over 17 converged
instances but at $-0.539$ within netflow and $-0.842$ within the regular
topology, a Simpson reversal and not a predictor.
% source: results/netflow_20260909/grid_4090.json,
% results/structure_20260909/structure_summary_4090.json
What replaces it is an exact symbolic side condition: steps that are legal on
total variation give $t s \|L\|^2 = 1.02267$ on netflow and are rejected in 24 of
48 configurations before execution.
% source: results/attack_plan_20260906/analysis/REPORT.md and analysis/summary.json
\citet{PushakHoos2022} report that sequential per-hyperparameter tuning ties
joint tuning in AutoML. Our catalog sweep reproduces that null, at 1.0000x on
CPU at both tolerances and at worst 1.0305x on the 4090, and the factorisation
explains it.

\subsection{Kernel autotuning, per-device selection and temporal blocking}

Per-device schedule search is established by Ansor, clSpMV and KernelFoundry
\citep{Zheng2020Ansor,Su2012clSpMV,KernelFoundry2026}, and temporal blocking,
which fuses several time steps into one pass over a tile, is standard in stencil
computing \citep{Wolf1991TemporalBlocking}. Concurrent work revisits
generated-kernel benchmarking \citep{KernelBenchVerified2026}. We claim neither
as new. What is left open is the shape of the payoff and the cost of guessing
it.
% source: results/kernel_fusion_20260909/e2e/full_4090_sz{128,256,384,496}.json
% and results/kernel_fusion_20260909/e2e_bw/full_bw_sz{128,256,384,496}.json
The wrong tile is worse than not fusing: at $n=496$ on the 4090 the best tile
gives 1.356x end to end and the worst gives 0.743x, and the two devices choose
different tiles at $n=384$ and $n=496$.
% source: results/kernel_fusion_20260909/large/large_4090_m2.json (992 px),
% m3.json (1488 px), m4.json (1984 px); 128 px from e2e/full_4090_sz128.json
The gain is U-shaped in problem size, at 1.571x for 128 px, 1.060x at 992 px and
1.410x at 1984 px, which is launch overhead removed at one end and HBM traffic
at the other.
% source: results/diagnostics_20260909/roofline_4090.json, model M1 recomputed
% against the e2e and large sweeps
A cheap a-priori model does not substitute for the search: roofline with
arithmetic intensity reaches mean Spearman 0.355 over seven sizes, picks the
best tile 1 time in 7, and is anti-correlated at $-0.943$ at 128 px.

\begin{table}[t]
\centering
\caption{Position of this work against the nearest lines. The right column
lists only measurements reported in Section~5. Every iteration count in this
table is a mean over the same train fold of six scans, data seed 20260907,
drawn from a pool of twelve.}
\label{tab:related}
\small
\begin{tabular}{@{}p{0.23\linewidth}p{0.34\linewidth}p{0.35\linewidth}@{}}
\toprule
Line of work & What it establishes & What is left open, and what we measure \\
\midrule
% source: results/tv_ffi_20260906/cadence_4090_v1, f32row_4090, f64host_4090,
% results/bw_20260909/sz256_bw
GPU first-order LP \citep{Applegate2021Practical,Lu2024HPRLP,cuOpt2026}
& Wall clock for one default configuration per solver
& At $n=256$, iteration count is unchanged across the two measured devices, both widths and both execution structures: 584.0 / 597.3 / 682.7 at $K=16$ / 64 / 256 \\
% source: results/tv_ffi_20260906/{cadence_4090_v1,f32row_4090,f64host_4090,px999999_4090}
Reduced width \citep{CarsonHigham2018,CuClarabel2024,cuOSQP2020}
& Reduced-width iterates with full-precision certification, and an accuracy ceiling theorem
& Width buys 1.000x per iteration in a traced loop and 1.569x in a chunked one, same device and cadence \\
% source: results/artifacts_20260910/fanova.json, regenerated from
% results/tv_ffi_20260906/cadence_4090_v1/factorial.json
Cadence tuning \citep{ORToolsPDLP2026}
& The constant is configurable and is tuned per solver
& It is the only gene in both factors, at 39.44\% and 60.98\% of variance against 0.00\% and 0.34\% for algorithm \\
% source: results/method_axis_20260909, results/netflow_20260909/grid_4090.json
Portfolios and features \citep{Xu2008SATzilla,BischlASlib2016}
& Instance features select among solvers
& Scheme identity is inert; an exact symbolic gate rejects 24 of 48 netflow configurations before execution \\
% source: results/attack_plan_20260906/analysis/summary.json
AutoML nulls \citep{PushakHoos2022}
& Sequential tuning ties joint tuning
& The same null, at 1.0000x to 1.0305x, explained by which axes fail to interact \\
% source: results/kernel_fusion_20260909/e2e, /large, results/diagnostics_20260909
Autotuning and temporal blocking \citep{Zheng2020Ansor,Su2012clSpMV}
& Per-device schedule search, time-step fusion in stencils
& U-shaped gain in size, 0.743x for the wrong tile, roofline picks the best tile 1 of 7 \\
\bottomrule
\end{tabular}
\end{table}

% SCOPE GUARDS. Keep these in the checklist when the section is edited:
% (1) two GPU architectures only, so no accelerator-portability claim is made
%     anywhere in this section;
% (2) no LLM-proposal advantage is claimed, here or elsewhere: the arms shared a
%     proposal stream with 8 of 8 byte-identical draws;
% (3) no joint-search advantage over sequential selection is claimed;
% (4) no 1.41x figure appears, since it sits below the published cuOSQP float32
%     against float64 envelope;
% (5) the 584 / 597 / 683 triple is stated for condat_vu_banked only. It is NOT
%     scheme-invariant at K=16, where pdhg_banked gives 578.67
%     (results/tv_ffi_20260906/cadence_4090_v1/factorial.json). The text says
%     "Condat-Vu with banked steps" for that reason;
% (5a) the triple is NOT claimed across kernels. kernelhost_4090 and catalog_4090
%     archive the custom kernel at K=64 and K=256 only, so no archived cell
%     gives 584.0 under a hand written kernel. The text says so;
% (6) width invariance of N is scoped to n=256. It fails on the Blackwell at
%     n=128 (576.0 against 586.7) and n=496 (1066.7 against 1077.3):
%     results/bw_20260909/sz128_bw/factorial.json and sz496_bw/factorial.json;
% (7) the 0 of 6 float32 result at 1e-6 is condat_vu-specific; pdhg_banked
%     solves 1 of 6. Both are stated above, and both are scoped to n=256: the
%     4090 archives K=16, 64 and 256 (f32row_4090) and the Blackwell archives
%     K=64 only at this size (bw_20260909/sz256_bw). The text says so;
% (8) omega0 and the step-ratio rule are NOT cited in this section. The
%     PDLP-rule regret figures did not reproduce from the saved artifacts and
%     are forbidden until scripts/diagnostics scoring is committed under
%     results/.

% Section 3. Method: recipes, the gate, and the execution layer.
% Every number carries a "% source:" comment naming a file under
% /home/miria/OptEvolve/v2/results/. Nothing here comes from atlas/.
% Style contract: no em-dash characters, no italics, plain declarative prose.

\section{Method: recipes, the gate, and the execution layer}
\label{sec:methods}
% sec:methods, sec:gate, sec:mixing and eq:composite are referenced by
% 4-theory.tex, which is carried over unchanged; the labels are placed here
% so those cross-references resolve to the corresponding new material.

A recipe is the complete executable object our search evaluates: a splitting
scheme, its step sizes, its wrappers, a precision policy, a certificate cadence,
a kernel, and an execution structure. Every block is typed, is validated before
any arithmetic runs, and is serialised with each measurement, so a reported time
is attributable to a specific object rather than to a solver name.

The compiler accepts supported instances of the composite form
\begin{equation}
  \min_{x}\; f(x) + g(x) + h(Lx),
  \label{eq:composite}
\end{equation}
with convex $L_f$-smooth $f$, proximable $g$ and $h$, and a linear map $L$ with a
computable adjoint and a supplied upper bound $\ell \geq \|L\|$. For anisotropic
total-variation denoising $f(x) = \tfrac12\|x-y\|^2$, $g = 0$, $h$ is the scaled
$\ell_1$ norm and $L$ is the discrete gradient. For the three linear-programming
families $f$ is the linear cost, $g$ is the indicator of the nonnegative orthant,
$h$ is the indicator of an affine equality set and $L$ is the constraint operator.

\subsection{The recipe}

% source: results/method_axis_20260909/factorial.json, first row of tolerances["0.0001"]
One record reads, verbatim from the results file,
\verb!{scheme: condat_vu, params: {tau: 0.25, sigma: 0.25, alpha: null, rho: 1.5},!
\verb!wrapper: {halpern: false, restart: none, restart_k: null}, mix: null,!
\verb!backend: {precision: schedule, K: 256, theta: 0.01, batch: false}, metric: null}!.
Every row carries the same block structure. Table~\ref{tab:recipe} gives the
blocks and the values the corpus exercises.

% source: results/attack_plan_20260906/source_v1/atlas/backend/CONTRACT.md, section 1 table.
% CAVEAT: that contract file is a 2026-09-06 snapshot. It predates the
% execution-structure gene and still describes precision as the gene that
% selects the execution policy. The prose below says so explicitly.
\begin{table}[t]
\centering
\caption{Recipe blocks and the values exercised in the recorded corpus.
Validation rejects with the violated condition and never clamps.}
\label{tab:recipe}
\begin{tabular}{@{}lll@{}}
\toprule
Block & Meaning & Recorded values \\
\midrule
scheme    & splitting identity        & \texttt{pdhg}, \texttt{condat\_vu} \\
params    & step sizes, relaxation    & \texttt{tau}, \texttt{sigma}, \texttt{alpha}, \texttt{rho} \\
wrapper   & anchoring and restart     & \texttt{halpern}, \texttt{restart}, \texttt{restart\_k} \\
precision & iterate arithmetic policy & \texttt{f64}, \texttt{f32\_cert64}, \texttt{schedule} \\
K         & certificate cadence       & integer, $1 \le K \le 10000$ \\
theta     & schedule switch threshold & float, $0 < \theta < 1$, schedule only \\
batch     & reserved population backend & bool, validates and is a no-op \\
kernel    & atom implementation       & \texttt{jax}, \texttt{tv\_fused128}, \texttt{tv\_fused256}, \texttt{tv\_recip256} \\
execution & compiled loop structure   & \texttt{device\_loop}, \texttt{host\_precision\_phases} \\
\bottomrule
\end{tabular}
\end{table}

The precision enum admits three values and rejects any spelling of
\texttt{bf16} or \texttt{f16} by name. The cadence gene $K$ is the interval at
which the float64 duality gap is evaluated, and the recorded evaluation points
are exactly the multiples of $K$: a solve at $K=256$ that stops after 4864
iterations records \texttt{gap\_iters} $= [256, 512, \ldots, 4864]$ and 19
checks.
% source: results/method_axis_20260909/factorial.json (gap_iters, certificate_checks);
% cadences 64,128,256,512 also in results/kernel_fusion_20260909/cadence/cad_4090_K*.json
The \texttt{schedule} policy is a two-phase wrapper, float32 iterates then
float64, with the switch governed by $\theta$ and its iteration recorded:
\texttt{iterate\_dtypes} $=[\text{float32},\text{float64}]$,
\texttt{switch\_iter} $=256$, $\theta = 0.01$.

% source: results/tv_ffi_20260906/DECOMPOSITION_REPORT.md, section 1
Execution structure is one XLA-visible difference. The device loop is
\texttt{jax.lax.fori\_loop(k, end, ...)} with a traced bound inside a
\texttt{while\_loop}, which XLA cannot unroll. The chunked structure is
\texttt{jax.lax.fori\_loop(0, K, ...)} with a static Python $K$, which XLA
unrolls and fuses across iterations. In the shipped catalog the choice is not
separately addressable, because it is dispatched from the precision gene:
\texttt{f64} and \texttt{f32\_cert64} always ran the device loop and
\texttt{schedule} always ran host precision phases.
% source: enumerated over all 15 factorial.json files in results/tv_ffi_20260906
% and results/method_axis_20260909, train and validation rows together: 2208
% (f64, device_loop), 144 (f32_cert64, device_loop), 1680 (schedule,
% host_precision_phases). Train rows alone are 1104, 72 and 840. No
% cross-pairing exists.
Enumerating the pairs over all fifteen factorial artifacts confirms that no
recorded cell separates them. We therefore promote execution structure to a
gene of its own and report it explicitly.

\subsection{The construction gate}
\label{sec:gate}

The gate is an exact-arithmetic side condition on the step sizes, checked
symbolically before execution. A rejected configuration carries the violated
condition and no timing at all.
% source: results/netflow_20260909/grid_4090.json, alg=pdhg_balanced and condat_vu_balanced rows
PDHG rejects with
\verb!side condition violated: t>0, s>0, t*s*||L||^2 < 1  (t*s*||L||^2=1.02267)!
and Condat-Vu with
\verb!side condition violated: t*L_f/2 + t*s*||L||^2 < 1  (lhs=1.02267)!.
On a linear cost the Lipschitz term vanishes and the two left-hand sides
coincide at every size, taking 1.02267, 1.03040, 1.03577 and 1.03781 at grid
sizes 12, 16, 24 and 32.

% source: results/structure_20260909/structure_summary_4090.json, 13 cells.
% norm_bound_gate/sigma_max recomputed per cell: minimum 1.0191856 at
% regular_s1024_d3, maximum 1.0200000 at grid_s12. ts_max_admissible times
% norm_bound_gate squared evaluates to 1.0 in float64 in all thirteen cells.
The gate consumes a norm bound, not the norm. The ratio of
\texttt{norm\_bound\_gate} to $\sigma_{\max}$ is 1.0200000, 1.0199999,
1.0199124 and 1.0199594 at grid sizes 12, 16, 24 and 32, and over all thirteen
cells it runs from 1.0191856 to 1.0200000. The admissible step product equals
the reciprocal of the squared bound exactly in double precision in every
cell. The admissible region is a measured function of family and degree
and is close to size-invariant within a family (Table~\ref{tab:gate}).

\begin{table}[t]
\centering
\caption{The admissible step product is set by problem structure. Per cell,
three instances each.}
\label{tab:gate}
% source: results/structure_20260909/structure_summary_4090.json, cells array
\begin{tabular}{@{}llll@{}}
\toprule
Family & Degree & $\sigma_{\max}$ range & Admissible $ts$ range \\
\midrule
grid    & 4 & 3.9658 to 3.9952 & 0.060223 to 0.061114 \\
regular & 2 & 2.828421 to 2.828427 & 0.120197 to 0.120208 \\
regular & 3 & 3.4116 to 3.4138 & 0.082559 to 0.082663 \\
regular & 4 & 3.8604 to 3.8630 & 0.064435 to 0.064535 \\
\bottomrule
\end{tabular}
\end{table}

% source: results/netflow_20260909/grid_4090.json (24 of 48 rows carry an "error" key)
% and results/method_axis_20260909/factorial.json (N and penalised SGM-1 ratios)
Family dependence is a difference in kind. The balanced preset
$(\tau,\sigma,\rho) = (0.25, 0.25, 1.5)$ is inadmissible on netflow, rejected in
24 of 48 configurations, while the same preset is legal on total-variation
denoising and merely slow there, costing 4.3125 times the iteration count of the
banked preset $(0.05, 1.0, 1.9)$.

% source: results/finish_20260910/{omega_4090,omega_ot_4090,omega_bw}.json,
% 972+729+972 = 2673 rows, every one requested at theta = 0.99. Recomputed
% tau*sigma*Lnorm^2 equals 0.99 exactly in 882 rows and differs by at most
% 3.331e-16 elsewhere; recomputed tau/sigma equals r exactly in 1575 rows and
% differs by at most 2.4e-16 elsewhere.
Reparameterising $(\tau, \sigma)$ as a saturation level $\theta = t s \|L\|^2$
and a ratio $r = \tau/\sigma$ makes the gate a function of $\theta$ alone.
Across all 2673 rows of the three sweeps, every one requested at
$\theta = 0.99$, the recomputed product agrees with 0.99 to within
$3.4\times10^{-16}$ and is exact in 882 rows, and the recomputed ratio agrees
with $r$ to within $2.4\times10^{-16}$ and is exact in 1575 rows.
% source: results/structure_20260909/structure_summary_4090.json,
% gate_is_function_of_theta_only block
% source: results/finish_20260910/omega_ot_4090.json, Lnorm constant within each cell
On optimal transport the bound needs no power iteration: the recorded norm is
8.160000, 11.539983 and 14.133535 at the $32\times32$, $64\times64$ and
$96\times96$ cells, which are $1.02\sqrt{m+n}$ to full double precision, so that
family's admissible region shrinks as the instance grows.

\subsection{Three assurance levels}
\label{sec:mixing}

% source: results/attack_plan_20260906/source_v1/atlas/backend/CONTRACT.md, section 3
The backend contract already states that construction assumptions,
implementation tests and runtime numerical checks are separate obligations, and
that changing dtype changes numerical errors so exact-arithmetic convergence
guarantees do not transfer verbatim to mixed precision. We keep the three
separate and never let one stand in for another.

\textbf{Level 1, construction.} The side condition is checked in exact
arithmetic before execution and recorded per run as
\texttt{construction\_satisfied}; on netflow, 24 rows are gate verdicts with no
measurement attached.
% source: results/netflow_20260909/grid_4090.json and
% results/method_axis_20260909/factorial.json

\textbf{Level 2, implementation admission.} A compiled update must reproduce a
float64 reference before use. Every kernel variant passes five per-atom checks
against float64 numpy references at its declared tolerance before entering any
phenotype, the generated LP CSR templates additionally pass an eight-trial
registry check and then per-instance reference-map admission, and non-finite
outputs fail admission. This level is empirical: numerical admission and
reference-trajectory agreement are not formal floating-point proofs, and the
reciprocal kernel variant deliberately changes rounding.
% source: CONTRACT.md current update and section 3; the limitation sentence is
% published verbatim in results/tv_ffi_20260906/analysis_v2/REPORT.md, Limits.

\textbf{Level 3, runtime certification.} The duality gap is evaluated in
float64 in original coordinates and termination is decided by that value alone.
The certificate dtype is float64 in all 2016 recorded per-instance train runs
across the fifteen factorial artifacts, and the profiled operator set names
\texttt{original\_forward} and \texttt{original\_adjoint} beside
\texttt{base\_step\_f64}, \texttt{base\_step\_f32} and \texttt{certificate\_f64}.
Enforcement is mechanical: the certified gap raises on any non-float64
certificate value, and the policy constructor refuses \texttt{bf16} and
\texttt{f16} by name.
% source: results/method_axis_20260909/profiles.json and factorial.json; CONTRACT.md section 3

\subsection{Hardware-aware backend genes}
\label{sec:hardware}
% [RESTORED VERBATIM FOR REVIEW BY NIKHIL, 2026-09-10]
% Source: github.com/miria00/OptAtlas_2026, main at commit 218ac81
% (2026-08-31), 3-methods.tex lines 142-172. Only LaTeX comment lines were
% added here; the rendered text of this subsection is identical to that commit.
% A different, never-pushed version titled "Hardware-aware execution" exists in
% v2/OptEvolve_tex_2026/3-methods.tex lines 85-109 and is NOT used here.
%
% PRECISION-FRAMING FLAGS [P1] [P2] [P3] mark sentences for review; the text
% itself is unedited. Common cause: in the shipped catalog the precision gene
% and the execution-structure gene are the same gene.
% atlas/schemes/base_schemes.py:813 sends f64 and f32_cert64 to the
% traced-bound device loop and only schedule to chunked host phases, so no
% catalog cell varies precision while holding loop structure fixed. Separated
% by hand (TV n=256, RTX 4090, tol 1e-4, K=256, cost per iteration): float32
% over float64 is 1.00x inside the device loop and 1.58x inside chunked phases.
% source: results/artifacts_20260910/fanova.json (per_iteration_us) and
% results/tv_ffi_20260906/DECOMPOSITION_REPORT.md sections 1 and 2.

% [P1] Names floating-point precision as a backend gene alongside batching,
%      sparse format and kernel fusion. In the shipped catalog precision is not
%      independent of loop structure; see the header note above.
A second family of genes controls not what iteration runs but
how it is computed: floating-point precision, batching across
instances, sparse-matrix format, and kernel-fusion strategy. The
% [P2] Flags "backend genes cannot affect correctness ... fusion, precision, and
%      block sizes change only its implementation." This holds for the
%      exact-arithmetic certificate on the map T. Numerically, precision is not
%      implementation-only: pure float32 iterates certify 0 of 6 (Condat-Vu) and
%      1 of 6 (PDHG) at tol 1e-6 where float64 certifies 5 of 6 (TV n=256, RTX
%      4090). Float32 also changes the iteration count at n=128 (float64 576.0,
%      float32 586.7) and n=496 (float64 1066.7, float32 1077.3), K=64, RTX PRO
%      6000 Blackwell.
%      source: results/tv_ffi_20260906/f32row_4090/factorial.json and
%      results/bw_20260909/sz128_bw, results/bw_20260909/sz496_bw.
architectural invariant is that backend genes cannot affect correctness at
the operator level: the construction certificate is a statement about the
map $T$, and fusion, precision, and block sizes change only its
implementation. Kernel correctness is instead established by per-atom unit
tests that require agreement with a reference backend to tolerance, in the
pattern of hardware-aware kernel generation for coding agents
\citep{Hawkeye2026}. The system thus has two free-form-but-verified layers
surrounding a certified core: heuristic schedules above the calculus,
checked by the safeguard monitor, and kernels below it, checked by
unit tests; the calculus between them is never mutated.

The target hardware is visible to the search: the device identifier
conditions the mutation context, so genomes can specialize to the
accelerator they will run on. Kernel fusion matters most exactly where
per-iteration work is small and launch overhead dominates, i.e.\ the
microsecond-budget cells; batched execution matters where many instances
% [P3] Flags "This paper ships the precision ... genes." The shipped precision
%      gene cannot be varied without also changing loop structure (header note),
%      so a catalog effect attributed to it is jointly a precision and a
%      loop-structure effect. The best measured configuration, float32 iterates
%      with chunked execution at K=256, is 2.11x over compiled float64 at K=64 and
%      is not expressible in the shipped genome.
%      source: results/tv_ffi_20260906/DECOMPOSITION_REPORT.md section 6.
are solved simultaneously. This paper ships the precision, batching, and
fused-versus-naive step genes; autogenerated fused kernels for the atoms are
future work.

\iclrTODO{Cross-hardware divergence table: evolve on two GPU architectures
with identical instance distributions and show the winning genomes differ
(consumer vs.\ data-center parts; FP64 throughput and memory bandwidth
drive different certified optima).}

\subsection{The temporal fusion kernel}

% source: results/kernel_fusion_20260909/tv_temporal.cu
The execution layer contains one construction no catalog gene expresses. The
fused kernel runs $T$ iterations per launch with the tile resident in shared
memory. For an output tile of edge $B$ it loads a shared tile of edge
$S = B + 2T$, a halo of width exactly $T$ on every side. Iteration \texttt{it}
computes the primal on $[\texttt{it}+1,\, S-\texttt{it})$ and the dual and the
commit on $[\texttt{it}+1,\, S-\texttt{it}-1)$, so the valid region shrinks by
exactly one cell per side per iteration. After the final iteration, at
$\texttt{it} = T-1$, the committed region is $[T,\, B+T)$, and the write-back
reads local index $T + o$ for $o \in [0, B)$. The interval closes exactly, with
no slack and no redundant computation. Each fused iteration is three
barrier-separated stages, a primal predict into scratch, a dual clip with an
in-place relaxation, and a primal commit; stages two and three branch on
$\rho = 1$ and write the unblended value, so at $\rho=1$ the blend never
executes and no extra rounding accumulates.

% source: build.py comment (5*(B+2T)^2*8 bytes) and the five shared pointers in
% tv_temporal.cu; S and redundancy cross-checked in
% results/diagnostics_20260909/roofline_4090.json. 64000 bytes at 32x4 is also
% stated in results/kernel_fusion_20260909/REPORT.md.
Shared memory holds five double arrays over the halo tile, $5(B+2T)^2 \cdot 8$
bytes, which caps the tile size; tiles above 48\,KB request the opt-in
carve-out. Table~\ref{tab:tiles} gives the measured geometry. Halo redundancy is
$S^2/B^2$ and the per-iteration traffic reduction is $T$ divided by it, so the
$24\times4$ tile runs four iterations at 1.7778 times the loads and moves 2.25
times less data per iteration.

\begin{table}[t]
\centering
\caption{Geometry of the measured tiles. Shared use is $5(B+2T)^2 \cdot 8$
bytes; redundancy is $S^2/B^2$.}
\label{tab:tiles}
% source: results/diagnostics_20260909/roofline_4090.json for B, T, S, redundancy.
% Shared bytes computed from the build.py formula.
\begin{tabular}{@{}lrrrrr@{}}
\toprule
Tile & $B$ & $T$ & $S$ & Shared bytes & Redundancy \\
\midrule
$8\times4$  &  8 & 4 & 16 & 10240 & 4.00000 \\
$16\times2$ & 16 & 2 & 20 & 16000 & 1.56250 \\
$16\times4$ & 16 & 4 & 24 & 23040 & 2.25000 \\
$24\times4$ & 24 & 4 & 32 & 40960 & 1.77778 \\
$32\times2$ & 32 & 2 & 36 & 51840 & 1.26563 \\
$32\times3$ & 32 & 3 & 38 & 57760 & 1.41016 \\
$32\times4$ & 32 & 4 & 40 & 64000 & 1.56250 \\
\bottomrule
\end{tabular}
\end{table}

% source: results/kernel_fusion_20260909/endtoend.py, "if args.K % T: continue";
% solve_chunked calls _run_phase with K//F and max_iters//F and reports k*F.
Fusion depth is coupled to cadence by a hard divisibility constraint: $T$ must
divide $K$, because the fused kernel replaces $T$ iterations of the chunk runner
and the chunk therefore takes $K/T$ steps. Measured tiles have
$T \in \{2,3,4\}$ and measured cadences lie in $\{64,128,256,512\}$. This is the
point at which the fusion choice and the certificate cadence stop being
independent. Portability across the two architectures is by forward JIT of
\texttt{compute\_80} PTX rather than by recompilation.
% source: build.py compile flags -gencode=arch=compute_80,code=compute_80
% CAVEAT for the coauthor: the temporal tile is NOT a gene in the shipped
% genome. endtoend.py builds the kernels itself and calls
% atlas.backend.precision._run_phase directly; the contract snapshot lists only
% jax, tv_fused128, tv_fused256, tv_recip256 and states the Stage-0.5 kernel gate
% "remains a hook only". Never assert that the tile is a gene.
The tile is not currently expressible as a gene. The measurements below drive it
from a harness, and the contract lists no temporal variant among the kernel
values. We argue it should become one.

\subsection{Numerical admission of the fused kernel}

% source: results/kernel_fusion_20260909/validate.py
Admission compares one fused launch against $T$ sequential launches of the
single-step reference kernel \texttt{tv\_fused128} over a declared grid: sizes
$64\times64$, $128\times128$ and $100\times130$, both the proximal and the
gradient form, and $\rho \in \{1, 1.9\}$, at $t=0.05$, $s=1.0$, $\lambda=0.1$
and $\text{inv}=1/(1+t)$. The size $100\times130$ is not a multiple of $B$, so
the partial-tile write-back guard is exercised. The pass criterion is that the
largest absolute primal and dual deviations fall below $10^{-12}$, with bitwise
equality reported separately.

% source: results/artifacts_20260910/kernel_admission_4090.log, 106 lines,
% 102 printed cases, all at rho=1.9, over 17 tile shapes, sizes 64x64, 128x128
% and 100x130 and both proximal settings, worst absolute deviation 4.649e-15,
% PASS. The worst is reached at tv_tf8x8t64 and tv_tf16x8t128 on 100x130, both
% T=8, log lines 13 and 55. validate.py prints a case only when it deviates by
% more than 1e-12 or is not bitwise equal, so the rho=1 half of the same grid,
% 102 further cases, is absent from the log precisely because every one of them
% was bitwise equal. No Blackwell admission log exists under results/, so no
% sm_120 figure is printed here.
The archived run scores 17 tile variants, a superset of the timed tiles of
Table~\ref{tab:tiles} that reaches $T=8$. Its worst absolute deviation is
$4.649\times10^{-15}$, at $\rho=1.9$, reached on the two variants that fuse
eight iterations, and it passes the declared $10^{-12}$ criterion on all 102
printed cases. The script prints a case only when it is not bitwise equal to
the sequential reference, and every printed case is at $\rho=1.9$, so at
$\rho=1$ the fused launch reproduced the reference bit for bit in all 102 cases
of that half of the grid. This record is the 4090 only.
No admission run on the second architecture is archived, and none is claimed.

% source: computed across the four files in results/kernel_fusion_20260909/e2e/
% and the four in e2e_bw/. Worst relative gap disagreement per size, identical
% on the two devices: 2.608e-11 at n=128, 5.655e-12 at 256, 1.224e-11 at 384,
% 5.111e-12 at 496. The fusion harness as run drew its six scans with seed 0, a
% different population from the catalog train fold; the manifest and the
% reconciliation are in results/artifacts_20260910/A1_RESOLUTION.md.
Archived data corroborates admission end to end: at every size the eight arms,
two baseline execution structures and six fused tiles, return identical
iteration counts on all six instances of the fusion harness draw, six scans at
data seed 0, and at $n=128$ the worst relative disagreement between certified
gaps is $2.608\times10^{-11}$. The timed arms run the gradient form with the
reciprocal kernel variant at $\rho=1.9$, the branch where agreement is
$4.649\times10^{-15}$ at worst rather than bitwise.
% source: results/kernel_fusion_20260909/endtoend.py: call(tgt, ..., proximal=0,
% rho=1.9, d["tv_inverse"]), key=(B,T,True) selects Reciprocal=True, timing is
% the median of `repeats` runs after one discarded warmup, from _run_phase.

\section{Convergence by construction}
\label{sec:theory}

For the base-and-relaxation grammar specified below, the Stage-0 gate of
Section~\ref{sec:gate} establishes sufficient exact-arithmetic convergence
conditions. Definition~\ref{def:certificate} defines the mathematical
certificate; Theorem~\ref{thm:soundness} and Corollary~\ref{cor:convergence}
give the guarantee under correct atom implementations. Numerical termination
and implementation admission are separate checks, not consequences of this
theorem.
Appendix~\ref{app:theory} gives the supporting lemmas and proofs.

\paragraph{Setup and notation.}
Throughout, $f$ is convex and differentiable with an $L_f$-Lipschitz
gradient, $g$ and $h$ are closed, proper, and convex, all spaces are
finite-dimensional, and the saddle (Karush--Kuhn--Tucker) set of
problem~\eqref{eq:composite} is nonempty. The gate is supplied with a
constant $\ell\geq\|L\|$ and checks inequalities using $\ell$. These conditions
are sufficient for the corresponding statements about $L$. Using the exact
norm removes conservatism in that bound but does not make the convergence
condition necessary. Let $\mathcal X^\star$ denote the primal solution set. For
$H\succ0$, define
$\|z\|_H^2=\langle z,Hz\rangle$ and
$\operatorname{dist}_H(z,C)=\min_{c\in C}\|z-c\|_H$. An operator $T$ is
\textbf{$\theta$-averaged in $\|\cdot\|_H$} for $\theta\in(0,1)$ if
$T=(1-\theta)I+\theta S$ for some $S$ that is nonexpansive in $\|\cdot\|_H$
\citep[\S2.3 and \S2.8]{RyuYin2022LargeScale}. A genome $G$ is a term of
the grammar $\mathcal G$ defined below. The subscript $G$ identifies data
attached to the genome: its iteration $T_G$, its state
$z\in\mathcal Z_G$, and its certificate. The state is a primal variable, a
Douglas--Rachford auxiliary variable, or a primal--dual pair.

\begin{definition}[Construction certificate]
\label{def:certificate}
A \textbf{construction certificate} for a compiled iteration
$T_G:\mathcal Z_G\to\mathcal Z_G$ is a triple $(H_G,\theta_G,\pi_G)$ consisting
of a metric $H_G\succ0$, an averagedness constant $\theta_G\in(0,1)$, and a
decoder $\pi_G$, such that
(C1) $T_G$ is $\theta_G$-averaged in $\|\cdot\|_{H_G}$;
(C2) $\operatorname{Fix}T_G\neq\emptyset$; and
(C3) $\pi_G(\operatorname{Fix}T_G)\subseteq\mathcal X^\star$.
\end{definition}

Conditions (C1) and (C2) concern the iteration. Condition (C3) specifies how a
fixed point is decoded as a solution. The certificate does not depend on an
implementation or a numerical run, so the gate can verify it symbolically.

\paragraph{The certified grammar.}
A genome is a term $G$ of the grammar
\begin{equation}
 \mathcal G:\qquad
 G \;::=\; \mathrm{base}_s(p)
 \;\mid\; \mathrm{relax}_{\rho}(G') ,
 \label{eq:grammar}
\end{equation}
Here $\mathrm{base}_s(p)$ selects one of the five base schemes in
Table~\ref{tab:certificates}, $s\in\{\mathrm{fbs},\mathrm{drs},
\mathrm{dys},\mathrm{pdhg},\mathrm{cv}\}$, together with its stepsizes $p$
($\alpha$ or the pair $(\tau,\sigma)$). The term
$\mathrm{relax}_\rho(G')$ over-relaxes a genome $G'$ by a weight $\rho>0$.
The champion of Section~\ref{sec:results} is
$\mathrm{relax}_{1.9}\bigl(\mathrm{base}_{\mathrm{cv}}(0.05,1.0)\bigr)$.
The mathematical cores in the hardware catalogs use this base-and-relaxation
grammar. Rescaling, numerical implementations and wrapper extensions have
additional obligations and are not automatically covered by this theorem.
Averaged blending transports certificates under the hypotheses of
Lemma~\ref{lem:blend}; the implementation admits only restricted parents for
which metric and fixed-point compatibility are established by the construction.
The broader derivation-tree transformations are outside the implemented grammar
and are not claimed here.

\begin{theorem}[Gate soundness]
\label{thm:soundness}
Every genome $G$ generated by the grammar $\mathcal G$ of
\eqref{eq:grammar} and accepted by the Stage-0 gate carries a construction
certificate $(H_G,\theta_G,\pi_G)$. Specifically, (i) for
$\mathrm{base}_s(p)$ the certificate is given in Table~\ref{tab:certificates},
with $\theta_G=1/(2-c_G)$, and the checked
inequality implies $H_G\succ0$ and $c_G<1$;
(ii) for $\mathrm{relax}_\rho(G')$ the gate checks $\rho\,\theta_{G'}<1$, and
the certificate is $(H_{G'},\,\rho\,\theta_{G'},\,\pi_{G'})$.
\end{theorem}

\begin{proof}
The result follows by structural induction on $G$. The base case is
Lemma~\ref{lem:base-certificates}, and the inductive step is
Lemma~\ref{lem:relaxation}. Each lemma uses the corresponding condition
checked by the gate.
\end{proof}

\begin{table}[t]
\centering
\small
\begin{tabular}{lccc}
\toprule
base scheme & metric $H_G$ & gate inequality & $c_G$ \\
\midrule
FBS & $I$ & $0<\alpha<2/L_f$ & $\alpha L_f/2$ \\
DRS & $I$ & $\alpha>0$ & $0$ \\
Davis--Yin & $I$ & $0<\alpha<2/L_f$ & $\alpha L_f/2$ \\
PDHG & $M_{\tau\sigma}$ & $\tau\sigma\ell^2<1$ & $0$ \\
Condat--V\~u & $M_{\tau\sigma}$ & $\tau L_f/2+\tau\sigma\ell^2<1$
  & $\dfrac{\tau L_f}{2\left(1-\tau\sigma\ell^2\right)}$ \\
\bottomrule
\end{tabular}
\caption{The certificates the gate emits for the base schemes, with
$M_{\tau\sigma}=\left[\begin{smallmatrix}\tau^{-1}I&-L^{\!\top}\\
-L&\sigma^{-1}I\end{smallmatrix}\right]$ and $\ell\geq\|L\|$ the supplied
operator-norm bound. In each row the averagedness constant is
$\theta_G=1/(2-c_G)$ and, for positive stepsizes, the gate inequality implies
$H_G\succ0$ and $c_G<1$. The forward-block smoothness constant belongs to
the accepted decomposition, including a dual decomposition where applicable.
These are sufficient conditions. The
two entries with
$c_G=0$ are the schemes with no forward step, which are firmly nonexpansive
in their metric. PDHG is the $L_f=0$ specialization of Condat--V\~u.}
\label{tab:certificates}
\end{table}

\begin{corollary}[Certified genomes converge, at a rate]
\label{cor:convergence}
Let $T_G$ carry a certificate $(H_G,\theta_G,\pi_G)$. Then the compiled
iteration
\begin{equation}
  z_{k+1}=T_Gz_k
  \label{eq:compiled-iteration}
\end{equation}
converges, from any $z_0$, to some $z^\star\in\operatorname{Fix}T_G$, and
$\pi_G(z^\star)$ solves problem~\eqref{eq:composite}. Moreover
$\|z_{k+1}-z_k\|_{H_G}$ is nonincreasing in $k$ and, for every $k\geq0$,
\begin{equation}
 \|z_{k+1}-z_k\|_{H_G}^2
 \;\leq\;
 \frac{\theta_G}{(1-\theta_G)\,(k+1)}\;
 \operatorname{dist}_{H_G}^2\!\left(z_0,\operatorname{Fix}T_G\right).
 \label{eq:residual-rate}
\end{equation}
\end{corollary}

\begin{proof}
By (C1) and (C2), Proposition~\ref{prop:km} in
Appendix~\ref{app:theory} applies to $T_G$. Condition (C3) shows that the
limit decodes to a solution.
\end{proof}

The rate applies to the theoretical step residual
$\|z_{k+1}-z_k\|_{H_G}$. The numerical TV and LP stopping criteria in
Section~\ref{sec:methods} are different quantities; the rate does not directly
bound their time to tolerance. For a genome
$\mathrm{relax}_\rho(G')$, this residual equals
$\rho\,\|T_{G'}z_k-z_k\|_{H_G}$. Thus, dividing
\eqref{eq:residual-rate} by $\rho^2$ bounds the unrelaxed fixed-point
residual (Remark~\ref{rem:monitored-residual}). Averagedness makes this
residual monotone, so the bound holds for the last iterate. The order in
\eqref{eq:residual-rate} is sharp.
Krasnosel'ski\u{\i}--Mann iterations achieve $o(1/k)$ on the squared residual
and, in the worst case, nothing faster
\citep{DavisYin2016Convergence,BravoCominetti2018Sharp}.

\begin{remark}[What the certificate buys the search]
\label{rem:relaxation-window}
Because $1/\theta_G=2-c_G$ for every base scheme, the certified
over-relaxation window is exactly $\rho\in(0,2-c_G)$: the textbook interval
$(0,2)$ for the schemes with no forward step, shrinking as the forward step
consumes the budget. For Condat--V\~u this window is not a conservative
surrogate. \citet[Theorem~3.1]{Condat2013Primal} proves convergence of the
same relaxed iteration for
$\rho_n\in\bigl(0,\delta\bigr)$ with
$\delta=2-\frac{\beta}{2}\bigl(\frac1\tau-\sigma\|L\|^2\bigr)^{-1}$, where
$\beta=L_f$ is the Lipschitz constant of the smooth gradient. Multiplying
numerator and denominator by $\tau$ gives
$\delta=2-c_\star$, where
$c_\star:=\frac{\tau L_f}{2(1-\tau\sigma\|L\|^2)}$ is the forward-step
constant computed with the exact operator norm. Since $\ell\geq\|L\|$, it follows
that $c_\star\leq c_G$ and hence $2-c_G\leq\delta$. The certified window
therefore coincides with Condat's when $\ell=\|L\|$, and is contained in it
when $\ell>\|L\|$. The discovered policy of Section~\ref{sec:results}
illustrates this effect. The certified textbook default
$(\tau,\sigma)=(0.5,0.186)$ has $c_G\approx0.97$ and therefore admits almost
no over-relaxation at all, $\rho<1.03$, whereas the evolved
$(\tau,\sigma)=(0.05,1.0)$ has $c_G=0.042$ and admits $\rho<1.96$. Shrinking
the primal step does not merely satisfy the gate, it widens the interval the
relaxation gene may search in, and the champion sits near its new edge at
$\rho=1.9$, with $\rho\,\theta=0.97$.
\end{remark}

\paragraph{Halpern genomes.}
An averaged map $T_G$ is nonexpansive. The implementation applies the
canonical Halpern combination directly to this map, including a permitted
relaxation, as recorded in Appendix~\ref{app:instantiation}. Without anchor
restarts, Proposition~\ref{prop:halpern} applies with that map in the role of
its nonexpansive operator. The normalized map
$S_G=I+\theta_G^{-1}(T_G-I)$ in the proposition is another theoretically
admissible choice; it is not the map selected by the current runtime.
The unrestarted guarantee does not establish the same rate or anchor
projection for periodically restarted trajectories. Neither Halpern nor
restarted Halpern is used in the current hardware factorial.

\paragraph{Finite-prefix switching.}
% PROPOSED to Nikhil: replaces summable-error safeguard with finite-prefix
% switch, matches implementation
Corollary~\ref{cor:convergence} assumes exact arithmetic and correct atom
implementations. Each returned numerical solution is also checked by
the numerical stopping criterion of Section~\ref{sec:methods}. Free-form
schedules, packaged as the $\mathrm{switch}_K$ gene of
Section~\ref{sec:mixing}, run outside $\mathcal G$ as a finite prefix: an
arbitrary map $T_A$ drives the first $K$ iterations, and the compiled
iteration then switches permanently to a certified map $T_B$, after a
finiteness check on the handoff state $z_K$. Because the tail is exactly
the certified iteration started at $z_K$, the sequence converges to a fixed
point of $T_B$ and the rate \eqref{eq:residual-rate} holds with $z_0$
replaced by $z_K$ (Lemma~\ref{lem:switchk}, Appendix~\ref{app:theory}); the
prefix enters only through the constant
$\operatorname{dist}_{H_G}(z_K,\operatorname{Fix}T_B)$, never the
asymptotics. This extension is not used in the experiments. All reported
runs have $K=0$.

Every admitted genome in $\mathcal G$ therefore has the stated
exact-arithmetic guarantee. Finite-budget numerical success remains a separate
measurement.

\FloatBarrier

% Section 5. Experiments.
% Every number carries a "% source:" comment naming a file under
% /home/miria/OptEvolve/v2/results/. Nothing here comes from atlas/.
% Style contract: no em-dash or en-dash characters, no triple-hyphen dash
% macro, no italic markup of any kind, pure ASCII, plain declarative prose.
% All figures below were recomputed from the raw JSON rows on 2026-09-10 by
% opening the cited file and agree with it to the digits printed. Where a
% recomputation disagreed with an earlier draft the earlier value was deleted,
% not softened.

\section{Experiments}
\label{sec:results}
\label{sec:exp}

All measurements run on an RTX 4090 (sm\_89) and an RTX PRO 6000 Blackwell
Max-Q (sm\_120), over four families: total variation denoising of six retinal
OCT scans, min-cost flow on grid and regular topologies, capacitated
multi-commodity flow, and optimal transport. Termination is decided by a
float64 relative duality gap in original coordinates at $10^{-4}$ unless stated
otherwise. Every time is wall clock to a certified gap and every iteration
count is the count at which that certificate passed. Iteration count is written
$N$ and per-iteration cost $t$, and we state which of the two any ratio is
measured on, because the thesis of the paper is $T = N \times t$.

\subsection{Problem families}
\label{sec:families}

% source: row and cell counts enumerated from every JSON under results/ carrying a
% per-instance 'cell' or 'family' key, plus the tolerance-keyed factorials for TV.
Table~\ref{tab:families} lists the five families on which every number in this
paper was measured, with the volume of measurement behind each. The three
linear-programming families share the same standard-form composite and are
solved by the same code path, which is what makes the transfer experiment of
Section~\ref{sec:exp:transfer} meaningful: the solver is held fixed and only the
operator changes. Total variation is the one non-LP family and the only one to
which the fused kernel of Section~\ref{sec:exp:fusion} applies, since that kernel is
a stencil.

\begin{table}[t]
\caption{Problem families with measurements. Rows count individual solves
recorded under \texttt{results/}. Cells count distinct problem configurations.}
\label{tab:families}
\begin{center}
\begin{tabular}{llrr}
\toprule
Family & Configurations & Rows & Cells \\
\midrule
Min-cost flow, grid & $k \times k$ grid, bidirectional 4-neighbour arcs & 5{,}645 & 4 \\
Min-cost flow, random regular & degrees 2, 3, 4; $n = 256, 512, 1024$ & 3{,}929 & 8 \\
Optimal transport & bipartite, $m = n \in \{32, 64, 96\}$ & 1{,}431 & 3 \\
Total-variation denoising & OCT scans, 128 to 1984 pixels & 1{,}003 & 7 sizes \\
Multi-commodity flow & capacity-coupled, $K \in \{2, 4\}$ & 648 & 2 \\
\bottomrule
\end{tabular}
\end{center}
\end{table}

The families are not variations on a single operator. Table~\ref{tab:opnorm}
gives the behaviour of $\|L\|$, which is the quantity the construction gate of
Section~\ref{sec:exp:method} tests and therefore the quantity that decides which
recipes are admissible at all. Two families have an operator norm bounded
independently of problem size, one grows without bound in the size, and one
contracts in its coupling parameter. This spread is the reason the admissible
region is family dependent rather than a property of the solver.

\begin{table}[t]
\caption{Operator-norm behaviour across the five families, measured. The
reported $\|L\|$ is the gate's bound, which carries a $1.02$ safety factor over
the spectral norm.}
\label{tab:opnorm}
\begin{center}
\begin{tabular}{lll}
\toprule
Family & $\|L\|$ behaviour & Measured \\
\midrule
% source: build_instances('grid', k) norm_bound; 528 to 9024 variables
Min-cost flow, grid & saturates at $4$, bounded in size & $4.045$ to $4.076$ over a $17\times$ size increase \\
% source: build_instances('regular', n, degree=d) norm_bound
Min-cost flow, regular & bounded by degree, $\leq \sqrt{4d}$ & $2.884$ at $d=2$, $3.935$ at $d=4$ \\
% source: atlas/bench/ot.py transport_composite norm_bound vs 1.02*sqrt(m+n)
Optimal transport & $\sqrt{m+n}$, unbounded in size & exact to $3 \times 10^{-14}$ relative \\
% source: atlas/bench/mcf.py mcf_composite norm_bound, grid size 8
Multi-commodity flow & grows with $K$, admissible $\tau\sigma$ contracts & $\|L\|^2$ from $17.1$ at $K{=}1$ to $150.1$ at $K{=}128$ \\
Total-variation denoising & gradient stencil, bounded & fixed by the discretisation \\
\bottomrule
\end{tabular}
\end{center}
\end{table}

The consequence is direct. On grid instances the admissible region is
effectively fixed: $\|L\|$ moves $0.76$ percent while the variable count grows
$17\times$, so the same step sizes remain legal at every size. On optimal
transport the admissible product $\tau\sigma$ must shrink as $1/(m+n)$, and both
catalog step-size settings become inadmissible at sizes far below those we
solve. On multi-commodity flow the admissible product falls from $0.0584$ at
$K=1$ to $0.00666$ at $K=128$. A step-size pair is therefore not portable across
families, and the gate detects this symbolically, before any iteration runs.

\subsection{The factorisation}
\label{sec:exp:factorisation}

% source: results/tv_ffi_20260906/cadence_4090_v1/factorial.json
% source: results/bw_20260909/sz256_bw/factorial.json
Table~\ref{tab:Nvectors} gives the per-instance iteration counts on TV at
$n=256$. The six-element vectors, not only their means, are identical element by
element on the two architectures at every cadence and under both schemes.

\begin{table}[t]
\centering
\caption{Per-instance iteration counts over the six OCT scans at $n=256$,
tolerance $10^{-4}$. Every vector is identical on the RTX 4090 and on the
Blackwell. The vectors are also identical across the float64, float32-iterate
and two-phase precision policies at $K=64$ and $K=256$. At $K=16$ they are not:
the float32-iterate policy moves PDHG from the float64 vector
$[1008, 496, 608, 400, 320, 640]$, mean $578.7$, to
$[1008, 512, 608, 416, 320, 640]$, mean $584.0$, while Condat-Vu is unchanged
at $584.0$.}
\label{tab:Nvectors}
\begin{tabular}{@{}llrr@{}}
\toprule
Cadence & Scheme & Per-instance $N$ & Mean \\
\midrule
$K=16$  & Condat-Vu & 1008, 512, 608, 416, 320, 640 & 584.0 \\
$K=16$  & PDHG      & 1008, 496, 608, 400, 320, 640 & 578.7 \\
$K=64$  & both      & 1024, 512, 640, 448, 320, 640 & 597.3 \\
$K=256$ & both      & 1024, 512, 768, 512, 512, 768 & 682.7 \\
\bottomrule
\end{tabular}
\end{table}

% source: results/tv_ffi_20260906/{cadence_4090_v1,f32row_4090,f64host_4090,kernelhost_4090}/factorial.json
% source: results/bw_20260909/{sz128_bw,sz496_bw}/factorial.json
% The width counterexamples are the reason this paragraph is scoped. They were
% recomputed from the two bw_20260909 files and from f32row_4090 named above.
The same vectors appear in four separate factorial artifacts covering the
float64 device loop, the float32-iterate float64-certificate device loop, the
host-phase policy and the hand written kernel in both execution structures. At
$n=256$ and $K \geq 64$ iteration count is invariant to arithmetic width,
execution structure, kernel and device, and is a function
of cadence and tolerance alone. Width is not inert in general. On the Blackwell
at $K=64$ float64 gives $576.0$ against float32's $586.7$ at $n=128$ and
$1066.7$ against $1077.3$ at $n=496$, and on the 4090 at $K=16$ float32 moves
PDHG from $578.7$ to $584.0$. Cadence quantises the width difference away at
$K \geq 64$ and $n=256$ and not everywhere.

% source: results/kernel_fusion_20260909/e2e/full_4090_sz{128,256,384,496}.json
% source: results/kernel_fusion_20260909/e2e_bw/full_bw_sz{128,256,384,496}.json
% Within-cell iteration identity and the 6-of-6 certification were rechecked
% arm by arm; the largest within-cell time spread over the eight cells is 2.6294.
The invariance survives into the deployed solver. Across eight arms, the traced
baseline, the static loop and six fusion tiles, at four sizes on both devices,
the per-instance vectors are identical within each cell and every arm certifies
6 of 6, while solve time within a cell varies by up to $2.629\times$. The sweep
requested seven tiles, $8\times4$, $16\times2$, $16\times4$, $24\times4$,
$32\times2$, $32\times3$ and $32\times4$. The $32\times3$ tile never runs,
because the harness requires the fused depth $T$ to divide the cadence $K$ and
no recorded cadence, $64$, $128$, $256$ or $512$, is divisible by 3. Six tiles
are measured everywhere in this paper.

% source: results/netflow_20260909/grid_4090.json
% 48 rows: 24 banked (all executed) and 24 balanced (all gate-rejected).
% The two rows at grid size 12, instance 1, hit the 20000 cap at gap 2.858e-4.
The second family agrees. All 24 banked min-cost flow rows return exactly equal
iteration counts under the static and the traced loop, and 22 of the 24
certify; the two rows at grid size 12, instance 1, agree only because both hit
the 20000-iteration cap at a gap of $2.858 \times 10^{-4}$. Over the 22
certified rows the traced-over-static wall-time ratio has geometric mean
$1.395$ and spans $1.293$ to $1.492$; over all 24, including the two capped
rows, it is $1.385$ and spans $1.277$ to $1.492$. All 12 matched instance pairs
return exactly equal counts under PDHG and Condat-Vu, of which 11 are certified
pairs and one is the pair that both schemes ran to the cap.

% source: results/method_axis_20260909/factorial.json (XLA half)
% source: results/tv_ffi_20260906/kernelhost_4090/factorial.json (kernel half)
% Both halves recomputed as solver_s / iters, averaged over the six images.
Per-iteration cost is invariant to scheme identity under XLA and is not under
the hand written kernel. At \texttt{f64\_K256} under XLA the two schemes return
the identical vector $[1024, 512, 768, 512, 512, 768]$, penalised SGM-1 times in
ratio $0.99494$ and per-iteration costs of $19.76$ against $19.76$
microseconds. In the custom-kernel artifact, with the same
\texttt{build\_digest} in both rows, Condat-Vu costs $19.21$ microseconds per
iteration and PDHG $25.02$, a factor of $1.303$, and the gap appears in all
four custom-kernel cells at $1.303$, $1.264$, $1.179$ and $1.197$. Scheme
identity is inert to $t$ on the compiler path and not on the kernel path.

% source: results/tv_ffi_20260906/cadence_4090_v1/factorial.json and
% results/bw_20260909/sz256_bw/factorial.json and
% results/method_axis_20260909/factorial.json
% CORRECTION relative to results/tv_ffi_20260906/DECOMPOSITION_REPORT.md, which
% states 584/597/683 for both schemes. That is true at K=64 and K=256 only.
Scheme identity is exactly inert to $N$ at coarse cadence and not at fine
cadence. At $K=16$ the schemes differ on two of the six instances, 512 against
496 and 416 against 400, and the split reproduces on the Blackwell. At
$10^{-6}$ and $K=64$ Condat-Vu averages $2860.0$ iterations and PDHG $2870.7$;
both certify 5 of 6, both means include the same capped instance, and the two
schemes differ on exactly one of the five certified instances, 2432 against
2496. A coarse cadence quantises a real scheme difference away, which is the
mechanism this paper argues for rather than an exception to it.

\subsection{Execution structure}
\label{sec:exp:execution}

% source: results/artifacts_20260910/fanova.json, block
% exec_x_prec_x_cadence__logT__tol0.0001, which carries all twelve cells to
% full precision and names its four source factorials:
%   device_loop|f64  -> tv_ffi_20260906/cadence_4090_v1/factorial.json
%   device_loop|f32  -> tv_ffi_20260906/f32row_4090/factorial.json
%   host_phases|f64  -> tv_ffi_20260906/f64host_4090/factorial.json
%   host_phases|f32  -> tv_ffi_20260906/px999999_4090/factorial.json
% The twelve values are 22.536573, 20.389612, 19.777049, 22.986976, 20.405764,
% 19.775866, 27.184947, 16.185603, 13.290695, 27.998956, 12.228629, 8.431614.
Table~\ref{tab:periter} decomposes per-iteration cost over execution structure,
arithmetic width and cadence.

\begin{table}[t]
\centering
\caption{Microseconds per iteration on the RTX 4090, TV at $n=256$, Condat-Vu,
averaged over the six scans. PDHG in the same artifacts gives $27.40$, $17.20$
and $14.15$ at host-phase float64 and $26.59$, $12.21$ and $8.36$ at host-phase
float32. The host-phase float32 row is the patched policy registry and not the
shipped genome, as disclosed at the end of this subsection.}
\label{tab:periter}
\begin{tabular}{@{}llrrr@{}}
\toprule
Execution & Width & $K=16$ & $K=64$ & $K=256$ \\
\midrule
Device loop  & float64                & 22.54 & 20.39 & 19.78 \\
Device loop  & float32 iterates       & 22.99 & 20.41 & 19.78 \\
Host phases  & float64                & 27.18 & 16.19 & 13.29 \\
Host phases  & float32 throughout     & 28.00 & 12.23 &  8.43 \\
\bottomrule
\end{tabular}
\end{table}

Width buys $1.00\times$ inside the device loop at $K=256$, where the measured
ratio is $1.00006$ and the section's own stability figure is $0.66$ percent, so
the third decimal is below resolution. It buys $1.576\times$ once the loop is
chunked into host phases, with the float32 host-phase cell taken from the
patched policy registry rather than the shipped genome. Execution structure is
worth $1.488\times$ at float64 and $2.345\times$ at float32 at $K=256$, and
reverses to $0.829\times$ at $K=16$, where chunking pays more in launches than
it saves. Under PDHG the float64 execution-structure gain at $K=256$ is
$1.398\times$ rather than $1.488\times$. Reporting one of these without the
other misattributes the gain.

% source: results/artifacts_20260910/fanova.json, same block, percent_variance.
% Response is log per-iteration cost; factors are execution, precision, cadence.
% This is NOT the algorithm-by-precision-by-cadence decomposition on log
% penalised SGM-1, whose algorithm share is 0.00 percent at 1e-4 and 0.34 at
% 1e-6 and which is reported in the introduction and in Section 2
% (2-related-work.tex lines 96 to 98). Do not merge the two decompositions.
A three-way decomposition of log per-iteration cost over those twelve cells
assigns $47.43$ percent of the variance to cadence, $25.96$ percent to the
execution-by-cadence interaction, $15.98$ percent to execution structure,
$3.26$ percent to execution by precision, $2.90$ percent to precision alone,
$2.43$ percent to precision by cadence and $2.04$ percent to the three-way
residual. The design is saturated at one observation per cell, so these are
shares of an exact decomposition and carry no error bar.

% source: results/finish_20260910/loop_mechanism.json, 48 rows, recomputed.
% The 48 rows are 2 sizes x 6 trip counts x 2 traced x 2 certificate settings,
% so the certificate term is a difference and there are 24 of them, not 48.
Table~\ref{tab:loopmech} settles the cause. It holds the total iteration budget
fixed at 512 and moves the number of outer trips from 32 down to 1, with the
cadence moving inversely.

\begin{table}[t]
\centering
\caption{Traced minus static cost, nanoseconds per iteration, certificate off.
The product of outer trips and cadence is 512 in every column. One measurement
per cell, no repeats.}
\label{tab:loopmech}
\begin{tabular}{@{}lrrrrrr@{}}
\toprule
Outer trips & 32 & 16 & 8 & 4 & 2 & 1 \\
\midrule
$n=128$ penalty & 11972 & 11365 & 10893 & 11014 & 10823 & 10769 \\
$n=128$ ratio   & 2.374 & 2.348 & 2.299 & 2.361 & 2.333 & 2.317 \\
$n=256$ penalty & 12943 & 12564 & 11604 & 11929 & 11202 & 11531 \\
$n=256$ ratio   & 2.091 & 2.126 & 1.989 & 2.060 & 1.938 & 1.997 \\
\bottomrule
\end{tabular}
\end{table}

The penalty is per iteration and nearly flat in the number of outer trips. It
falls $10.0$ percent at $n=128$, from $11972$ to $10769$\,ns, and $10.9$ percent
at $n=256$, from $12943$ to $11531$\,ns, while the outer count falls 32-fold,
and the static per-iteration cost itself moves only $6.2$ percent at $n=128$,
from $8714$ to $8175$\,ns, across the same 32-fold change in unroll depth. A
lost cross-iteration fusion opportunity would not be 10 percent of a 32-fold
change. The file carries one measurement per cell with no repeats, so any
residual trend is at the resolution of a single shot. Over the 24 matched
differences the certificate term, cost with the certificate on minus cost with
it off, spans $-239.90$ to $+844.56$\,ns against a loop penalty near
$11000$\,ns, so it is not a per-certificate-check cost either, which is the
explanation we previously advanced. The largest certificate term, $+844.56$\,ns,
falls at $n=128$ and 32 outer trips, the finest cadence in the file, where a
real certificate cost would appear, and is 7 percent of the loop penalty there.

% source: results/kernel_fusion_20260909/e2e/full_4090_sz{128,256,384,496}.json
% and e2e_bw/full_bw_sz{128,256,384,496}.json
End to end the static bound alone is worth $1.673\times$, $1.347\times$,
$1.291\times$ and $1.093\times$ at $n=128$, $256$, $384$ and $496$ on the 4090
and $1.228\times$, $1.382\times$, $1.401\times$ and $1.243\times$ on the
Blackwell. The 4090 benefit decays with size and the Blackwell benefit does
not.

% source: results/tv_ffi_20260906/kernelhost_4090/factorial.json
% 19.2068 / 19.7628 / 13.2732 / 12.9007 under Condat-Vu;
% 25.0212 / 19.7619 / 15.6475 / 13.7109 under PDHG.
Under Condat-Vu the custom TV kernel is at parity with XLA on this device in
both execution structures: $19.21$ against $19.76$ microseconds per iteration in
the device loop at $K=256$, an edge of $1.029\times$, and $13.27$ against the
schedule policy's $12.90$ in host phases, a loss of $0.972\times$. Under PDHG
the same kernel loses in both, $25.02$ against $19.76$, a loss of
$1.266\times$, and $15.65$ against $13.71$, a loss of $1.141\times$. The kernel
gains $1.447\times$ from chunked execution while XLA in the same artifact gains
$1.532\times$, so both gain from chunking, the gains are not equal, and
execution structure and kernel identity remain separate decisions.

% source: results/tv_ffi_20260906/px999999_4090/factorial.json (0.00576135) and
% cadence_4090_v1/factorial.json (0.01215536 baseline, 0.00789548 catalog best).
% 0.01215536 / 0.00576135 = 2.10981; 0.01215536 / 0.00789548 = 1.53953.
% These two cells were produced by patching the policy registry, not by the
% shipped genome. Label them as diagnostics, never as search output.
The best configuration measured is not expressible in the shipped genome: XLA
float32 with host phases at $K=256$ reaches $0.00576$\,s penalised SGM-1 against
the compiled float64 and $K=64$ baseline's $0.01216$\,s, a speedup of
$2.110\times$, where the best catalog cell reaches $0.00790$\,s and
$1.540\times$. That cell was produced by patching the policy registry rather
than by the shipped genome, so it is a diagnostic that motivates a schema
change and not a search result.

\subsection{Temporal kernel fusion}
\label{sec:exp:fusion}

% source: results/artifacts_20260910/kernel_admission_4090.log, 106 lines,
% header results/artifacts_20260910/kernel_admission_4090.header.
% 17 variants x 3 sizes x 2 proximal settings = 102 comparison rows, rho = 1.9
% throughout, bitwise column "no" in every row, final line PASS.
The fused kernel is numerically admitted before it is timed. The archived
admission run compares $T$ fused iterations against $T$ sequential single-step
reference launches over 102 cases, 17 built variants at three problem shapes
and both proximal settings, all at $\rho = 1.9$. The worst absolute deviation
over all 102 cases is $4.649 \times 10^{-15}$ and the run reports PASS. No case
is bitwise identical at this $\rho$. The variant names in the log carry the
thread count, so the two builds that share the label $16\times4$, at 128 and at
256 threads, are admitted separately, and all six timed tiles appear in the
log.

% source: results/kernel_fusion_20260909/e2e/full_4090_sz{128,256,384,496}.json
% source: results/kernel_fusion_20260909/e2e_bw/full_bw_sz{128,256,384,496}.json
% Totals recomputed by summing t_s over the six instances per arm.
% NOTE FOR COAUTHOR: the best-speedup and baseline rows here repeat
% Table~\ref{tab:intro-e2e} in the introduction. Keep one of the two, or drop
% the baseline rows here and keep only the worst-tile rows.
Table~\ref{tab:fusion} gives the end-to-end fusion result against the shipped
traced baseline, with the worst tile of the same sweep beside it.

\begin{table}[t]
\centering
\caption{Time to a certified $10^{-4}$ gap over six OCT scans, $K=256$,
relative to the shipped traced-bound baseline, three timing repeats per image.
Every arm certified 6 of 6 and every arm inside a cell ran the identical
per-instance iteration counts. The worst-tile identity at $n=256$ is not
stable: a 9-repeat re-measurement one day later puts $8\times4$ worst at
$1.366$ with $32\times2$ at $1.375$, so that one cell is inside the measurement
spread and the tile name is omitted.}
\label{tab:fusion}
\begin{tabular}{@{}lrrrr@{}}
\toprule
$n$ & 128 & 256 & 384 & 496 \\
\midrule
4090 baseline (s)  & 0.06647 & 0.07733 & 0.13486 & 0.22564 \\
4090 best tile     & $8\times4$ & $24\times4$ & $24\times4$ & $32\times4$ \\
4090 best speedup  & 2.629 & 1.867 & 1.634 & 1.356 \\
4090 worst tile    & $32\times2$ & not resolved & $8\times4$ & $8\times4$ \\
4090 worst speedup & 1.320 & 1.298 & 0.926 & 0.743 \\
\midrule
Blackwell baseline (s) & 0.05732 & 0.07299 & 0.13750 & 0.19964 \\
Blackwell best tile    & $8\times4$ & $24\times4$ & $32\times4$ & $16\times2$ \\
Blackwell best speedup & 1.732 & 1.450 & 1.625 & 1.149 \\
Blackwell worst speedup & 0.973 & 1.008 & 1.100 & 0.862 \\
\bottomrule
\end{tabular}
\end{table}

% source: results/kernel_fusion_20260909/e2e/full_4090_sz256.json against
% results/finish_20260910/fusion_r9_sz256.json (9 repeats, one day later).
The optimal tile moves with size and with device: the two devices agree at
$n=128$ and $n=256$ and disagree at $n=384$ and $n=496$. Tile selection is not
optional. The wrong tile is slower than not fusing at all at the two larger
sizes on the 4090, and the best-to-worst swing is $1.99\times$ at $n=128$ and
$1.82\times$ at $n=496$. The identity of the worst tile is $8\times4$ at both
large sizes and is not resolved at $n=256$, where three repeats and nine
repeats disagree about which of $8\times4$ and $32\times2$ is worst and the two
differ by 0.7 percent under the nine-repeat run.

% source: results/kernel_fusion_20260909/e2e_bw/full_bw_sz496.json
There is a cell where the correct decision is not to fuse. On the Blackwell at
$n=496$ the best tile reaches only $0.925\times$ of the plain static-loop path
and the worst reaches $0.862\times$ of the baseline, which belongs in the same
paragraph as the $1.149\times$ headline for that cell.

% source: results/kernel_fusion_20260909/e2e/full_4090_sz*.json (128 to 496),
% large/large_4090_m{2,3,4}.json and bw_day2/large_bw_m{2,3,4}.json (992 to 1984)
% CAVEAT: the three montage sizes use four images per arm, not six, and are
% built by tiling several distinct scans, because a raw scan caps an honest
% crop near 496 px. The pixel sizes 992/1488/1984 are the m2/m3/m4 montage
% convention; the args block of those three files records size 496 and n 4.
Table~\ref{tab:ushape} sweeps size to 1984 pixels. The fusion gain over the
static loop is U-shaped on the 4090: it falls monotonically from $1.571\times$
at 128 pixels to a minimum of $1.060\times$ at 992 pixels and rises to
$1.410\times$ at 1984 pixels. The Blackwell row is not monotone on either limb,
running $1.410$, $1.049$, $1.160$, $0.925$, $0.989$, $1.033$ and $1.123$, and
its minimum sits at 496 pixels. The 4090 trough is a single point and it sits
exactly at the boundary between the six-crop protocol and the four-montage
protocol, so we read the shape and not the location. Launch-overhead removal
accounts for the small end and reduced memory traffic for the large end.

\begin{table}[t]
\centering
\caption{Fusion gain over the static loop and static-loop gain over the traced
bound, by problem size. Sizes 992 and above are montages of several distinct
scans, four images per arm rather than six, so the two limbs of the 4090 row
use different instance construction.}
\label{tab:ushape}
\begin{tabular}{@{}lrrrrrrr@{}}
\toprule
Size (px) & 128 & 256 & 384 & 496 & 992 & 1488 & 1984 \\
\midrule
4090 fusion over static      & 1.571 & 1.386 & 1.266 & 1.240 & 1.060 & 1.346 & 1.410 \\
Blackwell fusion over static & 1.410 & 1.049 & 1.160 & 0.925 & 0.989 & 1.033 & 1.123 \\
4090 static over traced      & 1.673 & 1.347 & 1.291 & 1.093 & 1.016 & 0.989 & 0.993 \\
Blackwell static over traced & 1.228 & 1.382 & 1.401 & 1.243 & 1.063 & 1.060 & 1.037 \\
\bottomrule
\end{tabular}
\end{table}

The static bound decays monotonically on the 4090 and reverses past about 1000
pixels, to $0.989\times$ at 1488 and $0.993\times$ at 1984, while on the
Blackwell it does not reverse in this range. At the two largest 4090 sizes the
best recipe is the traced bound plus fusion rather than the static bound plus
fusion, so the two execution decisions are not separable at every scale.

% source: results/kernel_fusion_20260909/cadence/cad_4090_K{64,128,256,512}.json
% source: results/kernel_fusion_20260909/bw_day2/cad_bw_K{64,128,256,512}.json
% CORRECTION relative to results/kernel_fusion_20260909/REPORT.md section 6,
% which states the iteration count is 512 at every K in this sweep. It is not.
Across cadence the argmax is stable and its value is not. At $n=256$ on the
4090 the fusion gain over the static loop is $1.258\times$, $1.265\times$,
$1.336\times$ and $1.394\times$ at $K=64$, $128$, $256$ and $512$, and the best
tile is $24\times4$ at three of four cadences, losing at $K=128$ to $16\times4$
by 0.6 percent. The Blackwell selects $24\times4$ at all four with gains
$1.006\times$, $1.024\times$, $1.053\times$ and $1.062\times$. Iteration count
is not constant across this sweep, at $3712$, $3840$, $3840$ and $4608$ total
over the six images, so only the gain-over-static-loop column is a clean
per-iteration comparison and we report only that column.

\subsection{The method axis and the construction gate}
\label{sec:exp:method}

% source: results/method_axis_20260909/factorial.json
% 2944 / 682.6667 = 4.3125 exactly, at f64_K256 and tolerance 1e-4.
% SGM-1: condat_vu 0.05592390 / 0.01349855 = 4.142957;
%        pdhg      0.05635931 / 0.01343025 = 4.196444.
Scheme identity is inert on the method axis and step sizes are decisive. At
$n=256$ under \texttt{f64\_K256} and tolerance $10^{-4}$ the banked setting
$(\tau, \sigma, \rho) = (0.05, 1.0, 1.9)$ gives mean $682.7$ iterations under
both schemes and the balanced setting $(0.25, 0.25, 1.5)$ gives $2944.0$ under
both, a ratio of $4.3125$ on iteration count under both schemes, and on
penalised SGM-1 time $4.1430$ under Condat-Vu and $4.1964$ under PDHG. At
$10^{-6}$ the difference becomes feasibility: banked certifies 5 of 6 at mean
$2966.7$ iterations and balanced certifies 0 of 6, every instance running to the
5000-iteration cap, though balanced does certify 6 of 6 at $10^{-4}$.

% source: results/netflow_20260909/grid_4090.json (both error strings read
% verbatim: the Condat-Vu rows print "t*L_f/2 + t*s*||L||^2 < 1 (lhs=...)" and
% the PDHG rows print "t*s*||L||^2 = ..."; the recorded strings are 1.02267,
% 1.0304, 1.03577, 1.03781).
% source: results/structure_20260909/structure_summary_4090.json
% (raw_t025_s025_gate_lhs = 1.0226746, 1.0304043, 1.0357716, 1.0378125).
The admissible region is family dependent, and the difference is one of kind.
The same balanced parameters that are merely slow on TV are rejected before any
arithmetic runs on every min-cost flow instance: all 24 balanced netflow rows
fail with gate left-hand side $1.02267$, $1.0304$, $1.03577$ and $1.03781$ at
grid sizes 12, 16, 24 and 32 against a threshold of 1, and all 24 banked rows
run. The Condat-Vu rows record $t L_f / 2 + t s \lVert L \rVert^2$ and the PDHG
rows record $t s \lVert L \rVert^2$; the two coincide here because the min-cost
flow objective has no smooth term. The structure probe reproduces those four
values at $t = s = 0.25$ as $1.0226746$, $1.0304043$, $1.0357716$ and
$1.0378125$, against $0.51994$ to $0.52005$ across the three degree-2 regular
cells, which are admitted.

% source: results/structure_20260909/structure_summary_4090.json
% totals: runs 2106, admitted 1755, rejected 351, capped 257,
% gpu_exec_seconds 987.81.
Admission is a function of the gate product alone: over 2106 runs the gate
admits 351 of 351 settings at each of $\theta \in \{0.1, 0.3, 0.6, 0.9,
0.99\}$ and 0 of 351 at $\theta = 1.02$, admitting 1755 runs in total,
rejecting 351, capping 257 at the iteration limit and consuming 987.8 GPU
seconds. Rejected settings were never executed, so the gate is verified as
sufficient for non-divergence and never as necessary, and 257 of the admitted
runs still failed to certify inside the cap.

% source: results/structure_20260909/structure_summary_4090.json (norm bounds)
% source: results/finish_20260910/omega_ot_4090.json
% (Lnorm = 8.15999999999979, 11.5399826689641, 14.1335345897619;
%  1.02*sqrt(64) = 8.16, relative errors 2.6e-14, 3.2e-14, 1.1e-14)
How the region scales with instance size is also a property of the family.
Min-cost flow norm bounds are bounded in size, at $4.0451$ to $4.0749$ on the
grid over sizes 12 to 32 and, on the regular topology across 512, 1024 and 2048
nodes, $2.8843$ to $2.8844$ at degree 2, $3.4781$ to $3.4803$ at degree 3 and
$3.9364$ to $3.9395$ at degree 4. Optimal transport is the opposite case, and
measurement confirms its closed form: the recorded norm bound is $8.160000$,
$11.539983$ and $14.133535$ at the $32\times32$, $64\times64$ and $96\times96$
cells, which is $1.02\sqrt{m+n}$ to a relative error of $1.1 \times 10^{-14}$
to $3.2 \times 10^{-14}$, fourteen significant figures rather than bitwise
agreement, the $1.02$ being the safety inflation and the residual being
power-iteration accuracy. The admissible product is $0.0150183$, $0.0075091$
and $0.0050061$, in ratios $2.000$ and $1.500$ that match $m+n = 64$, $128$ and
$192$ exactly, so that region shrinks as $1/(m{+}n)$.

% source: results/finish_20260910/omega_mcf.json
% Lnorm 4.323183 at 2 commodities and 4.567979 at 4, both on the 16x16 grid,
% identical across the three seeds; 4.567979^2 / 4.323183^2 = 1.1165.
The corpus now contains capacitated multi-commodity flow runs, at two commodity
counts on one topology. Doubling the commodity count from 2 to 4 on the
$16\times16$ grid moves the norm bound from $4.323183$ to $4.567979$ and
shrinks the admissible product by a factor of $1.117$, not by a factor of 2. Two
points do not identify a rate, so we make no asymptotic claim about how that
family's admissible region contracts in the number of commodities.

\subsection{Predicting the step ratio from measured structure}
\label{sec:exp:omega}

% source: results/finish_20260910/omega_4090.json, omega_bw.json (native
%   omega0 over the 24 cell-and-seed groups: 6.9691 to 9.1254);
% source: results/structure_20260909/structure_summary_4090.json
%   (omega0 over the 13 probe cells: 7.5058 to 7.6733);
% source: results/finish_20260910/omega_ot_4090.json
%   (native omega0 on optimal transport: 0.4464 to 0.7928).
% sweep source scripts/diagnostics/omega_sweep.py
The step-size recipe looked unpredictable for a reason visible only once the
data are pooled. At native scaling the scale ratio $\omega_0 = \|c\| / \|b\|$
barely varies within a family: it spans $6.9691$ to $9.1254$ over the 24
cell-and-seed groups of the netflow sweep, $7.5058$ to $7.6733$ over the 13
cells of the structure probe and $0.4464$ to $0.7928$ on optimal transport. A
variable that never varied was never implicated.

% source: results/finish_20260910/omega_4090.json and omega_bw.json.
% Construction: tau = sqrt(theta*r)/||L||, sigma = sqrt(theta/r)/||L||, so the
% gate product is theta by construction and the ratio is r by construction.
% The sweep solves with PDHG at rho = 1 throughout.
The sweep moves it deliberately. Rescaling $c \mapsto \alpha c$ is exactly
$(\tau, \sigma) \mapsto (\alpha \tau, \sigma / \alpha)$, which leaves the gate
product fixed and moves $r = \tau / \sigma$ by $\alpha^2$. We sweep $\alpha$
and $r$ each over $\{0.01, 0.03, 0.1, 0.3, 1, 3, 10, 30, 100\}$, four decades,
on 8 cells and 3 seeds at $\theta = 0.99$, $K = 64$, tolerance $10^{-4}$ and a
cap of 60000 iterations, solving with PDHG at $\rho = 1$: 1944 runs, 1510
converged, over 216 instances of which 209 have at least one converged run. The
gate product $\tau \sigma \lVert L \rVert^2$ equals $0.99$ to within
$3.4 \times 10^{-16}$ in all 1944 rows and exactly in 774 of them, so the sweep
moves the ratio and never the admissible region. It is worth being plain about
what that check is and is not: $\theta = 0.99$ is a fixed input to the sweep and
$\tau$ and $\sigma$ are constructed from it, so the check verifies the
construction arithmetic and is not evidence that the gate depends on $\theta$
alone.

% source: results/finish_20260910/omega_4090.json and omega_bw.json, recomputed.
% PROTOCOL: converged rows only; a group is a (cell, seed) pair sharing a value
% of log10(r) + 2 log10(alpha) and containing at least two distinct alphas and
% at least two converged runs; a group qualifies for a window only if EVERY
% alpha in it lies inside the window. Counting singleton groups as trivially
% agreeing inflates these counts.
The exponent 2 in that rescaling identity, which is a different object from the
exponent $p$ of the rule fitted below, is a measured invariance and not a
fitted number. Group runs by
$\log_{10} r + 2 \log_{10} \alpha$ within each cell and seed, and count a group
only if it holds at least two distinct values of $\alpha$ and at least two
converged runs. A group qualifies for a window only if every $\alpha$ in it
lies inside that window. Every group whose values of $\alpha$ are all at least
1 returns identical iteration counts, 169 of 169 at $\alpha \geq 1$, 83 of 83 at
$\alpha \geq 3$ and 25 of 25 at $\alpha \geq 10$. Restricting $\alpha$ to
$[0.3, 3]$ gives agreement in 71 of 72 groups and to $[0.1, 10]$ in 230 of 237.
Over the whole grid agreement is 437 of 601, and the disagreement sits where
the objective is far smaller than the constraint data, at 12 of 120 groups
whose smallest $\alpha$ is $0.01$ and 71 of 120 whose smallest is $0.03$, which
is where the stopping test loses scale invariance.

% source: results/finish_20260910/omega_4090.json and omega_bw.json.
% THE PROTOCOL, fixed once for the whole paper and stated in words in every
% caption that reports these numbers (0-abstract, 1-introduction Table 2,
% 5-experiments Tables 4 and 6, 6-conclusion Table 2):
%   metric: certified iteration count N, not wall clock;
%   instance: a (cell, seed, alpha) triple;
%   oracle: the minimum certified N over that instance's own ratio grid;
%   capping: a run that fails to certify inside the 60000-iteration cap is
%     charged 60000 iterations, so every policy is scored on every instance;
%   rule: C/omega0^p snapped to the nearest grid ratio in log space;
%   fixed baseline: the single grid ratio minimising geometric mean regret on
%     the evaluation set itself, chosen separately for each set, which favours
%     the baseline.
% There is no committed analysis script; this was recomputed from the rows.
% On wall clock from the identical rows the same computation gives 1.626 and
% 16.92 for r = 0.01, which is also the best fixed ratio on that metric, and
% the censoring counts move to 88 at the floor and 11 at the ceiling. That is
% why the metric is named in the caption.
Table~\ref{tab:regret} scores policies against a per-instance oracle. The best
single global ratio leaves $1.851\times$ geometric-mean regret and $26.04\times$
worst case, and a rule that reads the ratio off measured structure,
$r = C / \omega_0^2$, cuts both.

\begin{table}[t]
\centering
\caption{Regret in iteration count against a per-instance oracle over the 209
instances of the nine-point ratio grid, under the protocol used in every table
of this paper that reports regret. The metric is $N$ and not wall clock; the
oracle is the minimum certified iteration count over that instance's own ratio
grid; a run that fails to certify inside the 60000-iteration cap is charged
60000 iterations, so all three policies are scored on all 209 instances. The
certified column is the count each policy solved inside the cap. Both constants
in the rule rows are fitted on these same 209 instances, so those two rows are
in sample; the constant carried onto unseen families is the separate fit of
Table~\ref{tab:transfer}.}
\label{tab:regret}
\begin{tabular}{@{}llrrr@{}}
\toprule
Policy & Constant & Geomean regret & Worst case & Certified \\
\midrule
Global constant ratio      & $r = 0.01$   & 1.851 & 26.04 & 208 \\
$r = C / \omega_0^2$, geomean-optimal & $C = 0.263$  & 1.181 & 3.87 & 209 \\
$r = C / \omega_0^2$, worst-case-optimal & $C = 0.037$ & 1.190 & 3.18 & 209 \\
\bottomrule
\end{tabular}
\end{table}

% Recomputed under the protocol above: 1.85083 / 1.18080 = 1.5675;
% 26.04167 / 3.86667 = 6.7349; 1.85083 / 1.18998 = 1.5553;
% 26.04167 / 3.17647 = 8.1984; r = 0.3 gives geomean 2.0403 and worst 6.5833,
% certifying 180 of 209, and 6.5833 / 3.17647 = 2.0725. The worst-case-optimal
% constant is flat: every C from 0.0295 to 0.0407 gives the identical worst case
% 3.17647, so only two digits of C are printed. Scoring each policy only where
% it certifies instead returns 1.8274 and 23.03 for r = 0.01 on 208 instances.
The improvement over the best fixed choice is $1.567\times$ on geometric mean
and $6.73\times$ on worst case at the geomean-optimal constant, and
$1.555\times$ and $8.20\times$ at the worst-case-optimal constant. Both of
those worst-case figures are measured against $r = 0.01$, which is the
geomean-optimal fixed ratio and not the worst-case-optimal one. The
worst-case-optimal fixed ratio is $r = 0.3$, with geometric mean $2.040$ and
worst case $6.583$, certifying 180 of 209, so the like-for-like worst-case
improvement is $2.07\times$. Scoring each policy only where it certifies,
instead of charging capped runs at the cap, moves the constant-ratio row to
$1.827\times$ and $23.03\times$ on 208 instances.

% source: results/finish_20260910/omega_4090.json and omega_bw.json
One caveat travels with this result and Section~\ref{sec:exp:transfer} takes it
up. The ratio grid censors the optimum: 91 of 209 instances have their best
ratio at the floor of $0.01$ and 11 at the ceiling of $100$, so 102 of 209
optima sit on a grid edge, counting an instance as censored when any ratio
attaining its minimum certified $N$ sits at an endpoint. A naive log-log fit of
the optimal ratio against $\omega_0$ then returns a slope of $-0.9365$ under the
geometric tie average this paper uses throughout, against $-0.9309$ with ties
broken to the smallest ratio and $-0.9421$ to the largest, none of which is the
exponent.

\subsection{Transfer of the step-ratio rule across families}
\label{sec:exp:transfer}

% source: results/finish_20260910/omega_wide_bw.json (fit set, 1404 rows)
% The functional form is not ours: r = C / omega0^2 with omega0 = ||c||/||b||
% is the primal weight of PDLP, cited in Section 2.
The functional form of this rule is not new. A step ratio set as
$C / \omega_0^2$ with $\omega_0 = \lVert c \rVert / \lVert b \rVert$ is the
primal-weight heuristic of PDLP \citep{Applegate2021Practical}, which this paper
already credits. What is measured here is how much it is worth against the best
fixed ratio and whether one constant transfers across families.

% source: results/finish_20260910/omega_wide_bw.json.
% Fitting p and C jointly on this set alone by minimising geometric-mean regret,
% over p on a 0.05 grid and C on a 0.01 dex grid, gives p = 2.00 and C = 0.9772
% at geometric-mean regret 1.18583. The argmin in C is a single grid point at
% that resolution: the neighbouring points 0.9550 and 1.0000 give 1.18894 and
% 1.19487. The earlier claim of a plateau from 0.9772 to 0.9886 does not
% reproduce and is deleted.
% Refining to p on a 0.01 grid and C on a 0.001 dex grid moves the argmin to
% p = 1.98, C = 1.040, at regret 1.18573, which is 0.008 percent lower. The
% coarser argmin is kept because that difference does not identify a rule.
% This pair, and no other, is used in the abstract, the introduction and the
% conclusion. The earlier drafts printed C = 1.10 at p = 2 (regret 1.1881) and
% C = 1.16 at p = 1.8 (regret 1.1981); both are worse on the set they claim to
% be fitted on and both are deleted.
We widen the netflow ratio grid to 13 points spanning $10^{-4}$ to $10^{2}$ and
fit $r = C / \omega_0^p$ on that set alone, by minimising geometric-mean regret
over $p$ on a $0.05$ grid and $C$ on a $0.01$ dex grid.
The fit returns $p = 2.00$ and $C = 0.977$, at geometric-mean regret $1.186$.
That pair is then carried unchanged onto three families the fit never saw, and
it is the only rule this paper reports.

% Table~\ref{tab:transfer} uses the protocol stated above Table~\ref{tab:regret},
% unchanged: charge a non-certifying run at the 60000-iteration cap, score every
% policy on every instance, and pick the fixed ratio per evaluation set.
Table~\ref{tab:transfer} is the result, under the protocol already used for
Table~\ref{tab:regret}. Policies are scored in iteration count against a
per-instance oracle, the minimum certified $N$ over that instance's own ratio
grid. An instance is a (cell, seed, $\alpha$) triple. A run that fails to
certify inside the 60000-iteration cap is charged 60000 iterations, so both
policies are scored on every instance of every set. The fixed-ratio arm is the
single grid ratio minimising geometric-mean regret on the evaluation set
itself, chosen separately for each set, which favours the baseline; its
certified count is reported alongside the rule's because the transferred rule
certifies every instance of every set.

\begin{table}[t]
\centering
\caption{The $\omega_0$ rule carried off the family it was fitted on. The rule
is $r = 0.977 / \omega_0^2$, fitted on netflow wide alone and then evaluated
unchanged. Regret is in iteration count against a per-instance oracle over that
instance's own ratio grid, with a run that fails to certify inside the
60000-iteration cap charged at the cap, so both policies score every instance of
every row. The fixed column is the single grid ratio minimising geometric-mean
regret on that set, chosen separately per set, and the certified columns are the
instances each of the two policies solved inside the cap. The pooled row is the
geometric mean over the 162 held-out instances, each family carrying its own
best fixed ratio.}
\label{tab:transfer}
\begin{tabular}{@{}lrlrrrrr@{}}
\toprule
& & & \multicolumn{2}{c}{geomean regret} & & \multicolumn{2}{c}{certified} \\
\cmidrule(lr){4-5} \cmidrule(lr){7-8}
Evaluation set & $n$ & $\omega_0$ range & fixed & rule & rule worst & rule & fixed \\
\midrule
netflow wide (fitted on) & 108 & 0.0697 to 825 & 2.779 & 1.186 & 2.12 & 108 & 100 \\
optimal transport        &  54 & 0.0054 to 79.3 & 1.428 & 1.179 & 1.56 &  54 & 54 \\
multi-commodity flow     &  54 & 0.0227 to 312 & 2.840 & 1.032 & 1.26 &  54 & 48 \\
held-out topologies      &  54 & 0.0735 to 806 & 2.514 & 1.190 & 1.83 &  54 & 44 \\
\midrule
pooled held out          & 162 & 0.0054 to 806 & 2.168 & 1.131 & 1.83 & 162 & 146 \\
\bottomrule
\end{tabular}
% source: results/finish_20260910/omega_wide_bw.json (108 instances, 13 ratios)
% source: results/finish_20260910/omega_ot_wide.json (54 instances, 13 ratios)
% source: results/finish_20260910/omega_mcf.json (54 instances, 12 ratios)
% source: results/finish_20260910/omega_holdout.json (54 instances, 13 ratios)
% Recomputed to four decimals: rule geomean 1.1858 / 1.1786 / 1.0318 / 1.1899
%   / 1.1311; rule worst 2.1250 / 1.5625 / 1.2623 / 1.8252 / 1.8252.
% Best fixed under the protocol in the caption: 2.7790 at r = 0.03 certifying
%   100 of 108; 1.4279 at r = 100 certifying 54 of 54; 2.8395 at r = 0.1
%   certifying 48 of 54; 2.5136 at r = 0.003 certifying 44 of 54; pooled 2.1681
%   certifying 146 of 162. A single ratio shared across the three held-out
%   families, drawn from the 11 ratios their grids have in common, is worse
%   still, r = 3 at 2.8204.
% Under the alternative protocol, scored only where certified with an 80 percent
%   coverage floor: 2.1513 at r = 0.0003 on 87 of 108; 1.4279 at r = 100 on
%   54 of 54; 2.1854 at r = 0.001 on 44 of 54; 1.9475 at r = 0.0001 on 45 of 54;
%   pooled per family 1.8249 on 143 of 162. One shared ratio under the same
%   80 percent floor gives r = 0.1 at 2.5039 on 130 of 162; with the floor
%   dropped, r = 30 at 1.7516 on 108 of 162.
% omega_0 endpoints to four decimals: 0.0697 to 825.2357; 0.0054 to 79.2830;
%   0.0227 to 311.5050; 0.0735 to 806.2762; 0.0054 to 806.2762. These are over
%   all seeds; computing them on one seed gives different and wrong endpoints.
\end{table}

The gain over the best fixed ratio is $2.34\times$ on the set the rule was
fitted on, $1.21\times$ on optimal transport, $2.75\times$ on multi-commodity
flow and $2.11\times$ on the held-out topologies, and $1.92\times$ pooled over
the 162 held-out instances. The comparison is protocol dependent and we give the
other reading here rather than in a second table. Under the alternative protocol
in which each policy is scored only on the instances it certifies, with the
fixed ratio restricted to those certifying at least 80 percent of the set, the
fixed baselines improve to $2.151\times$ on 87 of 108 netflow wide instances,
$1.428\times$ on 54 of 54 optimal transport, $2.185\times$ on 44 of 54
multi-commodity flow and $1.947\times$ on 45 of 54 held-out topologies, a pooled
$1.825\times$ on 143 of the 162 held-out instances. They are then measuring a
different and smaller problem on every row, which is why the tables do not use
that protocol. Dropping the coverage floor as well lets one shared ratio certify
barely two thirds of the pooled set, $r = 30$ at $1.752\times$ on 108 of 162.
The rule is untouched by the choice at $1.186$, $1.179$, $1.032$, $1.190$ and
$1.131$, because it certifies every instance of every set.

% source: results/finish_20260910/omega_wide_bw.json, omega_ot_wide.json,
% omega_mcf.json, omega_holdout.json. Native omega0 at alpha = 1:
% netflow 6.9691 to 8.2524; optimal transport 0.5419 to 0.7928;
% multi-commodity flow 2.2666 to 3.1151; held-out topologies 7.3548 to 8.0628.
The sample structure has to be stated plainly, because it bounds what transfer
means here. Each set is (cells $\times$ seeds) distinct problems replicated at
nine synthetic rescalings of $c$: netflow wide is 12 distinct problems at 9
values of $\alpha$, and each held-out set is 6 at 9. Variation in $\omega_0$
within a set is therefore variation in $\alpha$, and $c \mapsto \alpha c$ is
exactly $(\tau, \sigma) \mapsto (\alpha \tau, \sigma / \alpha)$. The native
$\omega_0$ spread inside netflow is only $6.97$ to $8.25$, a factor of $1.18$,
so the rule is being tested on rescalings and not on naturally occurring scale
variation. The part that is not a rescaling is the comparison across families:
optimal transport sits at native $\omega_0$ of $0.54$ to $0.79$ and
multi-commodity flow at $2.27$ to $3.12$, both genuinely away from netflow's
$7.5$, and the transferred constant is not refitted for either.

% source: the same four files. Refitting C per family at p = 2 by the same
% grid search: 0.977 on netflow wide, 31.62 on optimal transport, 1.380 on
% multi-commodity flow, 0.275 on the held-out topologies. 31.62 / 0.275 = 115.
One constant does not fit every family. Refitting $C$ per family at $p = 2$
gives $0.977$ on netflow wide, $31.62$ on optimal transport, $1.380$ on
multi-commodity flow and $0.275$ on the held-out topologies, a span of a factor
of 115 between the smallest and the largest. What Table~\ref{tab:transfer}
reports is the netflow constant used unchanged, and the residual gap on
optimal transport, $1.179$ against a per-family-optimal $1.026$, is the price
of that transfer.

% source: results/finish_20260910/omega_4090.json and omega_bw.json (narrow),
% omega_wide_bw.json (wide), omega_seeds8.json (8 seeds),
% omega_ultrawide_bw.json (ten decades of r, eleven alphas, 4 seeds),
% omega_fine_4090.json (finer ratio grid).
% CONVENTIONS, the same two used in 0-abstract, 1-introduction and 6-conclusion
% and named in the caption of Table~\ref{tab:exponent}:
%   censoring: an instance counts as censored when ANY ratio attaining its
%     minimum certified N sits at an endpoint of that instance's grid;
%   slope ties: the log-log fit of the argmin ratio against omega0 runs over the
%     instances with a converged run, with tied argmin ratios averaged
%     geometrically. Under the other two tie rules, ties to the smallest and
%     ties to the largest, the five slopes move by at most 0.011.
% Regret-optimal exponent by grid search over p on a 0.05 grid at the best C on
% a 0.01 dex grid, the same resolution as the fit above.
% Slopes to four decimals: -0.9365, -1.5412, -1.5781, -1.6834, -1.7380.
% Censoring to the convention above: 102 of 209 (91 floor, 11 ceiling), 33 of
% 108, 46 of 144, 15 of 80, 14 of 88.
Two estimators of the exponent disagree, and reducing censoring moves only one
of them. Table~\ref{tab:exponent} runs both over five ratio grids of increasing
width. The log-log slope of the argmin ratio against $\omega_0$ rises in
magnitude from $-0.94$ on the original grid, where 102 of 209 optima sit on a
grid edge, to $-1.74$ on the ultrawide ten-decade grid, where 14 of 88 do. The
regret-minimising exponent does not move with it: it is $2.10$, $2.00$, $1.90$,
$2.00$ and $2.00$ on the same five grids. The honest reading is that censoring
biases the slope estimator toward zero, that removing censoring moves it from
$-0.94$ toward $-1.74$ without reaching the regret estimator, and that the two
therefore bracket the exponent between about $1.7$ and $2.1$ and have not met.
That bracket is consistent with the value 2 implied by the primal weight, and at
its lower end below it; we do not claim the exponent is exactly 2, and no
estimator on any of the five grids returns $1.8$ as its argmin. The regret
profile is flat enough to say how finely the data resolve it. On the original
grid the set of exponents within 1 percent of the optimum is not even an
interval: it runs from $1.40$ to $2.90$ with an isolated point at $1.15$. On the
ultrawide grid, which is the least censored, that set is $1.80$ to $2.25$, with
$p = 1.8$ costing $0.8$ percent and $p = 2.4$ costing $2.2$ percent. That is the
regret estimator alone; the verdict this paper carries is the two-estimator
bracket above, about $1.7$ to $2.1$. The finer 19-point grid has now run and is
the fourth row of the table; it does not narrow
the bracket, returning $-1.68$ and $2.00$, between the wide and ultrawide grids
on both estimators.

\begin{table}[t]
\centering
\caption{Two estimators of the exponent in $r = C / \omega_0^p$, over five
ratio grids, under the two conventions used wherever this paper reports them.
Censoring is the count of instances for which any ratio attaining the minimum
certified $N$ sits at an endpoint of that instance's grid. The slope is a
log-log fit of the argmin ratio against $\omega_0$ over the instances with a
converged run, with tied argmin ratios averaged geometrically; the other two tie
rules move each slope by at most $0.011$. The regret exponent minimises
geometric-mean regret in iteration count, searched over $p$ on a $0.05$ grid at
the best constant on a $0.01$ dex grid for each exponent.}
\label{tab:exponent}
\begin{tabular}{@{}lrlrrr@{}}
\toprule
Grid & $n$ & Ratio range & Censored & Slope & Regret exponent \\
\midrule
original    & 209 & $10^{-2}$ to $10^{2}$  & 102 & $-0.94$ & 2.10 \\
netflow wide & 108 & $10^{-4}$ to $10^{2}$ & 33 & $-1.54$ & 2.00 \\
8 seeds     & 144 & $10^{-4}$ to $10^{3}$  & 46 & $-1.58$ & 1.90 \\
finer grid  &  80 & $10^{-5}$ to $3 \times 10^{2}$ & 15 & $-1.68$ & 2.00 \\
ultrawide   &  88 & $10^{-6}$ to $10^{4}$  & 14 & $-1.74$ & 2.00 \\
\bottomrule
\end{tabular}
% source: results/finish_20260910/omega_4090.json + omega_bw.json; omega_wide_bw.json;
% source: results/finish_20260910/omega_seeds8.json; omega_fine_4090.json; omega_ultrawide_bw.json
% Slopes with ties averaged geometrically: -0.9365, -1.5412, -1.5781, -1.6834,
%   -1.7380. Ties to the smallest ratio: -0.9309, -1.5307, -1.5716, -1.6762,
%   -1.7324. Ties to the largest: -0.9421, -1.5516, -1.5845, -1.6905, -1.7437.
% Best regret at the argmin exponent, C on a 0.01 dex grid: 1.1805, 1.1858,
%   1.1571, 1.1685, 1.0618.
% Ultrawide profile at the same resolution: p=1.8 gives 1.0702, p=2.0 gives
%   1.0618, p=2.4 gives 1.0853. On a 0.001 dex grid in C the p=2.0 value is
%   1.0574; every figure printed in the paper uses the 0.01 dex grid.
% Within 1 percent of the optimum, by grid: {1.15} together with 1.40 to 2.90;
%   1.85 to 2.15; 1.65 to 2.45; 1.75 to 2.05; 1.80 to 2.25.
% The finer grid is omega_fine_4090.json: 19 ratios from 1e-5 to 300, spaced by
%   a factor of about 2.6. It has run; no sentence in this paper says otherwise.
\end{table}

\subsection{Negative results}
\label{sec:exp:negative}

% source: results/diagnostics_20260909/roofline_4090.json (model inputs) and
% results/diagnostics_20260909/REPORT.md lines 26-29 (prose table, all four
% rows including the 0.567 fitted-launch-cost row). The four scores were
% reproduced on 2026-09-10 by running scripts/diagnostics/score_roofline.py,
% whose stdout is archived at
% results/artifacts_20260910/roofline_scores_4090.log, 26 lines: mean Spearman
% 0.355, 0.363, 0.282 and 0.567 with pick rates 1/7, 0/7, 0/7 and 1/7, and M1
% per-size correlations -0.943 at n=128 and +0.829 at 992, 1488 and 1984.
Cheap a-priori predictors fail on the task where $\omega_0$ succeeds.
Table~\ref{tab:roofline} scores four analytic models of the fusion tile against
measurement by rank correlation over 7 sizes and 6 tiles.

\begin{table}[t]
\centering
\caption{Analytic tile selection against measurement. Occupancy raises the rank
correlation from $0.355$ to $0.363$ and lowers the best-tile hit rate from 1 of
7 to 0 of 7. All four models pick the best tile at most 1 time in 7, which is
the load-bearing figure.}
\label{tab:roofline}
\begin{tabular}{@{}lrr@{}}
\toprule
Model & Mean Spearman & Picked the best tile \\
\midrule
Roofline only                        & 0.355 & 1 of 7 \\
Plus occupancy                       & 0.363 & 0 of 7 \\
Plus occupancy and wave quantisation & 0.282 & 0 of 7 \\
Plus a fitted per-launch cost of $10^{-5}$\,s & 0.567 & 1 of 7 \\
\bottomrule
\end{tabular}
\end{table}

% source: results/diagnostics_20260909/roofline_4090.json
% All four device constants verified: 746.6579 GB/s, 1.1402 TFLOP/s, ridge
% 1.5271 flop/byte, unfused intensity 0.5536, 24x4 intensity 2.7253.
The interpretable structure is a change of regime. The roofline-only model
scores $-0.943$ at $n=128$ and $+0.829$ at $n \geq 992$, so it is right only
once the kernel is bandwidth-bound and is anti-correlated when the kernel is
overhead-bound. The device constants are 746.7\,GB/s achievable float64
bandwidth, 1.140\,TFLOP/s float64 and a ridge point of 1.527 flop per byte, on
an accounting of 31 flops and 56 bytes per pixel-iteration; the unfused TV
iteration sits at 0.554 flop per byte and the $24\times4$ tile reaches 2.725,
above the ridge. Two caveats belong in the text and not in a comment: the flop
count is hand-counted from \texttt{tv\_update.cu} and not instrumented, and
occupancy is computed from assumed Ada AD102 constants rather than from a
device query.

% source: results/diagnostics_20260909/sharpness_full_4090.json, recomputed.
% Log-log Pearson over the converged rows: pooled +0.9426 (17), netflow -0.5387
% (13), regular -0.8423 (6), grid -0.1197 (7), TV +0.9825 (4).
Empirical sharpness looks predictive and is not. Over 17 converged instances the
correlation between log sharpness and log iteration count is $+0.943$ pooled,
$-0.539$ within min-cost flow over 13 instances, $-0.842$ within the regular
topology over 6, $-0.120$ within the grid topology over 7 and $+0.983$ within
TV over 4. That last figure is a four-point correlation spanning only two image
sizes and is really tracking size, as the source report's own caveat says. The
pooled figure only separates two clusters, min-cost flow at mean sharpness 92.1
and mean 10230 iterations against TV at 43.3 and 624. Within a family the sign
inverts.

% source: results/diagnostics_20260909/sharpness_full_4090.json, recomputed.
% Quantity: log of the fitted decay rate of the gap over the stated window of
% certificate checks, against log N, over the 13 converged min-cost flow rows.
% Recomputed: +0.5387, +0.2006, +0.0177, -0.2751, -0.3686, -0.6201, and
% -0.7551 for the last window when the thirteenth instance is scored on its
% truncated window instead of dropped.
A leakage-free version of the estimator has to watch most of the solve before it
becomes even moderately predictive. Estimating the gap decay rate on a fixed
index window of certificate checks and correlating its logarithm against log
iteration count over the min-cost flow instances gives $+0.539$ over the first 3
checks, $+0.201$ over 5, $+0.018$ over 10, $-0.275$ over checks 10 to 20,
$-0.369$ over 20 to 40 and $-0.620$ over 40 to 80. The sign flips halfway
through: early in the solve a faster measured decay goes with more iterations
and not fewer, and only over checks 40 to 80 does the windowed estimate acquire
the negative sign that the full-solve sharpness estimate carries within
min-cost flow, $-0.539$. That window is scored on the 12 of 13 instances
that record at least 80 certificate checks; scoring the thirteenth on its
truncated window instead gives $-0.755$. Checks are 64 iterations apart, so the
window ends at 5120 iterations against a median min-cost flow solve of 9600, and
the estimator reaches $|r| = 0.620$ only after 53 percent of the run.

% source: results/structure_20260909/structure_summary_4090.json
The one structure-dependent selector that does work is the construction gate.
It is exact, symbolic and free, and admits 1755 of 2106 settings purely as a
function of the gate product, 351 of 351 at each of $\theta \in \{0.1, 0.3,
0.6, 0.9, 0.99\}$ and 0 of 351 at $\theta = 1.02$. Rejected settings were never
executed, so the gate is verified as sufficient for non-divergence and never as
necessary, and 257 of the admitted runs still hit the iteration cap. It is a
feasibility test and not a rate test, so it tells a search which recipes are
legal and nothing about which legal recipe is fast. That is the hinge between
these negatives and the case for search.

\subsection{Measurement conditions}
\label{sec:exp:conditions}

% source: results/tv_ffi_20260906/DECOMPOSITION_REPORT.md lines 6 to 8. The
% 12-process stability run has NO saved artifact; it is prose in that report,
% for the f64_K64 cell at n=256.
% source: results/finish_20260910/fusion_r9_sz256.json and fusion_r9_sz496.json
% against results/kernel_fusion_20260909/e2e/full_4090_sz{256,496}.json.
% Recomputed drift: baseline +1.97 and +1.07 percent; static over traced
% 1.3470 -> 1.4008 (+3.99 percent) and 1.0933 -> 1.1788 (+7.81 percent);
% best tile 1.8675 -> 1.9396 and 1.3556 -> 1.4279.
Timing runs were not on exclusive machines. One dispersion figure is available
as prose only: the reference cell \texttt{f64\_K64} at $n=256$, re-measured in
12 independent processes, spans 20.297 to 20.432 microseconds per iteration, a
spread of 0.66 percent, and that run has no saved artifact. A recomputable
dispersion is available and is larger. Re-running the $n=256$ and $n=496$
fusion sweeps at 9 repeats one day later moves the traced baseline by 2.0 and
1.1 percent, the static-over-traced gain from $1.347$ to $1.401$, a change of
4.0 percent, and from $1.093$ to $1.179$, a change of 7.8 percent, and the
best-tile speedup from $1.868$ to $1.940$ and from $1.356$ to $1.428$. No
fusion figure in this section is stable in its third significant figure. The
fusion tables print three digits so that the arithmetic can be checked against
the artifacts, and only the first two should be read as measured.

% source: results/method_axis_20260909/data.json ('final_holdout_loaded': false)
% source: results/hardware_codex/REPORT.md line 7 ("The original
% pre-registration and final holdout remain untouched.")
The frozen final holdout has never been evaluated, and none of the figures in
this section touch it. The families described as held out in
Section~\ref{sec:exp:transfer} are held out from the rule fit only; they are not
the frozen paper holdout.

A vLLM server held 15.4\,GB resident on the timed 4090 at 0 percent
utilisation throughout. The Blackwell box is shared: the fusion microbenchmarks
used GPU 2 at 0 percent utilisation but not exclusively, and all four GPUs were
idle for the cross-architecture end-to-end runs.

\section{Discussion}

\subsection{What the factorisation buys}

The factorisation lets the two factors be searched separately. The evidence for
the first factor is exact rather than statistical. Across every end-to-end
fusion sweep, seven sizes from 128 to 1984 pixels on two devices, the
traced-bound baseline, the static chunked loop and every fusion tile reach the
same iteration count on the same instance and all certify: 72 of 72 instances,
eight arms each, zero iteration mismatches, zero certification failures. All
fourteen sweeps ran at $K=256$.
% source: results/kernel_fusion_20260909/e2e/full_4090_sz{128,256,384,496}.json
% source: results/kernel_fusion_20260909/e2e_bw/full_bw_sz{128,256,384,496}.json
% source: results/kernel_fusion_20260909/large/large_4090_m{2,3,4}.json
% source: results/kernel_fusion_20260909/bw_day2/large_bw_m{2,3,4}.json
% 6 instances in each full_* file, 4 in each large_* file, 8 arms each: 72 total.
The iteration count was also identical on the two architectures measured, at
every size run on both: $576.000$ and $682.667$ at $n{=}128$, $597.333$ and
$682.667$ at $n{=}256$, and $1066.667$ and $1152.000$ at $n{=}496$, for $K=64$
and $K=256$ at tolerance $10^{-4}$, on both sm\_89 and sm\_120. At $n{=}128$
both devices give the same per-instance vectors,
$[1216, 320, 448, 448, 384, 640]$ and $[1280, 512, 512, 512, 512, 768]$.
% source: results/tv_size_20260909/sz{128,256,496}_4090/factorial.json
% source: results/bw_20260909/sz{128,256,496}_bw/factorial.json
Kernel search can therefore hold $N$ fixed and time a single arithmetic budget,
within the two-device scope in which that equality was measured.

The converse half holds under fixed precision, conditioned on cadence. At
$K=256$, step sizes that change $N$ by $4.3125\times$ ($2944.000$ against
$682.667$) change the cost per iteration by 3.8 percent, from $19.0766$ to
$19.7976$ microseconds across the four \texttt{f64\_K256} arms, a factor of
$1.0378$ (Table~\ref{tab:leak}). At $K=64$ the same two step-size presets give
$4.839\times$ ($2890.667$ against $597.333$). Scheme identity is inert to within
one cadence block at $K=64$ and $K=256$: the stopping test is evaluated only at
multiples of $K$, so $N$ is the crossing iteration rounded up to the next
multiple of $K$, and at $K=256$ the equality is resolved only to 256 iterations.
Inside that resolution \texttt{condat\_vu\_banked} and \texttt{pdhg\_banked}
share the per-instance vector $[1024, 512, 768, 512, 512, 768]$ at $n{=}256$ and
agree on 11 of 12 certified netflow pairs; the twelfth, size 12 instance 1,
agrees at the 20000-iteration cap with \texttt{conv} false and gap
$2.858 \times 10^{-4}$ for both schemes, so it is an agreement about failure
rather than about $N$. At $K=16$ and $n{=}256$ the two schemes do differ,
$584.000$ against $578.667$.
% source: results/method_axis_20260909/factorial.json, tolerances["0.0001"].rows
% source: results/netflow_20260909/grid_4090.json
% source: results/bw_20260909/sz256_bw/factorial.json;
% source: results/tv_ffi_20260906/cadence_4090_v1/factorial.json

Cadence carries the largest single share of the variance in log cost per
iteration, 47.43 percent at $10^{-4}$ and 51.19 percent at $10^{-6}$.
% source: results/artifacts_20260910/fanova.json,
% decompositions.exec_x_prec_x_cadence__logT__tol0.0001.percent_variance.cadence
%   = 47.43020732030952
% decompositions.exec_x_prec_x_cadence__logT__tol1e-06.percent_variance.cadence
%   = 51.19205833962386
% 12 cells, execution x precision x cadence, response log per-iteration cost.
The mechanism is chunk length, not check frequency. Measured directly over the
24 certificate-on minus certificate-off differences in a loop-structure sweep,
the certificate term ranges from $-239.90$ to $+844.56$ nanoseconds per
iteration against a static per-iteration cost of $8.09$ to $11.94$ microseconds,
and at $n{=}128$, $K=64$ it is $26.5$ nanoseconds. At $K \geq 64$ the check
itself is inside noise; what cadence buys is the length of the chunk that can be
unrolled between checks.
% source: results/finish_20260910/loop_mechanism.json (48 cells, 24 cert differences)

\subsection{Where the axes stop separating}

The structure predicts its own failure modes and the measurements find both. An
adaptive precision schedule couples the factors, because the switch iteration is
chosen from the gap trajectory, so the fraction of the run spent in float32
depends on the step sizes: the 3.8 percent spread under fixed precision becomes
16.6 percent under the schedule (Table~\ref{tab:leak}). The two execution
decisions also stop being independent past about 1000 pixels. The static loop's
advantage over the traced bound falls from $1.673\times$ at 128 pixels through
$1.347$, $1.291$, $1.093$ and $1.016$ to a minimum of $0.989\times$ at 1488 and
$0.993\times$ at 1984, so at the two largest sizes the static bound is mildly
harmful rather than merely exhausted. Fusion over the static loop instead turns
U-shaped, at $1.571\times$, $1.060\times$ and $1.410\times$ across 128, 992 and
1984 pixels. Each of those is a ratio of summed wall time over the instances of
that sweep rather than a mean of per-instance ratios; the per-instance geometric
mean at 128 pixels is $1.580\times$. The 128-pixel sweep uses six raw crops and
the 992 to 1984 pixel sweeps use four montages each, of two to four distinct
scans, since a raw scan caps an honest crop near 496 pixels. The $1.410\times$
at 1984 pixels is a fusion gain over the static loop at float64 throughout,
$2.2407709$ seconds static over $1.5890337$ seconds for the $24 \times 4$ tile,
and is not a precision comparison. The execution structure and the tile must be
chosen jointly at large size and can be chosen independently at small.
% source: results/kernel_fusion_20260909/e2e/full_4090_sz128.json (128 px)
% source: results/kernel_fusion_20260909/e2e/full_4090_sz{256,384,496}.json
% source: results/kernel_fusion_20260909/large/large_4090_m{2,3,4}.json (992, 1488, 1984 px)
% source: results/kernel_fusion_20260909/REPORT.md section 7

\begin{table}[t]
\centering
\caption{The factorisation and its one measured leak. TV, $n{=}256$, tolerance
$10^{-4}$, mean over six OCT images. Under fixed precision the cost per
iteration spans 3.8 percent across arms whose iteration counts differ by
$4.3125\times$. Under the adaptive schedule it spans 16.6 percent.}
\label{tab:leak}
\begin{tabular}{llrr}
\toprule
policy & algorithm & mean $N$ & cost per iteration ($\mu$s) \\
\midrule
\texttt{f64\_K256} & \texttt{condat\_vu\_banked}   & 682.667  & 19.798 \\
\texttt{f64\_K256} & \texttt{pdhg\_banked}         & 682.667  & 19.683 \\
\texttt{f64\_K256} & \texttt{condat\_vu\_balanced} & 2944.000 & 19.077 \\
\texttt{f64\_K256} & \texttt{pdhg\_balanced}       & 2944.000 & 19.227 \\
\midrule
\texttt{schedule\_K256} & \texttt{condat\_vu\_banked}   & 682.667  & 11.052 \\
\texttt{schedule\_K256} & \texttt{pdhg\_banked}         & 682.667  & 11.232 \\
\texttt{schedule\_K256} & \texttt{condat\_vu\_balanced} & 2944.000 & 12.578 \\
\texttt{schedule\_K256} & \texttt{pdhg\_balanced}       & 2944.000 & 12.888 \\
\bottomrule
\end{tabular}
% source: results/method_axis_20260909/factorial.json, tolerances["0.0001"].rows.
% Cost per iteration is solver_s / iters per instance, averaged over the 6 training images.
% Recomputed to full precision: 19.797606 / 19.682557 / 19.076624 / 19.227495 /
% 11.051901 / 11.231891 / 12.577807 / 12.888450 microseconds.
% Spread 19.076624 to 19.797606 is a factor 1.0378; 11.051901 to 12.888450 is 1.1662.
\end{table}

\subsection{Predicting the step ratio from problem structure}

Rescaling the objective vector as $c \to \alpha c$ is identically the step
rescaling $(\tau, \sigma) \to (\alpha \tau, \sigma / \alpha)$, so the
construction gate is invariant under it while the step ratio $r = \tau / \sigma$
and the conditioning proxy $\omega_0 = \lVert c \rVert / \lVert b \rVert$ both
move. Sweeping $\alpha$ over decades therefore separates the ratio from the
gate. Fitting a rule of the form $r = C / \omega_0^{p}$ on the netflow wide
sweep alone, by minimising geometric-mean regret over its 13-point ratio grid,
with $p$ searched on a $0.05$ grid and $C$ on a $0.01$ dex grid, gives
$p = 2.00$ and $C = 0.977$. That is the only rule this paper reports, and it is
the pair used in the abstract, the introduction and Section~5. Evaluated
unchanged on three held-out families, it beats the best single fixed ratio on
every one of them (Table~\ref{tab:omega}), and the cleanest statement is
coverage rather than regret: the rule certifies every instance of all four sets,
while the best fixed ratio for each set certifies 100 of 108, 54 of 54, 48 of 54
and 44 of 54.

The comparison is protocol dependent and both protocols are reported, in the
same words as Section~5. Under the protocol of Table~\ref{tab:omega}, which is
the one every table in this paper uses, a run that fails to certify inside the
60000-iteration cap is charged at the cap, so both policies are scored on the
identical instance set. Under the alternative protocol in which each policy is
scored only on the instances it certifies, with the fixed ratio restricted to
those certifying at least 80 percent of the set, the fixed baselines improve to
$2.151\times$ on 87 of 108 netflow wide instances, $1.428\times$ on 54 of 54
optimal transport, $2.185\times$ on 44 of 54 multi-commodity flow and
$1.947\times$ on 45 of 54 held-out topologies, a pooled $1.825\times$ on 143 of
the 162 held-out instances. They are then measuring a different and smaller
problem on every row. Dropping the coverage floor as well lets one shared ratio
certify barely two thirds of the pooled set, $r = 30$ at $1.752\times$ on 108 of
162. The exponent 2 is the value implied by the primal-weight heuristic in PDLP;
the evidence here is consistent with it without pinning it down, and
Section~\ref{sec:limits} gives the identification caveat in full.
% source: results/finish_20260910/omega_wide_bw.json (fit), omega_ot_wide.json,
% source: results/finish_20260910/omega_mcf.json, omega_holdout.json
% p and C fitted jointly by grid search on omega_wide_bw only, p on a 0.05 grid
% and C on a 0.01 dex grid: p = 2.00, C = 0.9772, geometric mean regret 1.18583.
% The alternative-protocol figures above recompute as 2.1513 on 87 of 108,
% 1.4279 on 54 of 54, 2.1854 on 44 of 54, 1.9475 on 45 of 54, pooled 1.8249 on
% 143 of 162, and with the coverage floor dropped 1.7516 at r = 30 on 108 of 162.
% They are identical to the figures printed in Section 5, by construction.

\begin{table}[t]
\centering
\caption{The $\omega_0$ rule transfers off the family it was fitted on. The rule
is $r = C / \omega_0^{2}$ with $C = 0.977$, fitted on netflow wide alone and
then evaluated unchanged. This is the protocol used in every table of this paper
that reports regret. Regret is against the per-instance oracle over that
instance's own ratio grid, with a run that fails to certify inside the
60000-iteration cap charged at the cap, so both policies are scored on every
instance of every row. The fixed column is the single grid ratio minimising
geometric-mean regret on that set, chosen separately per set, and both certified
columns count the instances that same policy solved inside the cap. The pooled
row is the geometric mean over the 162 held-out instances, each family carrying
its own best fixed ratio.}
\label{tab:omega}
\begin{tabular}{lrlrrrrr}
\toprule
& & & \multicolumn{2}{c}{geomean regret} & & \multicolumn{2}{c}{certified} \\
\cmidrule(lr){4-5} \cmidrule(lr){7-8}
evaluation set & $n$ & $\omega_0$ range & fixed & rule & rule worst & rule & fixed \\
\midrule
netflow wide (fitted on) & 108 & 0.0697 to 825 & 2.779 & 1.186 & 2.125 & 108 & 100 \\
optimal transport        &  54 & 0.0054 to 79.3 & 1.428 & 1.179 & 1.562 &  54 &  54 \\
multi-commodity flow     &  54 & 0.0227 to 312 & 2.840 & 1.032 & 1.262 &  54 &  48 \\
held-out topologies      &  54 & 0.0735 to 806 & 2.514 & 1.190 & 1.825 &  54 &  44 \\
\midrule
pooled held out          & 162 & 0.0054 to 806 & 2.168 & 1.131 & 1.825 & 162 & 146 \\
\bottomrule
\end{tabular}
% source: results/finish_20260910/omega_wide_bw.json (108 instances, 13 ratios, 1404 rows)
% source: results/finish_20260910/omega_ot_wide.json (54, 13 ratios, 702 rows)
% source: results/finish_20260910/omega_mcf.json (54, 12 ratios, 648 rows)
% source: results/finish_20260910/omega_holdout.json (54, 13 ratios, 702 rows)
% Recomputed to 4 decimals: rule 1.1858 / 1.1786 / 1.0318 / 1.1899 / 1.1311;
% worst 2.1250 / 1.5625 / 1.2623 / 1.8252 / 1.8252;
% best fixed 2.7790 (r=0.03) / 1.4279 (r=100) / 2.8395 (r=0.1) / 2.5136 (r=0.003),
%   pooled 2.1681 as the geometric mean of the three held-out rows over their
%   162 instances; certified by those same ratios 100 / 54 / 48 / 44 / 146.
% The pooled row of earlier drafts used one shared ratio on the 11-point common
%   grid, 2.7843 at r = 3, and the rule value 1.1105 that goes with restricting
%   the oracle to that grid. Both are dropped: the table now keeps each
%   instance's own grid throughout, which gives 2.1681 and 1.1311.
% One shared ratio, drawn from the 11 ratios the three held-out grids have in
%   common but scored against each instance's own oracle, gives 2.8204 at r = 3.
% omega_0 ranges to 4 decimals: 0.0697 to 825.2357; 0.0054 to 79.2830;
%   0.0227 to 311.5050; 0.0735 to 806.2762; 0.0054 to 806.2762.
\end{table}

\subsection{Limitations}
\label{sec:limits}

Two GPU architectures are in scope for every timing claim, an RTX 4090
(sm\_89) and an RTX PRO 6000 Blackwell Max-Q (sm\_120), so no
accelerator-portability claim is available. An RTX 5090 was measured earlier and
is excluded from scope; its rows survive in the archive and are not used here.
The two in-scope devices disagree about whether to fuse at all. At $n{=}496$ the
best Blackwell tile runs at $0.925\times$ the plain static loop and at 992 pixels
at $0.989\times$, so on that device the correct decision at two sizes is not to
fuse. The devices agree on the tile at 128 pixels ($8 \times 4$) and 256 pixels
($24 \times 4$), disagree at 384 ($24 \times 4$ against $32 \times 4$), 496
($32 \times 4$ against $16 \times 2$), 992 ($32 \times 2$ against $24 \times 4$)
and 1488 ($24 \times 4$ against $32 \times 2$), and agree again at 1984
($24 \times 4$).
% source: results/kernel_fusion_20260909/e2e_bw/full_bw_sz{128,256,384,496}.json
% source: results/kernel_fusion_20260909/bw_day2/large_bw_m{2,3,4}.json
% source: results/hardware_codex/REPORT.md (oct_5090_v1, lp_5090_v2);
% source: results/tv_ffi_20260906/catalog_5090; results/kernel_fusion_20260909/REPORT.md line 4

The frozen final holdout has never been evaluated and the pre-registration is
untouched. The field \texttt{final\_holdout\_loaded} is false in all 23
\texttt{data.json} records under \texttt{results/} that carry it; the nine older
\texttt{hardware\_codex} records predate that field and carry only a
three-patch per-run diagnostic split, whose report states that the
pre-registration and final holdout remain untouched. Every measurement here is
exploratory, and the LP diagnostics use only air05, trento1 and neos8 from the
frozen calibration partition under a 20000-iteration cap.
% source: results/*/data.json (32 files: 23 with the field, all false; 9 without it)
% source: results/hardware_codex/REPORT.md lines 5 and 7

Temporal fusion is untested on the sparse families rather than proved
inapplicable to them. The kernel is a dense two-dimensional stencil with tile
edge $S = B + 2T$ and a halo of width $T$ whose valid region shrinks one cell
per side per iteration. Transfer to a sparse incidence operator was not
attempted and no fusion measurement on netflow, OT or MCF exists, so this is a
scope limit and not an impossibility result.
% source: results/kernel_fusion_20260909/tv_temporal.cu lines 2-3, 18 and 26

The $\omega_0$ exponent is censored by the ratio grid, and the estimate has
moved as that censoring has fallen. Two conventions are fixed for every count in
this paragraph and named again in the caption of the exponent table in
Section~5: an instance counts as censored when any ratio attaining its minimum
certified $N$ sits at an endpoint of that instance's grid, and tied argmin
ratios are averaged geometrically before the log-log fit. The narrow grid has
nine points from 0.01 to 100 and places the oracle at the floor for 91 of 209
instances and at the ceiling for 11; its log-log fit of the oracle ratio on
$\omega_0$ gives slope $-0.94$, or $-1.14$ over the 118 instances left after the
floor-censored ones are dropped. Four wider or finer grids now exist, and the
slope magnitude rises across them: $-0.94$ with 102 of 209 on an edge (9
points), $-1.54$ with 33 of 108 (13 points), $-1.58$ with 46 of 144 (12 points,
8 seeds), $-1.68$ with 15 of 80 (19 points) and $-1.74$ with 14 of 88 (11 points
over ten decades). It does not rise monotonically in the censored fraction: the
eight-seed grid censors a larger fraction than the 13-point grid and still
returns the larger slope. The regret-minimising exponent moves in the other
direction and far less: $2.10$ on the nine-point pooled grid, $2.00$ on the
13-point wide grid, $1.90$ on the 12-point eight-seed grid, $2.00$ on the
19-point fine grid and $2.00$ on the 11-point ultrawide grid. It is not
identified. At the best constant for each exponent, every exponent from $1.85$
to $2.15$ is within 1 percent of the optimum on the wide grid, $1.65$ to $2.45$
on the eight-seed grid, $1.75$ to $2.05$ on the fine grid and $1.80$ to $2.25$
on the ultrawide grid; on the narrow pooled grid that set is not an interval at
all, running from $1.40$ to $2.90$ with an isolated point at $1.15$. The two
estimators bracket the exponent between about $1.7$ and $2.1$ and approach each
other from opposite directions as censoring falls, but they have not met. Both
remain consistent with the value 2 implied by PDLP's primal weight, and the
bracket lies at its lower end below it. The 19-point grid has run and is
included above; it does not narrow the bracket. No claim here should be read as
pinning the exponent tighter than that bracket, and none of the five grids
returns $1.8$ as the argmin of either estimator.
% source: results/finish_20260910/omega_{4090,bw}.json (9 points, 0.01 to 100, 1944 rows,
%   216 instances, 209 with a converged run)
% source: results/finish_20260910/omega_wide_bw.json (13 points, 1e-4 to 100, 1404 rows)
% source: results/finish_20260910/omega_holdout.json (13 points, 702 rows)
% source: results/finish_20260910/omega_ot_wide.json (13 points, 702 rows)
% source: results/finish_20260910/omega_mcf.json (12 points, 1e-4 to 1000, 648 rows)
% source: results/finish_20260910/omega_ultrawide_bw.json (11 points, 1e-6 to 1e4, 968 rows)
% source: results/finish_20260910/omega_fine_4090.json (19 points, 1e-5 to 300, 1520 rows)
% Slopes recomputed with tied argmin ratios averaged geometrically, the
% convention used in every section: -0.9365 / -1.5412 / -1.5781 / -1.6834 /
% -1.7380. Ties broken to the smallest ratio instead gives -0.9309 / -1.5307 /
% -1.5716 / -1.6762 / -1.7324, and to the largest -0.9421 / -1.5516 / -1.5845 /
% -1.6905 / -1.7437; the choice moves each slope by at most 0.011.
% Excluding the floor-censored instances on the narrow grid, -1.1368 over 118.
% Censoring counts under the convention above: 102 of 209 (91 floor, 11
% ceiling), 33 of 108, 46 of 144, 15 of 80, 14 of 88.
% Regret exponents at C on a 0.01 dex grid and p on a 0.05 grid: 2.10, 2.00,
% 1.90, 2.00, 2.00, at regret 1.1805, 1.1858, 1.1571, 1.1685, 1.0618.

A clairvoyant policy that knows the cell and the scaling still carries
$1.112\times$ geometric-mean and $1.814\times$ worst-case regret against the
per-seed oracle over 142 of 144 instances. That measurement covers 2 cells and
8 seeds, while the transfer results of Table~\ref{tab:omega} cover 10 cells,
and 54 of the seed sweep's 144 instances also appear in the netflow wide row of
that table, so the two are not computed on disjoint instance sets. It is an
indication of the seed noise present in this kind of measurement and not a
floor for any regret figure in that table.
% source: results/finish_20260910/omega_seeds_4090.json (2 cells, 8 seeds,
% 9 alphas, 144 instances, 142 with a converged run; recomputed regret 1.1123
% geometric mean and 1.8140 worst). The 10 cells are the distinct cells of
% omega_wide_bw.json (4), omega_ot_wide.json (2), omega_mcf.json (2) and
% omega_holdout.json (2). Overlap on (cell, seed, alpha) with the seed sweep is
% 54 instances, all inside netflow wide. No results JSON carries a device
% field, so no device count is claimed here.

No LLM advantage is claimed: the arms shared a fallback RNG stream and were
positively correlated rather than independent, with 336 of 336 fallback-sourced
LLM-arm candidates byte-identical to candidates of the matched random run.
% source: results/sprint/phase2/analysis/gate1_numbers.json, budgets.arm_coupling.per_seed
% (74 + 99 + 48 + 68 + 47 = 336, all byte_identical_in_matched_random_run)
No joint-search advantage is claimed. On LP, backend-only, sequential,
reverse-sequential and joint selection all choose \texttt{frozen\_B0} /
\texttt{f64\_K50} at $0.392627$\,s over two instances. On TV, joint,
implementation-only and implementation-then-algorithm all choose
\texttt{condat\_vu\_banked} / \texttt{schedule\_K256} at $0.0076265$\,s
penalised SGM-1, while algorithm-then-implementation chooses
\texttt{pdhg\_banked} / \texttt{schedule\_K256} at $0.0077793$\,s, a
$1.020\times$ gap. Sequential selection therefore ties joint in one order and
not the other, over an exhaustively enumerated 12-cell catalog whose own
protocol record declares it development only and not search efficiency.
% source: results/hardware_codex/REPORT.md, joint_lp_4090_v2 table
% source: results/method_axis_20260909/factorial.json selections and rows
% source: results/method_axis_20260909/protocol.json status field

Five narrower boundaries belong on the record. Scheme identity is inert at
$K=64$ and $K=256$ at $10^{-4}$ but not everywhere: at $K=16$ and $n{=}256$, on
both devices, \texttt{condat\_vu\_banked} averages $584.000$ iterations against
$578.667$ for \texttt{pdhg\_banked}, with per-instance vectors
$[1008, 512, 608, 416, 320, 640]$ against $[1008, 496, 608, 400, 320, 640]$. At
$n{=}128$ and $n{=}496$ on the Blackwell the two schemes agree at $K=16$, at
$552.000$ and $1040.000$, so the split is size-specific as well as
cadence-specific.
% source: results/bw_20260909/sz{128,256,496}_bw/factorial.json
% source: results/tv_ffi_20260906/cadence_4090_v1/factorial.json
Rejected recipes were never run, so the gate is verified as sufficient and not
as necessary. The dense gate grid records the decision alone: 1573 gate
evaluations, 121 distinct $(\tau, \sigma)$ pairs across 13 cells for scheme
\texttt{pdhg} only, admitting 1220 and rejecting 353, with $N$,
\texttt{exec\_s} and \texttt{converged} null in every row because no solve was
attached to it. Sufficiency also carries a budget scope: across the two solving
sweeps on that device, 1755 of 2106 runs were admitted and 257 of those hit the
150000-iteration cap, so admission does not imply convergence within budget.
% source: results/structure_20260909/rawgrid_4090gate.json
%   (1573 runs, args.no_solve true, gate_admit only)
% source: results/structure_20260909/structure_summary_4090.json totals block,
%   which sums grid_4090.json (648 runs, 540 admitted, 18 capped) and
%   regular_4090.json (1458, 1215, 239)
Sample sizes are small: six OCT images per size in the end-to-end sweeps, four
montages in each of the larger sizes, three timing repeats per instance, and
three seeds per cell in the four sweeps of Table~\ref{tab:omega} and the two
that make up the narrow grid, rising to four seeds in the ultrawide and fine
grids and eight in the two seed-noise grids. All fourteen end-to-end fusion
sweeps ran at $K=256$, so the iteration agreement they establish is resolved
only to the nearest 256 against measured $N$ between 256 and 2048.
% source: args blocks of results/kernel_fusion_20260909/{e2e,e2e_bw,large,bw_day2}/*.json
%   (args.n = 6 in the eight full_* files, args.n = 4 in the six montage files,
%    args.repeats = 3 throughout, args.K = 256 throughout; iters over all 72
%    instances x 8 arms take the values 256, 512, 768, 1024, 1280, 1536, 2048)
% source: args blocks of results/finish_20260910/omega_*.json: args.seeds = 3 in
%   omega_{4090,bw,wide_bw,ot_wide,mcf,holdout}.json, 4 in
%   omega_{ultrawide_bw,fine_4090}.json, 8 in omega_{seeds8,seeds_4090}.json
Neither device was exclusive during timing. On the 4090 a resident server held
15427\,MiB of 24564\,MiB at 0 percent utilization; the Blackwell box is shared
and these runs used a GPU that was idle but not reserved. No occupancy snapshot
was recorded for the Blackwell, so that second claim rests on the run report
rather than on a machine-readable field. The 4090 reference cell is stable to
0.66 percent across 12 independent processes, spanning $20.297$ to $20.432$
microseconds per iteration, a factor of $1.00665$.
% source: results/method_axis_20260909/protocol.json gpu_snapshot
%   ("NVIDIA GeForce RTX 4090, 15427 MiB, 24564 MiB, 0 %"); no protocol.json
%   under results/ carries a Blackwell gpu_snapshot
% source: results/kernel_fusion_20260909/REPORT.md lines 82-84
% source: results/tv_ffi_20260906/DECOMPOSITION_REPORT.md lines 5-8
Finally, cadence is not our discovery as an engineering choice. OR-Tools ships
\texttt{termination\_check\_frequency} with a default of 64 and an explicit
annotation that the check involves extra work; what is new here is its status as
the one gene the two factors share, not the fact that it is tuned.
% source: survey record at 2-related-work.tex lines 89-90, citing
% ORToolsPDLP2026. This is an external fact about OR-Tools and is not measured
% under results/; DECOMPOSITION_REPORT.md names neither OR-Tools nor
% termination_check_frequency and is not a source for it.

\subsection{What would falsify the thesis}

Table~\ref{tab:falsify} lists the five tests we consider decisive. The two
negatives are the sharpest, because a better instrument would overturn them
tomorrow. Roofline is anti-correlated at $n{=}128$ ($-0.943$) and positively
correlated at $n \geq 992$ ($+0.829$), with per-size Spearman coefficients of
$-0.943$, $-0.143$, $+0.143$, $+0.943$ and then $+0.829$ at each of the three
largest sizes, for a mean of $0.355$. Sharpness pools to a log-log
$r = +0.943$ over 17 converged instances, though the same pooled correlation
unlogged is only $+0.724$, and the sign is not stable within families:
$-0.539$ within netflow ($n=13$), $-0.842$ within the regular topology
($n=6$), $-0.120$ within grid ($n=7$) and $+0.983$ within TV ($n=4$). One of
the four subgroups therefore does not invert. The four are not disjoint: the
netflow subgroup is exactly the union of the grid and regular topology
subgroups, and the subgroups of size 4 and 6 carry no inference on their own.
% source: results/diagnostics_20260909/roofline_4090.json, pred_s_per_pixel_iter per
%   tile, rank-correlated against summed arm times in the seven 4090 timing files
%   (e2e/full_4090_sz{128,256,384,496}.json, large/large_4090_m{2,3,4}.json).
%   No Spearman value is stored in any JSON; it is recomputed from those files
%   by scripts/diagnostics/score_roofline.py, whose output is archived at
%   results/artifacts_20260910/roofline_scores_4090.log and tabulated in
%   results/diagnostics_20260909/REPORT.md lines 26-29, the source for 0.567.
% source: results/diagnostics_20260909/sharpness_4090.json, 17 of 18 rows
%   converged (grid12d4 instance 1 hit the 20000 cap). Families: grid12d4,
%   grid16d4, grid24d4, grid32d4 give the 7 grid rows; regular256d3,
%   regular256d4, regular512d4 give the 6 regular rows; netflow is those 13
%   together; tv128 and tv256 give the 4 TV rows. Recomputed log-log Pearson r
%   on (tau, N): pooled +0.943, unlogged +0.724, netflow -0.539,
%   regular -0.842, grid -0.120, TV +0.983.

\begin{table}[t]
\centering
\caption{Falsifiers, with the measurement each would overturn.}
\label{tab:falsify}
\begin{tabular}{p{0.46\linewidth}p{0.46\linewidth}}
\toprule
observation that falsifies & current measurement \\
\midrule
A fusion tile or execution structure that changes the iteration count on any
instance
& 72 of 72 instances agree exactly at $K=256$, zero mismatches; a nine-repeat
rerun reproduces the same means, 640 at $n{=}256$ and 1024 at $n{=}496$, while
the timing ratios move 4 to 5 percent \\
\addlinespace
An adaptive-schedule leak that grows with problem size or family
& 16.6 percent against 3.8 percent under fixed precision, at one cadence, one
size, one family and one tolerance; the second size and second family were not
measured, so this falsifier is not yet evaluable \\
\addlinespace
A cheap a-priori predictor that picks the fusion tile
& the a-priori roofline model reaches 0.355 mean Spearman and 1 correct pick of
7; adding a per-launch cost fitted on the same sweep raises it to 0.567 and
still picks 1 of 7 \\
\addlinespace
A family on which the $r = C/\omega_0^{2}$ rule loses to a fixed ratio
& on three held-out families it does not: $1.131\times$ against $2.168\times$
pooled geomean regret over 162 instances, certifying 162 of 162 against 146 \\
\addlinespace
A device or family on which cadence is not the dominant variance share
& 47.43 and 51.19 percent of variance in cost per iteration, from a 12-cell
fANOVA on one device, one family and one size, with one observation per cell and
no residual term, and with cadence carrying more levels than either competitor \\
\bottomrule
\end{tabular}
% source: results/kernel_fusion_20260909/{e2e,e2e_bw,large,bw_day2}/*.json;
% source: results/finish_20260910/fusion_r9_sz{256,496}.log (best arm 1.940x against
%   the three-repeat 1.867x at n=256, 1.428x against 1.356x at n=496)
% source: results/method_axis_20260909/factorial.json
% source: results/diagnostics_20260909/roofline_4090.json and the seven 4090 timing files
% source: results/finish_20260910/omega_{wide_bw,ot_wide,mcf,holdout}.json
% source: results/artifacts_20260910/fanova.json
\end{table}

\section{Conclusion}

Time to a certified solution factorises into an iteration count set by the
method, the certificate cadence and the tolerance, and a cost per iteration set
by the hardware, the execution structure, the kernel and the size. At fixed size
and fixed cadence the first factor is invariant to the hardware, the execution
structure and the kernel, exactly, in 72 of 72 measured instances across two
architectures and seven sizes. It is not invariant to size: at $K=64$ in
float64 with \texttt{condat\_vu\_banked}, $N$ is $576.000$ at $n{=}128$,
$597.333$ at $n{=}256$, $736.000$ at $n{=}384$ and $1066.667$ at $n{=}496$.

Cadence is the gene the two factors share within the recipe, and it carries the
largest single variance share in cost per iteration. It is not the only gene in
both. Size appears in both factors as well, and precision changes $N$ at some
sizes: on the Blackwell at $K=64$, \texttt{f32\_cert64} gives $586.667$ against
float64's $576.000$ at $n{=}128$ and $1077.333$ against $1066.667$ at $n{=}496$,
while at $n{=}256$ the two agree at $597.333$.
% source: results/tv_size_20260909/sz{128,256,384,496}_4090/factorial.json
% source: results/bw_20260909/sz{128,256,496}_bw/factorial.json

Three consequences follow. Kernel search runs against a fixed arithmetic budget:
on the two architectures measured, the iteration count was identical at every
size run on both, so $N$ can be held fixed within that scope while the second
factor is timed. Legality is decided symbolically before any solve, by a gate
that is an exact function of $\theta$ and family dependent in its threshold,
verified as sufficient for admission and never as necessary; it admits on TV
what it rejects on netflow, where the PDHG gate product
$\tau \sigma \lVert L \rVert^2$ runs $1.02267$, $1.03040$, $1.03577$ and
$1.03781$ at sizes 12, 16, 24 and 32. Every rejected row is a balanced
configuration, 24 of 24, and no banked configuration is ever rejected. And the
recipe is predictable from measured structure where a fixed choice is not:
scaling the step ratio as $C / \omega_0^{2}$ with $C = 0.977$ fitted on netflow
alone cuts pooled geomean regret on 162 held-out instances from $2.168\times$ to
$1.131\times$ and lifts coverage from 146 to 162 of 162, at an exponent
consistent with, though not pinned to, the value 2 implied by PDLP's primal
weight.
% source: results/structure_20260909/structure_summary_4090.json
%   (gate_is_function_of_theta_only: 351 of 351 admitted at theta 0.1, 0.3, 0.6,
%    0.9 and 0.99, and 0 of 351 at 1.02)
% source: results/netflow_20260909/grid_4090.json (48 rows, 24 rejected, all balanced;
%   for the Condat-Vu rows the recorded value is the LHS of
%   t*L_f/2 + t*s*||L||^2 < 1, not the product)
% source: results/finish_20260910/omega_{wide_bw,ot_wide,mcf,holdout}.json
%   (fit on omega_wide_bw only, C = 0.977, exponent 2)

The numerical admission run for the temporal fusion kernel is archived with the
other artifacts: 106 lines holding 102 checked cases, 17 tile and fused-depth
variants at three image shapes including a non-square 100 by 130, in both
proximal settings. The worst absolute deviation against sequential single-step
reference launches is $4.649 \times 10^{-15}$ and the verdict is PASS within
float64 tolerance. No case is bitwise identical, so the admission is numerical
and not a floating-point equivalence proof.
% source: results/artifacts_20260910/kernel_admission_4090.log
%   (102 rows beginning tv_tf; variants tv_tf{8x4,8x6,8x8}t64, tv_tf{16x2,16x3,
%    16x4,16x6,16x8}t128, tv_tf16x4t256, tv_tf{24x3,24x4,24x6,32x2,32x3,32x4,
%    32x6,40x2}t256; sizes 64x64, 128x128, 100x130; bitwise = no on all 102)


\section*{Reproducibility Statement}
Every numeric claim in this paper names the artifact that produced it in a
LaTeX source comment of the form \texttt{\% source: results/...}. Iteration
counts additionally name their instance population, because scan selection
depends on both the sampling seed and the number of scans drawn: the catalog
fold is six scans at data seed 20260907 drawn from a pool of twelve, and the
fusion sweep is six scans at seed 0, and the two selections share no images.
Running the fusion harness on the catalog draw reproduces the catalog
iteration counts. The frozen final holdout has not been evaluated and no
number in this paper comes from it.

\section*{Ethics Statement}
The imaging experiments solve a convex denoising objective on an existing OCT
dataset. Optimization accuracy does not establish preservation of pathology or
suitability for clinical decisions. We report no clinical efficacy result.

\section*{AI Use Statement}
Language models assisted literature discovery, experiment planning, code
implementation, debugging, manuscript drafting and interpretation of results.
The authors are responsible for verifying the mathematics, citations, code and
reported measurements.

\bibliographystyle{iclr2027/iclr2027_conference}
\bibliography{refs}

\appendix
% NOTE (2026-08-30): theory appendix created in the audit-driven restructure.
% Contains the proofs for Section 4 (sec:theory). Style rules: never use the
% em-dash; upright (non-italic) body for theorem-like environments.
\section{Proofs for Section~\ref{sec:theory}}
\label{app:theory}

Throughout this appendix the standing assumptions of the setup paragraph of
Section~\ref{sec:theory} are in force. Because every metric $H\succ0$ makes
$\langle u,Hv\rangle$ an inner product on a finite-dimensional space, all
statements are statements about a Hilbert space, and we work directly in
$\|\cdot\|_H$ without rescaling.

\subsection{Convergence and rate for certified iterations}

\begin{proposition}[Krasnosel'ski\u{\i}--Mann convergence with a
last-iterate rate]
\label{prop:km}
Let $H\succ0$, let $T$ be $\theta$-averaged in $\|\cdot\|_H$ with
$\theta\in(0,1)$, and let $\operatorname{Fix}T\neq\emptyset$. Then the
iteration $z_{k+1}=Tz_k$ converges, from any $z_0$, to some
$z^\star\in\operatorname{Fix}T$; the residual $\|z_{k+1}-z_k\|_H$ is
nonincreasing in $k$; and, for every $k\geq0$,
\begin{equation}
 \|z_{k+1}-z_k\|_H^2
 \;\leq\;
 \frac{\theta}{(1-\theta)\,(k+1)}\;
 \operatorname{dist}_H^2\!\left(z_0,\operatorname{Fix}T\right).
 \label{eq:km-rate}
\end{equation}
\end{proposition}

\begin{proof}
Write $T=(1-\theta)I+\theta S$ with $S$ nonexpansive in $\|\cdot\|_H$; all
norms in this proof are $\|\cdot\|_H$. Fix any
$z^\star\in\operatorname{Fix}T$. Expanding $z_{k+1}=(1-\theta)z_k+\theta
Sz_k$ with the identity
$\|(1-\theta)a+\theta b\|^2=(1-\theta)\|a\|^2+\theta\|b\|^2
-\theta(1-\theta)\|a-b\|^2$, at $a=z_k-z^\star$ and $b=Sz_k-z^\star$, and
bounding $\|Sz_k-z^\star\|\leq\|z_k-z^\star\|$ by nonexpansiveness gives the
Fej\'er inequality
\begin{equation}
 \|z_{k+1}-z^\star\|^2
 \;\leq\;
 \|z_k-z^\star\|^2-\frac{1-\theta}{\theta}\,\|z_{k+1}-z_k\|^2 .
 \label{eq:fejer}
\end{equation}
The map $T$ is nonexpansive, so
$\|z_{k+2}-z_{k+1}\|\leq\|z_{k+1}-z_k\|$. Summing \eqref{eq:fejer} over the
first $k+1$ iterations, keeping only the last (smallest) residual, and
minimizing over $z^\star$ yields \eqref{eq:km-rate}. By \eqref{eq:fejer} the
iterates are bounded and $\|z_{k+1}-z_k\|\to0$, so by continuity every
cluster point of $(z_k)$ is fixed by $T$; reusing \eqref{eq:fejer} with such
a cluster point as $z^\star$ shows $\|z_k-z^\star\|$ is nonincreasing and
vanishes along a subsequence, hence $z_k\to z^\star$. This argument is the
Krasnosel'ski\u{\i}--Mann theorem with its modern rate
\citep[Theorem~1]{RyuYin2022LargeScale}; see also
\citet{Krasnoselskii1955Two}, \citet{Mann1953Mean}, and
\citet{BauschkeCombettes2017Convex}.
\end{proof}

\begin{remark}[Monitored residual, normalization, and summability]
\label{rem:monitored-residual}
The theorem controls $\|z_{k+1}-z_k\|_{H_G}$, the step residual of the
certified map that is actually iterated. For a genome
$\mathrm{relax}_\rho(G')$ the compiled map is
$T_G=(1-\rho)I+\rho T_{G'}$, so
$z_{k+1}-z_k=\rho\,(T_{G'}z_k-z_k)$: the two residuals differ by the factor
$\rho$, which matters when comparing genomes with different relaxation genes,
and \eqref{eq:km-rate} divided by $\rho^2$ bounds the unrelaxed fixed-point
residual $\|T_{G'}z_k-z_k\|_{H_G}^2$. Summing \eqref{eq:fejer} also gives the
summability estimate
$\sum_{k=0}^{\infty}\|z_{k+1}-z_k\|_{H_G}^2\leq
\frac{\theta_G}{1-\theta_G}\operatorname{dist}_{H_G}^2
(z_0,\operatorname{Fix}T_G)$.
\end{remark}

\subsection{Base-scheme certificates}

\begin{lemma}[Base schemes]
\label{lem:base-certificates}
For each base scheme of Table~\ref{tab:certificates} with positive stepsizes
satisfying the gate inequality, the pair $(H_G,\theta_G)$ of the table, with
$\theta_G=1/(2-c_G)$, together with the scheme's standard decoder $\pi_G$,
is a construction certificate in the sense of
Definition~\ref{def:certificate}.
\end{lemma}

\begin{proof}
We first establish every row with the true norm $\|L\|$ in place of
$\ell$, then account for the conservative bound.

Since $\nabla f$ is $(1/L_f)$-cocoercive by the Baillon--Haddad theorem
\citep{BaillonHaddad1977Quelques}, the forward step $I-\alpha\nabla f$ is
$(\alpha L_f/2)$-averaged for $\alpha\in(0,2/L_f)$, and the proximal step is
$\tfrac12$-averaged. A composition of $\theta_1$- and $\theta_2$-averaged
maps is $(\theta_1+\theta_2-2\theta_1\theta_2)/(1-\theta_1\theta_2)$-averaged
\citep[Theorem~27]{RyuYin2022LargeScale}, a constant due to
\citet{OguraYamada2002Nonstrictly}; see also
\citet{CombettesYamada2015Compositions}. At $\theta_1=\tfrac12$ and
$\theta_2=c$ it is $1/(2-c)$; this is the FBS row, and the Davis--Yin row is
the same constant \citep[Theorem~28]{RyuYin2022LargeScale}, established by
\citet{DavisYin2017Three}. DRS is
$\tfrac12\bigl(I+R_{\alpha A}R_{\alpha B}\bigr)$, a composition of reflected
resolvents, hence firmly nonexpansive for every $\alpha>0$
\citep[\S2.7]{RyuYin2022LargeScale}, a fact going back to
\citet{LionsMercier1979Splitting}. For the primal--dual rows,
$M_{\tau\sigma}\succ0$ exactly when $\tau\sigma\|L\|^2<1$, and in that metric
PDHG \citep{ChambollePock2011First} is a proximal-point iteration on the
saddle subdifferential, hence firmly nonexpansive in
$\|\cdot\|_{M_{\tau\sigma}}$ \citep[\S3.3]{RyuYin2022LargeScale}; see also
the contraction analysis of \citet{HeYuan2012Convergence}. Condat--V\~u
\citep{Condat2013Primal,Vu2013Splitting} is variable-metric
forward--backward in the same metric \citep{CombettesVu2014Variable}, and its
forward step is $c$-averaged in $\|\cdot\|_{M_{\tau\sigma}}$ with
$c=\tau L_f/\bigl(2(1-\tau\sigma\|L\|^2)\bigr)$, so composition gives
$1/(2-c)$ again \citep[\S3.3]{RyuYin2022LargeScale}. The inequality
$\tau L_f/2+\tau\sigma\|L\|^2<1$ is equivalent, for $\tau,\sigma>0$, to the
conjunction of $\tau\sigma\|L\|^2<1$ and $c<1$, so a single inequality
delivers both (C1) and a positive-definite metric.

Now let $\ell\geq\|L\|$ be the supplied bound and let $c_G$ denote the
tabulated forward-step constant computed from $\ell$. The checked
inequalities in $\ell$ imply the corresponding inequalities in $\|L\|$, so
$H_G\succ0$, and the true forward-step constant $c$ computed from $\|L\|$
satisfies $c\leq c_G<1$, so the compiled iteration is $\theta$-averaged with
$\theta=1/(2-c)$. A $\theta$-averaged map is also $\theta'$-averaged for
every $\theta'\in[\theta,1)$: writing $T=(1-\theta)I+\theta S$ gives
$T=(1-\theta')I+\theta'S'$ with $S'=(1-\theta/\theta')I+(\theta/\theta')S$,
a convex combination of nonexpansive maps. Since
$\theta=1/(2-c)\leq1/(2-c_G)=\theta_G<1$, the pair $(H_G,\theta_G)$
certifies (C1), exactly when $\ell=\|L\|$ and conservatively otherwise.

For (C2) and (C3), no bijection between fixed points and solutions is
needed; it suffices that fixed points exist and that every fixed point
decodes to a solution. For the primal--dual schemes the state is the saddle
pair and, because $M_{\tau\sigma}\succ0$, the fixed points of the
(variable-metric) proximal-point and forward--backward iterations are
exactly the zeros of the saddle operator, i.e.\ the saddle set, which is
nonempty by assumption; $\pi_G$ is the primal coordinate
\citep[\S3.3]{RyuYin2022LargeScale}. For FBS and Davis--Yin the fixed points
are exactly the solutions of the associated inclusion (for Davis--Yin after
applying the decoder $\pi_G=\operatorname{prox}_{\alpha g}$), and for DRS
every fixed point $z$ satisfies
$\operatorname{prox}_{\alpha g}(z)\in\mathcal X^\star$ while every solution
arises from at least one fixed point; these standard correspondences are
\citep[\S2.7]{RyuYin2022LargeScale}. In each case the nonempty saddle set
yields a fixed point through the correspondence, giving (C2), and the
decoder maps all of $\operatorname{Fix}T_G$ into $\mathcal X^\star$, giving
(C3).
\end{proof}

\iclrTODO{Base-scheme instantiation table: for each row of
Table~\ref{tab:certificates}, give the compiled state space, the exact
update the compiler emits for problem~\eqref{eq:composite} (including how
FBS and DRS instantiate the three-term objective), the decoder, and the
fixed-point correspondence, and verify the entries against the
implementation. Without this table, Lemma~\ref{lem:base-certificates}
certifies the textbook operators, not yet the emitted code.}

\subsection{Closure rules}

\begin{lemma}[Relaxation]
\label{lem:relaxation}
Let $T$ be $\theta$-averaged in $\|\cdot\|_H$ and let $\rho>0$ satisfy
$\rho\theta<1$. Then $(1-\rho)I+\rho T$ is $\rho\theta$-averaged in
$\|\cdot\|_H$ with $\operatorname{Fix}\bigl((1-\rho)I+\rho
T\bigr)=\operatorname{Fix}T$.
\end{lemma}

\begin{proof}
Writing $T=(1-\theta)I+\theta S$ with $S$ nonexpansive gives
$(1-\rho)I+\rho T=(1-\rho\theta)I+\rho\theta S$, which is
$\rho\theta$-averaged because $\rho\theta\in(0,1)$, and $\rho>0$ makes the
fixed-point equations equivalent. The decoder and metric are untouched, so
the certificate transports as stated in Theorem~\ref{thm:soundness}(ii).
\end{proof}

% PROPOSED to Nikhil: the operational account of blending (mechanism, the
% semantic fixed-point caveat, and the empirical negative) now lives in
% Appendix app:blending, demoted from the main text; the old iclrTODO asking
% for one consistent account is resolved by that demotion.
The next lemma is the certificate rule for the blend extension discussed in
Appendix~\ref{app:blending}. It is excluded from
Theorem~\ref{thm:soundness} because its hypothesis of a nonempty common
fixed-point set is semantic: the current gate has no syntactic test for it.

\begin{lemma}[Blends]
\label{lem:blend}
Let $T_1,T_2$ be $\theta_1$- and $\theta_2$-averaged in a common
$\|\cdot\|_H$ with
$\operatorname{Fix}T_1\cap\operatorname{Fix}T_2\neq\emptyset$, and let
$\lambda\in(0,1)$. Then $T=\lambda T_1+(1-\lambda)T_2$ is
$\bigl(\lambda\theta_1+(1-\lambda)\theta_2\bigr)$-averaged in $\|\cdot\|_H$
and $\operatorname{Fix}T=\operatorname{Fix}T_1\cap\operatorname{Fix}T_2$.
\end{lemma}

\begin{proof}
Write $T_i=(1-\theta_i)I+\theta_iS_i$ and set
$\theta=\lambda\theta_1+(1-\lambda)\theta_2$. Then
$T=(1-\theta)I+\theta\bigl(\tfrac{\lambda\theta_1}{\theta}S_1
+\tfrac{(1-\lambda)\theta_2}{\theta}S_2\bigr)$, and the bracket is a convex
combination of nonexpansive maps, hence nonexpansive; this is the
combination rule of \citet{CombettesYamada2015Compositions}. The inclusion
$\operatorname{Fix}T\supseteq\operatorname{Fix}T_1\cap\operatorname{Fix}T_2$
is immediate. Conversely, let $z\in\operatorname{Fix}T$ and
$w\in\operatorname{Fix}T_1\cap\operatorname{Fix}T_2$. Then
$\|z-w\|\leq\lambda\|T_1z-w\|+(1-\lambda)\|T_2z-w\|\leq\|z-w\|$, so both
inequalities are equalities and $\|T_iz-w\|=\|z-w\|$ for $i=1,2$; the
Fej\'er inequality \eqref{eq:fejer} applied to $T_i$ at $z$ then forces
$\|T_iz-z\|=0$. Hence $z\in\operatorname{Fix}T_1\cap\operatorname{Fix}T_2$.
\end{proof}

\subsection{The Halpern wrapper}

\begin{proposition}[Halpern genomes]
\label{prop:halpern}
Let $T_G$ carry a certificate $(H_G,\theta_G,\pi_G)$ and set
$S_G=I+\theta_G^{-1}(T_G-I)$. Then $S_G$ is nonexpansive in
$\|\cdot\|_{H_G}$ with $\operatorname{Fix}S_G=\operatorname{Fix}T_G$, and the
compiled iteration
\begin{equation}
 z_{k+1}=\lambda_{k+1}z_{\rm anc}+(1-\lambda_{k+1})S_Gz_k ,
 \label{eq:halpern}
\end{equation}
under any admitted schedule, i.e.\ $\lambda_k\to0$,
$\sum_k\lambda_k=\infty$, and $\sum_k|\lambda_{k+1}-\lambda_k|<\infty$,
converges to $P^{H_G}_{\operatorname{Fix}T_G}(z_{\rm anc})$, the projection
of the anchor onto $\operatorname{Fix}T_G$ in $\|\cdot\|_{H_G}$
\citep{Halpern1967Fixed,Wittmann1992Approximation}. The canonical schedule
$\lambda_{k+1}=1/(k+2)$ meets all three conditions and, when the wrapper
anchors at the initial iterate ($z_{\rm anc}=z_0$), in addition gives
\begin{equation}
 \|S_Gz_k-z_k\|_{H_G}^2
 \;\leq\;
 \frac{4\operatorname{dist}_{H_G}^2\!\left(z_0,\operatorname{Fix}
 T_G\right)}{(k+1)^2}
 \label{eq:halpern-rate}
\end{equation}
\citep{Lieder2021Optimized,Kim2021Accelerated}.
\end{proposition}

\begin{proof}
$S_G$ is precisely the nonexpansive map in the averaged decomposition
$T_G=(1-\theta_G)I+\theta_GS_G$ of (C1), and $T_Gz=z$ if and only if
$S_Gz=z$ because $\theta_G>0$. Since $(\mathcal Z_G,\langle\cdot,H_G
\cdot\rangle)$ is a Hilbert space, the cited Halpern results apply verbatim
in $\|\cdot\|_{H_G}$: convergence to the projection of the anchor is
\citet{Wittmann1992Approximation}, and the rate \eqref{eq:halpern-rate} is
\citet{Lieder2021Optimized}.
\end{proof}

Rate \eqref{eq:halpern-rate} is an $O(1/k^2)$ bound on the same squared
residual that \eqref{eq:km-rate} bounds at $O(1/k)$, and it is exactly
optimal in the following sense: \citet{ParkRyu2022Exact} exhibit a matching
lower bound for deterministic fixed-point methods that access $S_G$ only
through evaluations, with iterates constrained to the span of the anchor and
past evaluations. The claim is optimality within that oracle class, not an
unrestricted impossibility statement.

\subsection{Finite-prefix switching}
% PROPOSED to Nikhil: replaces summable-error safeguard with finite-prefix
% switch, matches implementation

\begin{lemma}[Finite-prefix switching]
\label{lem:switchk}
Let $T_B$ be $\theta$-averaged in $\|\cdot\|_H$ on a state space
$\mathcal Z$ with $\operatorname{Fix}T_B\neq\emptyset$, let
$T_A:\mathcal Z\to\mathcal Z$ be an arbitrary map, and fix
$K\in\mathbb N$. Define $z_{k+1}=T_Az_k$ for $k<K$ and $z_{k+1}=T_Bz_k$
for $k\geq K$. If $z_K$ is finite, then $(z_k)$ converges to some
$z^\star\in\operatorname{Fix}T_B$, and for every $k\geq K$ the residual
bound \eqref{eq:km-rate} holds with $z_0$ replaced by $z_K$ and $k+1$
replaced by $k-K+1$.
\end{lemma}

\begin{proof}
For $k\geq K$ the sequence is exactly the certified iteration
$z_{k+1}=T_Bz_k$ started at the finite point $z_K$.
Proposition~\ref{prop:km} applies to this orbit and gives convergence to a
point of $\operatorname{Fix}T_B$ together with the rate from the new
starting point. The prefix enters only through the constant
$\operatorname{dist}_H(z_K,\operatorname{Fix}T_B)$, so it cannot affect the
asymptotics.
\end{proof}

The compiler enforces the hypotheses: $K$ is a bounded gene, and the
handoff is guarded by a finiteness check on $z_K$. If a prefix produces
a nonfinite state, the implementation starts the tail from the last finite
state, or the initial state if necessary, and records that recovery. This
recovered trajectory is distinct from the unrecovered prefix in the lemma.
The rule deliberately trades any guarantee about the prefix for total
freedom in choosing it: $T_A$ needs no certificate, no error budget, and no
per-step monitor, in contrast to summable-error analyses of inexact
Krasnosel'ski\u{\i}--Mann iterations
\citep{Combettes2001QuasiFejerian,Combettes2004Solving}, inexact
Douglas--Rachford and ADMM \citep{EcksteinBertsekas1992Douglas}, and
safeguarded Anderson acceleration \citep{Zhang2020Globally}, which admit
perturbations at every iteration at the price of a per-step error test.

\FloatBarrier

\section{Averaged blending}
\label{app:blending}

The extended genome can represent
$\mathrm{avg}(\lambda;T_1,T_2)=\lambda T_1+(1-\lambda)T_2$.
For averaged maps in a common metric with a nonempty common fixed-point set,
Lemma~\ref{lem:blend} gives the blended certificate
\citep{RyuYin2022LargeScale,CombettesYamada2015Compositions}.
The current gate admits restricted cases: compatible FBS parents can use
different steps, DRS/Davis--Yin parents must preserve the auxiliary fixed-point
encoding through identical steps, and saddle parents require identical
$(\tau,\sigma)$ so their metric and saddle encoding agree. The common
fixed-point condition remains a semantic mathematical obligation; the checker
is not a general decision procedure for that property.

Blending lies outside the grammar of Theorem~\ref{thm:soundness} and is not
used in the hardware factorial. In the standard LP saddle instantiation,
PDHG and Condat--V\~u with identical parameters are the same map, so blending
those parents adds no mathematical variation. The expression contains two
parent evaluations, but actual execution cost depends on compiler elimination
and fusion and must be measured. We make no empirical discovery claim for this
extension in the current evaluation.

% AUTO-GENERATED by scripts/gen_instantiation_table.py; do not edit by
% hand. Regenerate from the v2 root with
%   JAX_PLATFORMS=cpu python scripts/gen_instantiation_table.py
% Certificate conditions and citations are extracted from atlas.typing_;
% every display equation is anchored to the builder source in
% atlas/schemes/base_schemes.py (generation fails if the code drifts); the
% deviation numbers are re-measured at generation time by layer A of
% tests/test_compiler_equivalence.py.
% Style rules: never use the em-dash; no italics commands.
\section{Instantiation of the certificates in the compiler}
\label{app:instantiation}

This appendix connects Lemma~\ref{lem:base-certificates} to the emitted
code. For each base scheme of Table~\ref{tab:certificates} it records how
the compiler decomposes problem~\eqref{eq:composite}, the compiled state
space $\mathcal Z_G$ and metric $H_G$, the exact update the compiler
emits, the decoder $\pi_G$, the fixed-point correspondence, and the
certificate condition the gate checks (extracted verbatim, up to LaTeX
conversion, from the typing rules in the implementation). The compiler
dispatches on the shape of the composite: FBS and DRS instantiate the
three-term objective either on a primal two-block form or, when $h$ is
composed with $L$, on the dual problem
\begin{equation}
 \min_{u}\; f^*(-L^{\top}u)+h^*(u),
 \label{eq:inst-dual}
\end{equation}
which requires $g=0$ and $f$ to be $\mu$-strongly convex, so that
$F=f^*\circ(-L^{\top})$ is $(\ell^2/\mu)$-smooth and $\nabla f^*$ is
single-valued; the implemented strongly convex atoms have affine
$\nabla f^*$, which the dual updates below use. PDHG and Condat--V\~u
instantiate the saddle form directly, and Davis--Yin requires the
three-block form without $L$. A structurally inapplicable (problem,
scheme) pair is rejected by the gate as a certified construction failure.
Table~\ref{tab:instantiation} summarizes; the updates and correspondences
follow, one block per scheme.

\begin{table}[h]
\centering
\small
\setlength{\tabcolsep}{3pt}
\begin{tabular}{llllll}
\toprule
scheme & decomposition & state $\mathcal Z_G$ & $H_G$ & update &
decoder $\pi_G$ \\
\midrule
FBS & $f$ smooth $+$ one prox block & primal $x$ &
 $I$ & \eqref{eq:inst-fbs-primal} & $x$ \\
FBS & dual \eqref{eq:inst-dual} & dual $u$ &
 $I$ & \eqref{eq:inst-fbs-dual} & $\nabla f^*(-L^{\top}u)$ \\
DRS & $g$ prox $+$ $h$ prox ($f=0$, no $L$) & auxiliary $z$ &
 $I$ & \eqref{eq:inst-drs-primal} & $\operatorname{prox}_{\alpha g}(z)$ \\
DRS & dual \eqref{eq:inst-dual} & auxiliary $z$ &
 $I$ & \eqref{eq:inst-drs-dual} &
 $\nabla f^*(-L^{\top}\operatorname{prox}_{\alpha h^*}(z))$ \\
PDHG & one prox block $\varphi$ $+$ $h(Lx)$ & saddle $(x,u)$ &
 $M_{\tau\sigma}$ & \eqref{eq:inst-pdhg} & $x$ \\
Condat--V\~u & $f$ smooth $+$ $g$ prox $+$ $h(Lx)$ & saddle $(x,u)$ &
 $M_{\tau\sigma}$ & \eqref{eq:inst-cv} & $x$ \\
Davis--Yin & $f$ smooth $+$ $g$ prox $+$ $h$ prox (no $L$) &
 auxiliary $z$ & $I$ & \eqref{eq:inst-dys} &
 $\operatorname{prox}_{\alpha g}(z)$ \\
\bottomrule
\end{tabular}
\caption{Instantiation of the certificates of Table~\ref{tab:certificates}
in the compiler, per base scheme and accepted decomposition of
problem~\eqref{eq:composite}. $M_{\tau\sigma}$ is the primal--dual metric
of Table~\ref{tab:certificates}; on the dual instantiations the FBS/DRS
certificate is evaluated with the dual smoothness constant
$L_{F}=\ell^2/\mu$ in place of $L_f$.}
\label{tab:instantiation}
\end{table}

\paragraph{FBS.}
Primal form ($L$ absent, exactly one nonzero prox block
$p\in\{g,h\}$):
\begin{equation}
 x_{k+1}=\operatorname{prox}_{\alpha p}\!\bigl(x_k-\alpha\nabla
 f(x_k)\bigr).
 \label{eq:inst-fbs-primal}
\end{equation}
Dual form \eqref{eq:inst-dual} (with $\nabla f^*$ affine, so
$\nabla F(u)=\tfrac1\mu LL^{\top}u-L\nabla f^*(0)$):
\begin{equation}
 u_{k+1}=\operatorname{prox}_{\alpha h^*}\!\Bigl(u_k-\alpha\bigl(
 \tfrac{1}{\mu}LL^{\top}u_k-L\nabla f^*(0)\bigr)\Bigr).
 \label{eq:inst-fbs-dual}
\end{equation}
Certificate condition (from the typing rules): $0 < \alpha < 2/L_f$, evaluated
with $L_f$ on the primal form and $L_F=\ell^2/\mu$ on the dual form.
Fixed-point correspondence: on the primal form
$\operatorname{Fix}T=\mathcal X^\star$; on the dual form fixed points are
the dual solutions, and $\pi_G(u)=\nabla f^*(-L^{\top}u)$ recovers the
unique primal solution (strong convexity of $f$). On the TV form the
decoder specializes to $x=y-L^{\top}\operatorname{clip}(u;[-\lambda,
\lambda])$.

\paragraph{DRS.}
Primal two-prox form ($f=0$, $L$ absent):
\begin{equation}
 x_g^k=\operatorname{prox}_{\alpha g}(z_k),\qquad
 x_h^k=\operatorname{prox}_{\alpha h}(2x_g^k-z_k),\qquad
 z_{k+1}=z_k+x_h^k-x_g^k.
 \label{eq:inst-drs-primal}
\end{equation}
Dual form \eqref{eq:inst-dual} with $A=\partial h^*$ and $B=\nabla F$:
$J_{\alpha A}=\operatorname{prox}_{\alpha h^*}$, and because
$\nabla f^*$ is affine the resolvent $J_{\alpha B}$ is the exact linear
solve
\begin{equation}
 u^k=\operatorname{prox}_{\alpha h^*}(z_k),\quad
 w^k=\Bigl(I+\tfrac{\alpha}{\mu}LL^{\top}\Bigr)^{-1}
 \bigl(2u^k-z_k+\alpha L\nabla f^*(0)\bigr),\quad
 z_{k+1}=z_k+w^k-u^k,
 \label{eq:inst-drs-dual}
\end{equation}
computed by conjugate gradients to relative residual $10^{-12}$.
Certificate condition: $0 < \alpha < \infty$.
Fixed-point correspondence: every fixed point $z$ decodes to a solution
through the resolvent ($\operatorname{prox}_{\alpha g}(z)$ on the
primal form, the dual solution $\operatorname{prox}_{\alpha h^*}(z)$
then $\nabla f^*(-L^{\top}\cdot)$ on the dual form), and every solution
arises from at least one fixed point.

\paragraph{PDHG.}
The primal block $\varphi$ is $f$ or $g$ (whichever is nonzero), or, when
$f$ is affine and $g$ is prox-friendly, the exact combined block
$\varphi=f+g$ through the linear-tilt identity
$\operatorname{prox}_{\tau(q+p)}=\operatorname{prox}_{\tau p}
\circ\operatorname{prox}_{\tau q}$ for affine $q$; this is how the
standard-form LP block $\langle c,x\rangle+I_{x\geq0}$ compiles to
$\max(x-\tau c,0)$. The emitted update is
\begin{equation}
 x_{k+1}=\operatorname{prox}_{\tau\varphi}\!\bigl(x_k-\tau
 L^{\top}u_k\bigr),\qquad
 u_{k+1}=\operatorname{prox}_{\sigma h^*}\!\bigl(u_k+\sigma
 L(2x_{k+1}-x_k)\bigr).
 \label{eq:inst-pdhg}
\end{equation}
Certificate condition: $\tau>0$, $\sigma>0$, $\tau\cdot \sigma\cdot \ell^2 < 1$. Fixed points are the saddle points
of the $(\varphi,h,L)$ Lagrangian; $\pi_G$ is the primal coordinate.

\paragraph{Condat--V\~u.}
\begin{equation}
 x_{k+1}=\operatorname{prox}_{\tau g}\!\bigl(x_k-\tau\nabla
 f(x_k)-\tau L^{\top}u_k\bigr),\qquad
 u_{k+1}=\operatorname{prox}_{\sigma h^*}\!\bigl(u_k+\sigma
 L(2x_{k+1}-x_k)\bigr).
 \label{eq:inst-cv}
\end{equation}
Certificate condition: $\tau\cdot L_f/2 + \tau\cdot \sigma\cdot \ell^2 < 1$. Fixed points are the saddle
points of problem~\eqref{eq:composite}; $\pi_G$ is the primal
coordinate. With $L_f=0$ the update and its condition reduce exactly to
PDHG, which the LP instantiation exercises.

\paragraph{Davis--Yin.}
\begin{equation}
 x_g^k=\operatorname{prox}_{\alpha g}(z_k),\qquad
 x_h^k=\operatorname{prox}_{\alpha h}\!\bigl(2x_g^k-z_k-\alpha\nabla
 f(x_g^k)\bigr),\qquad
 z_{k+1}=z_k+x_h^k-x_g^k.
 \label{eq:inst-dys}
\end{equation}
Certificate condition: $0 < \alpha < 2/L_f$. Fixed points correspond to
solutions through $\pi_G(z)=\operatorname{prox}_{\alpha g}(z)$.

\paragraph{Wrapper genes.}
The wrappers compile to state-space transformations of the certified map
$T_G$, exactly as certified: relaxation emits
$z_{k+1}=(1-\rho)z_k+\rho\,T_Gz_k$ (Lemma~\ref{lem:relaxation}); the
Halpern gene emits
\begin{equation}
 z_{k+1}=\tfrac{1}{j_k+2}\,z_{\rm anc}+\tfrac{j_k+1}{j_k+2}\,T_Gz_k,
 \label{eq:inst-halpern}
\end{equation}
where $j_k$ counts iterations since the last anchor reset and the restart
gene resets $z_{\rm anc}$ to the current iterate and $j_k$ to zero every
$k_{\rm r}$ iterations. The compiled anchor combination is applied to
$T_G$ itself, which is nonexpansive by (C1), rather than to the
normalized map $S_G$ of Proposition~\ref{prop:halpern}; the proposition
applies with $T_G$ in the role of the nonexpansive map when no anchor
restarts are used. It does not automatically give the same projection or rate
for periodically restarted trajectories.

\paragraph{Machine check.}
The correspondence between this table and the emitted code is checked
numerically, not formally proved: tests/test\_compiler\_equivalence.py holds an independent
reference implementation of every update above (plain NumPy, double
precision, written directly from the formulas, sharing no code with the
compiler) and asserts agreement with the compiled iteration
iterate-by-iterate for 50 iterations on 3 random small problems per
scheme, drawn from a TV $16\times16$ instance, a tiny netflow LP, and
random dense composites, at relative tolerance $10^{-10}$ in the
maximum norm, for the full solver state and the decoded primal iterate.
The battery covers the plain map, the production genome-to-phenotype path,
and each wrapper gene (relaxation, Halpern, Halpern with fixed-$k$
restart). The largest deviation observed at generation time was
$5.0\!\times\!10^{-13}$ (DRS dual form, where the reference solves the resolvent
system exactly while the compiled path runs conjugate gradients to
$10^{-12}$); per scheme: FBS $7.1\!\times\!10^{-17}$, DRS $5.0\!\times\!10^{-13}$, PDHG
$7.5\!\times\!10^{-16}$, Condat--V\~u $6.9\!\times\!10^{-16}$, Davis--Yin
$1.8\!\times\!10^{-16}$. This appendix is generated by
scripts/gen\_instantiation\_table.py, which extracts the certificate
conditions from the typing rules, verifies every display equation against
literal fragments of the builder source, and refuses to regenerate if a
scheme is added to the compiler without a corresponding entry and
reference implementation.


\end{document}
